% Options for packages loaded elsewhere
\PassOptionsToPackage{unicode}{hyperref}
\PassOptionsToPackage{hyphens}{url}
%
\documentclass[
  11pt,
  letterpaper,
]{article}
\usepackage{amsmath,amssymb}
\usepackage{iftex}
\ifPDFTeX
  \usepackage[T1]{fontenc}
  \usepackage[utf8]{inputenc}
  \usepackage{textcomp} % provide euro and other symbols
\else % if luatex or xetex
  \usepackage{unicode-math} % this also loads fontspec
  \defaultfontfeatures{Scale=MatchLowercase}
  \defaultfontfeatures[\rmfamily]{Ligatures=TeX,Scale=1}
\fi
\usepackage{lmodern}
\ifPDFTeX\else
  % xetex/luatex font selection
  \setmainfont[]{Arial}
  \setmonofont[]{Menlo}
\fi
% Use upquote if available, for straight quotes in verbatim environments
\IfFileExists{upquote.sty}{\usepackage{upquote}}{}
\IfFileExists{microtype.sty}{% use microtype if available
  \usepackage[]{microtype}
  \UseMicrotypeSet[protrusion]{basicmath} % disable protrusion for tt fonts
}{}
\makeatletter
\@ifundefined{KOMAClassName}{% if non-KOMA class
  \IfFileExists{parskip.sty}{%
    \usepackage{parskip}
  }{% else
    \setlength{\parindent}{0pt}
    \setlength{\parskip}{6pt plus 2pt minus 1pt}}
}{% if KOMA class
  \KOMAoptions{parskip=half}}
\makeatother
\usepackage{xcolor}
\usepackage[margin=0.8in]{geometry}
\usepackage{color}
\usepackage{fancyvrb}
\newcommand{\VerbBar}{|}
\newcommand{\VERB}{\Verb[commandchars=\\\{\}]}
\DefineVerbatimEnvironment{Highlighting}{Verbatim}{commandchars=\\\{\}}
% Add ',fontsize=\small' for more characters per line
\newenvironment{Shaded}{}{}
\newcommand{\AlertTok}[1]{\textcolor[rgb]{1.00,0.00,0.00}{\textbf{#1}}}
\newcommand{\AnnotationTok}[1]{\textcolor[rgb]{0.38,0.63,0.69}{\textbf{\textit{#1}}}}
\newcommand{\AttributeTok}[1]{\textcolor[rgb]{0.49,0.56,0.16}{#1}}
\newcommand{\BaseNTok}[1]{\textcolor[rgb]{0.25,0.63,0.44}{#1}}
\newcommand{\BuiltInTok}[1]{\textcolor[rgb]{0.00,0.50,0.00}{#1}}
\newcommand{\CharTok}[1]{\textcolor[rgb]{0.25,0.44,0.63}{#1}}
\newcommand{\CommentTok}[1]{\textcolor[rgb]{0.38,0.63,0.69}{\textit{#1}}}
\newcommand{\CommentVarTok}[1]{\textcolor[rgb]{0.38,0.63,0.69}{\textbf{\textit{#1}}}}
\newcommand{\ConstantTok}[1]{\textcolor[rgb]{0.53,0.00,0.00}{#1}}
\newcommand{\ControlFlowTok}[1]{\textcolor[rgb]{0.00,0.44,0.13}{\textbf{#1}}}
\newcommand{\DataTypeTok}[1]{\textcolor[rgb]{0.56,0.13,0.00}{#1}}
\newcommand{\DecValTok}[1]{\textcolor[rgb]{0.25,0.63,0.44}{#1}}
\newcommand{\DocumentationTok}[1]{\textcolor[rgb]{0.73,0.13,0.13}{\textit{#1}}}
\newcommand{\ErrorTok}[1]{\textcolor[rgb]{1.00,0.00,0.00}{\textbf{#1}}}
\newcommand{\ExtensionTok}[1]{#1}
\newcommand{\FloatTok}[1]{\textcolor[rgb]{0.25,0.63,0.44}{#1}}
\newcommand{\FunctionTok}[1]{\textcolor[rgb]{0.02,0.16,0.49}{#1}}
\newcommand{\ImportTok}[1]{\textcolor[rgb]{0.00,0.50,0.00}{\textbf{#1}}}
\newcommand{\InformationTok}[1]{\textcolor[rgb]{0.38,0.63,0.69}{\textbf{\textit{#1}}}}
\newcommand{\KeywordTok}[1]{\textcolor[rgb]{0.00,0.44,0.13}{\textbf{#1}}}
\newcommand{\NormalTok}[1]{#1}
\newcommand{\OperatorTok}[1]{\textcolor[rgb]{0.40,0.40,0.40}{#1}}
\newcommand{\OtherTok}[1]{\textcolor[rgb]{0.00,0.44,0.13}{#1}}
\newcommand{\PreprocessorTok}[1]{\textcolor[rgb]{0.74,0.48,0.00}{#1}}
\newcommand{\RegionMarkerTok}[1]{#1}
\newcommand{\SpecialCharTok}[1]{\textcolor[rgb]{0.25,0.44,0.63}{#1}}
\newcommand{\SpecialStringTok}[1]{\textcolor[rgb]{0.73,0.40,0.53}{#1}}
\newcommand{\StringTok}[1]{\textcolor[rgb]{0.25,0.44,0.63}{#1}}
\newcommand{\VariableTok}[1]{\textcolor[rgb]{0.10,0.09,0.49}{#1}}
\newcommand{\VerbatimStringTok}[1]{\textcolor[rgb]{0.25,0.44,0.63}{#1}}
\newcommand{\WarningTok}[1]{\textcolor[rgb]{0.38,0.63,0.69}{\textbf{\textit{#1}}}}
\usepackage{longtable,booktabs,array}
\usepackage{calc} % for calculating minipage widths
% Correct order of tables after \paragraph or \subparagraph
\usepackage{etoolbox}
\makeatletter
\patchcmd\longtable{\par}{\if@noskipsec\mbox{}\fi\par}{}{}
\makeatother
% Allow footnotes in longtable head/foot
\IfFileExists{footnotehyper.sty}{\usepackage{footnotehyper}}{\usepackage{footnote}}
\makesavenoteenv{longtable}
\setlength{\emergencystretch}{3em} % prevent overfull lines
\providecommand{\tightlist}{%
  \setlength{\itemsep}{0pt}\setlength{\parskip}{0pt}}
\setcounter{secnumdepth}{-\maxdimen} % remove section numbering
\ifLuaTeX
\usepackage[bidi=basic]{babel}
\else
\usepackage[bidi=default]{babel}
\fi
\babelprovide[main,import]{american}
\ifPDFTeX
\else
\babelfont{rm}[]{Arial}
\fi
% get rid of language-specific shorthands (see #6817):
\let\LanguageShortHands\languageshorthands
\def\languageshorthands#1{}
\ifLuaTeX
  \usepackage{selnolig}  % disable illegal ligatures
\fi
\IfFileExists{bookmark.sty}{\usepackage{bookmark}}{\usepackage{hyperref}}
\IfFileExists{xurl.sty}{\usepackage{xurl}}{} % add URL line breaks if available
\urlstyle{same}
\hypersetup{
  pdftitle={Mass filter equations and software guide},
  pdflang={en-US},
  hidelinks,
  pdfcreator={LaTeX via pandoc}}

\title{Mass filter equations and software guide}
\usepackage{etoolbox}
\makeatletter
\providecommand{\subtitle}[1]{% add subtitle to \maketitle
  \apptocmd{\@title}{\par {\large #1 \par}}{}{}
}
\makeatother
\subtitle{Sources and reproducible conventions for the website and
TrapLab}
\author{}
\date{2026-10-24}

\begin{document}
\maketitle

Public edition 2026-10-24. Open a terminal in the separately installed
companion application folder before running its commands. Module labels
identify its entry points and are not download links. Equations and
official references are preserved from the checked source edition.
Canonical Markdown SHA-256:
\texttt{f858a87ed961a5e62a4cfe32468744f8fa0f219578702bf1290b65d009cb88a6}.

This guide records all 38 shared mass-related quantities, the separate
native legacy W indicator, and 28 legacy binding-energy recipes in the
appendix, with exact stencils, units, signs, practical uses, sources and
program entry points. Choose the physical question first, supply
consistent masses, then interpret the output and its uncertainty.
Equations follow the current source inventory; literature definitions
and implementation conventions are distinguished. This is not a new mass
evaluation, model-accuracy benchmark or measurement.

Machine-readable inventory and sources: current inventory ·
sources.json.

\section{Start with one reproducible
question}\label{start-with-one-reproducible-question}

To study neutron binding along the tin isotopes, choose s2n and fix Z at
50. The reference is AME2020. If the 116Sn mass excess is deliberately
increased by 10 keV, the same primitive-mass stencil changes S₂ₙ of
116Sn by −10 keV, S₂ₙ of 118Sn by +10 keV, and the 116Sn two-neutron
shell gap by −20 keV. This explicitly synthetic exercise demonstrates
neighbouring sensitivity; it is not a new measurement. A separately
rounded published reaction table can retain a small printed-value
difference.

\subsection{Website workflow}\label{website-workflow}

Open the website Lab page and its chart of nuclides. Select a key from
the chain-plot quantity menu, choose a fixed-Z, fixed-N or fixed-A
chain, and select the reference nuclide and range. Colour layers,
classed legends, manual colour limits, magic lines and N = Z are display
controls; they do not change masses. When choosing AME, new inputs,
mixed stencils or models, inspect the series source and uncertainty
basis together.

\href{https://gezhuang0717.github.io/traps.github.io/lab/}{Website Lab}

Current website source:
\texttt{layouts/\_shortcodes/nuclide-chart.html};
\texttt{static/js/nuclide-chart.js}.

\subsection{Native applications and individual
code}\label{native-applications-and-individual-code}

Use the existing launcher from the GUI checkout below. Both standalone
entry points share the mass-workbench engine; the independent plotter
generates data and static figures. These are operating instructions:
this documentation build did not start Qt or rerun an analysis. Choose a
fresh output directory and preserve the original inputs.

\begin{Shaded}
\begin{Highlighting}[]

\ExtensionTok{./Run\_private.command}\NormalTok{ apps/ame\_nubase.py }\AttributeTok{{-}{-}quantity}\NormalTok{ s2n }\AttributeTok{{-}{-}z}\NormalTok{ 50 }\AttributeTok{{-}{-}minimum}\NormalTok{ 60 }\AttributeTok{{-}{-}maximum}\NormalTok{ 74}
\ExtensionTok{./Run\_private.command}\NormalTok{ apps/mulberry\_mass.py }\AttributeTok{{-}{-}quantity}\NormalTok{ s2n }\AttributeTok{{-}{-}z}\NormalTok{ 50 }\AttributeTok{{-}{-}minimum}\NormalTok{ 60 }\AttributeTok{{-}{-}maximum}\NormalTok{ 74}
\ExtensionTok{./Run\_private.command}\NormalTok{ scripts/plot\_mass\_workbench.py }\AttributeTok{{-}{-}help}
\end{Highlighting}
\end{Shaded}

Each entry and the code index identify its functions, files and launch
arguments. GUI session JSON preserves inputs and the view; CSV preserves
editable masses and results; PNG, PDF and SVG are figures. Saving does
not validate measurement provenance, and static mass comparisons need no
physical-motion video. Windows execution and native file-dialog
acceptance are outside this documentation build.

\section{Typed data input workflow}\label{typed-data-input-workflow}

The current website My data panel treats a pasted table as a typed
observable, never as an automatic mass replacement. For headerless
\texttt{50,50,8}, the first two columns identify Z = 50 and N = 50 and
the third is the value 8. Choose the quantity and input unit explicitly.
The initial type is Custom with an explicit custom-unit label; it is not
ME. Per-row quantity, unit, label or reference columns are not applied
automatically. Split tables containing different physical quantities or
units.

\begin{enumerate}
\def\labelenumi{\arabic{enumi}.}
\tightlist
\item
  Open My data, paste \texttt{50,50,8}, select BE/A and MeV/nucleon,
  choose the value column, leave sigma unknown, and enter a series label
  and reference. Press Apply. The own point is 8 MeV/nucleon and the
  conversion preview is total BE = 800 MeV for A = 100. Selecting total
  BE instead means 8 MeV total. Selecting ME and keV makes the chain
  point 0.008 MeV. These are deliberately chosen input examples, not
  measured records.
\item
  To supply uncertainties, add a fourth numeric column and select it as
  sigma. Sigma must be nonnegative and use the same input unit as the
  value. A blank or unselected sigma is unknown; an entered zero is a
  reported zero. Both value and sigma undergo the same unit conversion.
\item
  Own points, AME and selected model checkboxes are independent. All 38
  mass-filter quantities can be supplied as direct values and compared
  with the same available reference quantity. Radii and moments compare
  with the existing IAEA data; E2 and upward B(E2) use NNDC data. A
  custom observable has its own series only. Direct-input residuals are
  not calculated, and mass/max-sigma filters apply to reference/model
  series while direct points stay visible.
\item
  Use Append as measured masses or Replace measured masses only when you
  explicitly want a mass-table change. ME, total BE, BE/A and neutral
  atomic mass are the only adoptable types; ground-state catalogue
  identities are required. Append keeps prior rows and the current
  duplicate policy; Replace is explicit. All rows are validated before
  mutation. Direct Sn or other filter/property values remain observables
  and cannot become masses. Weighted duplicates containing unknown/zero
  sigma have no active weighted aggregate; choose last-row policy or
  provide positive sigmas.
\item
  Save inputs and settings JSON, then reopen through File or paste it
  and Apply. Raw text, selected value/sigma columns, quantity/unit,
  label/reference and own/reference/model choices are retained. Chain
  CSV records direct-input provenance; mass-table CSV/JSON additionally
  retain original quantity/value/unit, uncertainty status and input
  origin. Existing plot PNG/CSV controls operate on the plotted
  quantity; their native downloads are a separate browser acceptance
  task.
\end{enumerate}

Energy chain values normalize to MeV, BE/A to MeV/nucleon, atomic mass
to u, radii to fm, magnetic moments to μ\_N, quadrupole moments to b, E2
to keV and upward B(E2) to e²b². The adapter scales e²fm⁴ by 1e-4 to
e²b² and quadrupole fm² by 0.01 to b. It introduces no new fit,
covariance estimate or mass-model equation.

Implementation: \texttt{static/js/chart-typed-data.js} provides units,
adapt, massRows, portable and restore;
\texttt{static/js/nuclide-chart.js} binds the controls; measured-masses
and measured-masses-ui handle explicit adoption. The source-freeze
receipt records 88 targeted Node and 14 targeted Python tests. This is
source/fixture evidence, not a claim that file selection, downloads or a
full browser workflow were verified by the manual build.

Receipt: typed input implementation and source freeze.

\section{Notation units and
uncertainty}\label{notation-units-and-uncertainty}

Z is proton number, N is neutron number and \(A=Z+N\). M in the
equations denotes neutral-atom mass excess ME, normally in keV; some
plots divide both value and uncertainty by 1000 to show MeV. B or BE is
positive total binding energy and BEA is binding energy per nucleon.
Symbols n, H, d and α denote the mass excesses of a neutron, neutral
hydrogen, deuterium and helium. State assignments must agree, and a
missing neighbour must not be replaced by zero.

\[
\mathrm{ME}(Z,N)=\bigl[M_{\mathrm{atom}}(Z,N)-Au\bigr]c^2,\qquad A=Z+N
\]

\[
F(Z,N)=\sum_i c_i\,\mathrm{ME}(Z+\Delta Z_i,N+\Delta N_i)+C
\]

Stencil coefficients are dimensionless. Offset C must use the same
energy unit as the output. Combine repeated occurrences of an identical
primitive mass before propagating uncertainty. Full covariance and its
independent approximation are:

\[
\sigma_F^2=\boldsymbol{c}^{\mathsf T}\mathbf C\boldsymbol{c},\qquad \sigma_{F,\mathrm{ind}}^2=\sum_i c_i^2\sigma_i^2
\]

Directly published AME energies and errors retain evaluation
correlations. Marginal mass errors used for new inputs or arbitrary
stencils are not full evaluated covariance. All fourteen model mass
uncertainties remain unknown; extra model columns or model spread must
not be relabelled experimental uncertainty. The following explanation
reuses the existing bilingual AME audit guide.

Let a quantity be a linear combination of mass excesses. Its variance
depends on both individual mass errors and correlations:

\[
q=\sum_i c_i\,\mathrm{ME}_i,\qquad
\sigma_q^2=\sum_i c_i^2\sigma_i^2+2\sum_{i<j}c_i c_j\,\mathrm{Cov}(\mathrm{ME}_i,\mathrm{ME}_j).
\]

For two masses in a decay-energy difference:

\[
q=\mathrm{ME}_P-\mathrm{ME}_D,\qquad
\sigma_q^2=\sigma_P^2+\sigma_D^2-2\,\mathrm{Cov}(\mathrm{ME}_P,\mathrm{ME}_D).
\]

The AME least-squares evaluation links masses through shared
mass-spectrometric, decay and reaction constraints. Its published energy
uncertainties retain these evaluation correlations. Setting every
off-diagonal covariance to zero is a separate approximation, documented
as independent propagation. A positive parent--daughter covariance
reduces the uncertainty of their difference. A coefficient with the same
sign changes that effect. The figure compares the published result with
the diagonal approximation; it does not reconstruct the covariance
matrix. See \href{https://doi.org/10.1088/1674-1137/abddb0}{AME2020 Part
I}.

Atomic masses and mass excesses give the following ground-state
expressions. ME denotes mass excess expressed as an energy, and the
neutron ME uses the same AME source:

\[
Q_{\beta^-}=\mathrm{ME}(Z,N)-\mathrm{ME}(Z+1,N-1),
\]

\[
Q_{2\beta^-}=\mathrm{ME}(Z,N)-\mathrm{ME}(Z+2,N-2),\qquad
S_{2n}=\mathrm{ME}(Z,N-2)+2\mathrm{ME}_n-\mathrm{ME}(Z,N).
\]

Other double-beta channels need the correct electron thresholds: the
neutral-atom double-electron-capture mass difference is the parent minus
the daughter at Z−2; EC plus beta-plus subtracts two electron rest
energies, and double beta-plus subtracts four. Capture-hole energies and
daughter excitation energies are additional channel-specific
requirements. A ground-state mass difference alone does not establish a
resonant capture condition, decay rate or branching ratio.

\subsection{Constants and implementation
differences}\label{constants-and-implementation-differences}

\begin{longtable}[]{@{}llr@{}}
\toprule\noalign{}
Profile & Constant key & Value \\
\midrule\noalign{}
\endhead
\bottomrule\noalign{}
\endlastfoot
web & u\_keV & 931494.10242 \\
web & electron\_u & 0.000548579909065 \\
web & u\_kg & 1.6605390666e-27 \\
web & electron\_keV & 510.99895000014743 \\
native & u\_keV & 931494.10242 \\
native & electron\_u & 0.0005485799090441 \\
native & u\_kg & 1.66053906892e-27 \\
native & electron\_keV & 510.99894998067913 \\
legacy\_gsi & ME\_n\_keV & 8071.3181 \\
legacy\_gsi & ME\_H\_keV & 7288.97106 \\
legacy\_gsi & electron\_u & 0.000548579909065 \\
\end{longtable}

u\_keV is one atomic-mass unit in rest energy (keV); electron\_u is
electron mass in u and electron\_keV its rest energy; u\_kg is in kg.
Legacy ME\_n\_keV and ME\_H\_keV are in keV.

\begin{longtable}[]{@{}lrr@{}}
\toprule\noalign{}
Reference particle & ME keV & σ keV \\
\midrule\noalign{}
\endhead
\bottomrule\noalign{}
\endlastfoot
n & 8071.31806 & 0.00044 \\
H & 7288.971064 & 1.3e-05 \\
d & 13135.722895 & 1.5e-05 \\
alpha & 2424.91587 & 0.00015 \\
\end{longtable}

These constants describe the current code snapshot. The website and
native default use different CODATA editions for electron rest energy;
preserve the documented difference rather than forcing numerical
identity. Positron channels need their atomic-mass electron thresholds;
an EC mass difference is not a shell-resolved neutrino energy.

\textbf{legacy-W} · Native legacy W uses fixed quarter-normalized D2pS2n
at adjacent isobar centers. Shared wig uses parity-dependent vpn at N±2
with fixed Z. R W1/W2/W3 are three additional distinct stencils.

\textbf{pn-parity} · Legacy native dVpn is an alias for fixed D2pS2n,
not the general parity-dependent vpn. They agree at even-even centers
only. Other fixed mixed filters match the physical automatic parity
branch as stated.

\textbf{electron-constants} · Web uses AME2020/CODATA2018
electron\_u=5.48579909065e-4. Native default uses CODATA2022
electron\_u=5.485799090441e-4 with AME2020 u\_keV=931494.10242. qecbp
and q2bp differ by the corresponding 2e/4e threshold. Constants are
fixed in these errors; no independent electron/u uncertainty is added.

\textbf{SI-u} · Web u\_kg=1.66053906660e-27 vs native
default1.66053906892e-27. This affects mass/frequency conversions, not
pure keV mass stencils. Native AME2020\_COMPAT is separately selectable
in APIs.

\textbf{input-sigma-hash} · Native experimental workbench permits
separate trailing \# on value and uncertainty. Web measured-masses.js
accepts \# on value but requires plain numeric positive sigma. Public
AME primitives preserve separate value and sigma flags in both.

\textbf{light-models} · Web mass-surface uses experimental n/H/alpha
(and experimental deuterium fallback) even for models. Native ModelData
fills missing light species from AME with unknown sigma, but does not
override a light species already present in a model payload.
Constants.py rounded legacy light values are separate from the richer
canonical reference data.

\textbf{published-vs-raw} · Both default canonical reference surfaces
prefer published AME values/evaluated uncertainties or exact reverse
mass differences. Supplying custom/model maps disables these overrides.
Small value differences from independently rounded source columns are
retained.

\textbf{legacy-Gsym} · R metadata prose states 2δ−δ−δ without /2, but
executable CSV weights and canonical Gsym have /2. Inventory documents
executable behavior.

\textbf{recipe-units} · R-derived recipes use BE inputs and mixed output
units MeV/keV, fixed R H=7288.97106 keV, and explicit Qec offset.
Canonical surface uses ME primitives and all outputs keV except BEA. A
same-looking key is not a unit guarantee.

\textbf{recipe-covariance} · Canonical Surface has diagonal primitive
covariance only, plus direct evaluated override. GSI
apply\_energy\_operator optionally accepts a user-supplied input-row
covariance; science.nuclear\_mass\_operator instead expects covariance
in original coefficient order, including zero weights. These matrix
layouts must not be interchanged.

\textbf{old-manual-constant-description} · Existing
NUCLEAR\_REFERENCE\_EN.md says website CODATA2018 thresholds for the
adapter; current native default uses the mixed compatibility profile
stated above. Treat older claims of complete numeric web/native parity
as historical evidence at their source hashes.

\textbf{edition-uncertainty} · Chart edition difference is
ME2020−MEedition and hypot of marginal sigmas when both known.
Evaluations share data/correlations, so this is a diagonal
approximation, not a published evaluated cross-edition uncertainty.

\textbf{model-residual-error} · Chart dmod shows MEAME−MEmodel and
experimental sigma only; this is conditional data uncertainty, not total
theory predictive error. Native model residual uncertainty remains
unknown when model errors are absent.

\subsection{Supporting conversion and processing
equations}\label{supporting-conversion-and-processing-equations}

\subsubsection{linear\_uncertainty}\label{linear_uncertainty}

\[
F=\sum_i c_i M_i+C,\quad \sigma_F^2=\sum_i c_i^2\sigma_i^2\quad(\mathrm{diagonal})
\]

Combine identical primitive identities before squaring. A
zero-coefficient absent/unknown mass has no effect. Unknown nonzero
sigma makes sigma unknown. Genuine zero sigma stays zero. Constant
offsets are fixed.

\subsubsection{full\_covariance}\label{full_covariance}

\[
\sigma_F^2=\boldsymbol c^\mathsf{T}\mathbf C\boldsymbol c=\sum_i c_i^2C_{ii}+2\sum_{i<j}c_ic_jC_{ij}
\]

Implemented only when explicit covariance is supplied to GSI adapters;
direct evaluated AME errors are retained without reconstructing C.

\subsubsection{atomic\_mass}\label{atomic_mass}

\[
M_{\mathrm{atom}}/u=A+M(Z,N)/(u c^2),\quad u c^2=931494.10242\ {\mathrm{keV}}
\]

Mass excess M in energy keV, not absolute atomic mass. Atomic mass
sigma=sigmaME/uKeV; fixed conversion constant.

\subsubsection{ion\_mass}\label{ion_mass}

\[
M_{\mathrm{ion}}/u=M_{\mathrm{atom}}/u-q m_e/u+E_{\mathrm{ion}}/(u c^2)
\]

Signed q. Binding/ionization correction positive increases ion mass
relative to pure electron subtraction. Binding input eV requires
denominator1000*uKeV. No binding or constant uncertainty invented. Web
catalogue resolve omits these binding energies.

\subsubsection{ratio\_mass}\label{ratio_mass}

\[
r=\nu_{\mathrm{ref}}/\nu_{\mathrm{IOI}},\quad M_{\mathrm{IOI,ion}}=r(q_{\mathrm{IOI}}/q_{\mathrm{ref}})M_{\mathrm{ref,ion}}
\]

Neutralization adds qIOI electron masses and subtracts IOI ionization
correction. Charges are positive magnitudes in ratio\_mass\_result.
Independent sigma contributions factor\emph{MrefIon}sigma\_r and
factor\emph{r}sigma\_ref are added in quadrature.

\subsubsection{weighted\_duplicates}\label{weighted_duplicates}

\[
\bar M=\frac{\sum_i M_i/\sigma_i^2}{\sum_i1/\sigma_i^2},\quad \sigma_{\bar M}=(\sum_i1/\sigma_i^2)^{-1/2}\max[1,\sqrt{\chi^2/(k-1)}]
\]

Independent duplicate input rows only; last-row policy is an
alternative. No facility/systematic covariance or global AME adjustment.

\subsubsection{residual\_cancellation}\label{residual_cancellation}

\[
\Delta F=F_{\mathrm{New}}-F_{\mathrm{AME}},\quad \sigma_{\Delta F}^2=\sum_{\text{replaced }i}c_i^2(\sigma_{{\mathrm{New}},i}^2+\sigma_{{\mathrm{AME}},i}^2)
\]

Native Workbench residual rule; unchanged shared AME primitives cancel.
It is a declared diagonal comparison, not propagation of independent
nested reaction errors.

\subsubsection{single\_beta\_plus}\label{single_beta_plus}

\[
Q_{\beta^+}=M(Z,N)-M(Z-1,N+1)-2e
\]

Implemented native physics.q\_beta\_plus and decay.evaluate, separate
from the canonical qecbp double-weak mode.

\subsubsection{state\_Qstar}\label{state_qstar}

\[
Q^*=Q_{\mathrm{selected\ states}}-E_x,\quad \sigma_{Q^*}=\sqrt{\sigma_Q^2+\sigma_{E_x}^2}
\]

Native decay.evaluate. Selected daughter-state excitation already in its
mass; Ex is ADDITIONAL energy. Ground-ground evaluated lookup when
available; otherwise selected NUBASE masses diagonal. Missing sigma
remains unknown.

\subsubsection{general\_reaction}\label{general_reaction}

\[
Q=\sum_{\mathrm{reactants}}a_i M_i-\sum_{\mathrm{products}}b_j M_j
\]

NuclearData.reaction checks total A and Z conservation for neutral
atoms, combines repeated identities and propagates diagonal primitive
sigma. No evaluated covariance or molecular binding correction.

\subsubsection{drip\_boundary}\label{drip_boundary}

\[
S_n>0,\ S_n(N+1)\leq0\quad\text{or analogous }S_{2n},S_p,S_{2p}
\]

dripBoundary returns last positive integer-index sample only if its
immediate index+1 neighbor is finite and nonpositive. Missing samples
and source-table edges are not physical drip lines; no interpolation.

\subsubsection{legacy\_four\_mass\_tool}\label{legacy_four_mass_tool}

\[
S_n=M_{N-1}+n-M_N,\quad S_{2n}=M_{N-2}+2n-M_N,\quad Q_{\mathrm{EC}}=M_N-M_{\mathrm{daughter}}
\]

gsi\_mrtof.science.mass\_quantities: optional four independent
nonnegative input sigmas; fixed historical neutron ME offset. No
covariance lookup.

\section{Equations and practical use by
quantity}\label{equations-and-practical-use-by-quantity}

Every entry gives the actual equation and primitive-mass stencil. Keys
remain unchanged so they can be searched in menus, command lines and
source. Use the equation's sign and parity factors rather than inferring
them from a name. A positive Q only establishes an energy allowance for
its stated channel; a difference or cusp does not identify a unique
structural mechanism.

\subsection{1 me Mass excess}\label{me-mass-excess}

Neutral atomic mass excess

Use and what it reveals: Compare the atomic mass with the integer mass
number; inspect evaluation residuals.

\[
M(Z,N)
\]

Output unit: keV. Inputs are all masses in the state-consistent stencil
below; missing terms make the result unavailable.

\begin{longtable}[]{@{}rrr@{}}
\toprule\noalign{}
ΔZ & ΔN & Coefficient \\
\midrule\noalign{}
\endhead
\bottomrule\noalign{}
\endlastfoot
0 & 0 & 1 \\
\end{longtable}

Constant offset:

\[
C=0
\]

Convention: Neutral atomic mass excess, ground states, fixed Z/N
coordinate convention.

Source index: \texttt{ame2020}, \texttt{filters-implementation}.

\subsubsection{Checkable example}\label{checkable-example}

\textbf{116Sn} · Calculated from retained source data

-91525.979 ± 0.096 keV

Default evaluated reference: -91525.979 ± 0.096 keV

\href{https://www-nds.iaea.org/amdc/ame2020/mass_1.mas20.txt}{Original
source table}

Recalculated and evaluated values/errors are distinct; source rounding
is retained. Full primitive inputs, flags and hashes are in this entry's
worked\_example in inventory-delivery.json.

Uncertainty: use the primitive-identity combination rule above. Keep
directly evaluated reference errors separate from independently
recalculated errors; model uncertainty is unavailable.

\subsubsection{Program mapping}\label{program-mapping}

\begin{itemize}
\tightlist
\item
  energies.py ·
  \texttt{Surface.get\ /\ Surface.\_linear\ /\ Reference.published} ·
  line 35
\item
  mass-surface.js · \texttt{create.derive} · line 26
\end{itemize}

\begin{Shaded}
\begin{Highlighting}[]
\ExtensionTok{./Run\_private.command} \AttributeTok{{-}m}\NormalTok{ traplab\_nuclear quantity 116Sn me }\AttributeTok{{-}{-}format}\NormalTok{ csv}
\end{Highlighting}
\end{Shaded}

\subsection{2 BE Total binding energy}\label{be-total-binding-energy}

Positive total nuclear binding energy in the atomic-hydrogen convention

Use and what it reveals: Compare total binding; neutral hydrogen
accounts for atomic electrons.

\[
B(Z,N)=Z H+N n-M(Z,N)
\]

Output unit: keV. Inputs are all masses in the state-consistent stencil
below; missing terms make the result unavailable.

\begin{longtable}[]{@{}rrr@{}}
\toprule\noalign{}
ΔZ & ΔN & Coefficient \\
\midrule\noalign{}
\endhead
\bottomrule\noalign{}
\endlastfoot
0 & 0 & -1 \\
\end{longtable}

Light-particle terms

\begin{longtable}[]{@{}lr@{}}
\toprule\noalign{}
Particle & Coefficient \\
\midrule\noalign{}
\endhead
\bottomrule\noalign{}
\endlastfoot
H & Z \\
n & N \\
\end{longtable}

Constant offset:

\[
C=0
\]

Convention: Neutral atomic mass excess, ground states, fixed Z/N
coordinate convention.

Source index: \texttt{ame2020}, \texttt{filters-implementation}.

\subsubsection{Checkable example}\label{checkable-example-1}

\textbf{116Sn} · Calculated from retained source data

988681.52416 ± 0.10029827565815876 keV

Default evaluated reference: 988681.5256 ± 0.1044 keV

\href{https://www-nds.iaea.org/amdc/ame2020/mass_1.mas20.txt}{Original
source table}

Recalculated and evaluated values/errors are distinct; source rounding
is retained. Full primitive inputs, flags and hashes are in this entry's
worked\_example in inventory-delivery.json.

Uncertainty: use the primitive-identity combination rule above. Keep
directly evaluated reference errors separate from independently
recalculated errors; model uncertainty is unavailable.

\subsubsection{Program mapping}\label{program-mapping-1}

\begin{itemize}
\tightlist
\item
  energies.py ·
  \texttt{Surface.get\ /\ Surface.\_linear\ /\ Reference.published} ·
  line 35
\item
  mass-surface.js · \texttt{create.derive} · line 26
\end{itemize}

\begin{Shaded}
\begin{Highlighting}[]
\ExtensionTok{./Run\_private.command} \AttributeTok{{-}m}\NormalTok{ traplab\_nuclear quantity 116Sn BE }\AttributeTok{{-}{-}format}\NormalTok{ csv}
\end{Highlighting}
\end{Shaded}

\subsection{3 BEA Binding energy /
nucleon}\label{bea-binding-energy-nucleon}

Binding energy per nucleon

Use and what it reveals: Compare binding on a per-nucleon scale; divide
by A, not by Z or N.

\[
B(Z,N)/A,\quad A=Z+N
\]

Output unit: keV/nucleon. Inputs are all masses in the state-consistent
stencil below; missing terms make the result unavailable.

\begin{longtable}[]{@{}rrr@{}}
\toprule\noalign{}
ΔZ & ΔN & Coefficient \\
\midrule\noalign{}
\endhead
\bottomrule\noalign{}
\endlastfoot
0 & 0 & -1/A \\
\end{longtable}

Light-particle terms

\begin{longtable}[]{@{}lr@{}}
\toprule\noalign{}
Particle & Coefficient \\
\midrule\noalign{}
\endhead
\bottomrule\noalign{}
\endlastfoot
H & Z/A \\
n & N/A \\
\end{longtable}

Constant offset:

\[
C=0
\]

Convention: Neutral atomic mass excess, ground states, fixed Z/N
coordinate convention.

Source index: \texttt{ame2020}, \texttt{filters-implementation}.

\subsubsection{Checkable example}\label{checkable-example-2}

\textbf{116Sn} · Calculated from retained source data

8523.116587586206 ± 0.0008646403073979202 keV/nucleon

Default evaluated reference: 8523.1166 ± 0.0009 keV/nucleon

\href{https://www-nds.iaea.org/amdc/ame2020/mass_1.mas20.txt}{Original
source table}

Recalculated and evaluated values/errors are distinct; source rounding
is retained. Full primitive inputs, flags and hashes are in this entry's
worked\_example in inventory-delivery.json.

Uncertainty: use the primitive-identity combination rule above. Keep
directly evaluated reference errors separate from independently
recalculated errors; model uncertainty is unavailable.

\subsubsection{Program mapping}\label{program-mapping-2}

\begin{itemize}
\tightlist
\item
  energies.py ·
  \texttt{Surface.get\ /\ Surface.\_linear\ /\ Reference.published} ·
  line 35
\item
  mass-surface.js · \texttt{create.derive} · line 26
\end{itemize}

\begin{Shaded}
\begin{Highlighting}[]
\ExtensionTok{./Run\_private.command} \AttributeTok{{-}m}\NormalTok{ traplab\_nuclear quantity 116Sn BEA }\AttributeTok{{-}{-}format}\NormalTok{ csv}
\end{Highlighting}
\end{Shaded}

\subsection{4 sn One-neutron separation
Sn}\label{sn-one-neutron-separation-sn}

Energy required to remove one neutron

Use and what it reveals: Inspect one-neutron thresholds and local
changes in binding.

\[
S_n=M(Z,N-1)+n-M(Z,N)
\]

Output unit: keV. Inputs are all masses in the state-consistent stencil
below; missing terms make the result unavailable.

Menu or code aliases: \texttt{S1n}.

\begin{longtable}[]{@{}rrr@{}}
\toprule\noalign{}
ΔZ & ΔN & Coefficient \\
\midrule\noalign{}
\endhead
\bottomrule\noalign{}
\endlastfoot
0 & -1 & 1 \\
0 & 0 & -1 \\
\end{longtable}

Light-particle terms

\begin{longtable}[]{@{}lr@{}}
\toprule\noalign{}
Particle & Coefficient \\
\midrule\noalign{}
\endhead
\bottomrule\noalign{}
\endlastfoot
n & 1 \\
\end{longtable}

Constant offset:

\[
C=0
\]

Convention: Neutral atomic mass excess, ground states, fixed Z/N
coordinate convention.

Source index: \texttt{ame2020}, \texttt{filters-implementation}.

\subsubsection{Checkable example}\label{checkable-example-3}

\textbf{116Sn} · Calculated from retained source data

9563.451060000001 ± 0.09716580468456996 keV

Default evaluated reference: 9563.4512 ± 0.0948 keV

\href{https://www-nds.iaea.org/amdc/ame2020/rct2_1.mas20.txt}{Original
source table}

Recalculated and evaluated values/errors are distinct; source rounding
is retained. Full primitive inputs, flags and hashes are in this entry's
worked\_example in inventory-delivery.json.

Uncertainty: use the primitive-identity combination rule above. Keep
directly evaluated reference errors separate from independently
recalculated errors; model uncertainty is unavailable.

\subsubsection{Program mapping}\label{program-mapping-3}

\begin{itemize}
\tightlist
\item
  energies.py ·
  \texttt{Surface.get\ /\ Surface.\_linear\ /\ Reference.published} ·
  line 35
\item
  mass-surface.js · \texttt{create.derive} · line 27
\end{itemize}

\begin{Shaded}
\begin{Highlighting}[]
\ExtensionTok{./Run\_private.command} \AttributeTok{{-}m}\NormalTok{ traplab\_nuclear quantity 116Sn sn }\AttributeTok{{-}{-}format}\NormalTok{ csv}
\end{Highlighting}
\end{Shaded}

\subsection{5 s2n Two-neutron separation
S2n}\label{s2n-two-neutron-separation-s2n}

Energy required to remove two neutrons

Use and what it reveals: Track isotopic trends with reduced odd-even
staggering.

\[
S_{2n}=M(Z,N-2)+2n-M(Z,N)
\]

Output unit: keV. Inputs are all masses in the state-consistent stencil
below; missing terms make the result unavailable.

Menu or code aliases: \texttt{S2n}.

\begin{longtable}[]{@{}rrr@{}}
\toprule\noalign{}
ΔZ & ΔN & Coefficient \\
\midrule\noalign{}
\endhead
\bottomrule\noalign{}
\endlastfoot
0 & -2 & 1 \\
0 & 0 & -1 \\
\end{longtable}

Light-particle terms

\begin{longtable}[]{@{}lr@{}}
\toprule\noalign{}
Particle & Coefficient \\
\midrule\noalign{}
\endhead
\bottomrule\noalign{}
\endlastfoot
n & 2 \\
\end{longtable}

Constant offset:

\[
C=0
\]

Convention: Neutral atomic mass excess, ground states, fixed Z/N
coordinate convention.

Source index: \texttt{ame2020}, \texttt{filters-implementation}.

\subsubsection{Checkable example}\label{checkable-example-4}

\textbf{116Sn} · Calculated from retained source data

17108.880120000005 ± 0.10028845596577904 keV

Default evaluated reference: 17108.8797 ± 0.098 keV

\href{https://www-nds.iaea.org/amdc/ame2020/rct1.mas20.txt}{Original
source table}

Recalculated and evaluated values/errors are distinct; source rounding
is retained. Full primitive inputs, flags and hashes are in this entry's
worked\_example in inventory-delivery.json.

Uncertainty: use the primitive-identity combination rule above. Keep
directly evaluated reference errors separate from independently
recalculated errors; model uncertainty is unavailable.

\subsubsection{Program mapping}\label{program-mapping-4}

\begin{itemize}
\tightlist
\item
  energies.py ·
  \texttt{Surface.get\ /\ Surface.\_linear\ /\ Reference.published} ·
  line 35
\item
  mass-surface.js · \texttt{create.derive} · line 27
\end{itemize}

\begin{Shaded}
\begin{Highlighting}[]
\ExtensionTok{./Run\_private.command} \AttributeTok{{-}m}\NormalTok{ traplab\_nuclear quantity 116Sn s2n }\AttributeTok{{-}{-}format}\NormalTok{ csv}
\end{Highlighting}
\end{Shaded}

\subsection{6 sp One-proton separation
Sp}\label{sp-one-proton-separation-sp}

Atomic-convention one-proton separation energy

Use and what it reveals: Inspect proton-emission thresholds. H means
neutral 1H, not bare-proton mass excess.

\[
S_p=M(Z-1,N)+H-M(Z,N)
\]

Output unit: keV. Inputs are all masses in the state-consistent stencil
below; missing terms make the result unavailable.

Menu or code aliases: \texttt{S1p}.

\begin{longtable}[]{@{}rrr@{}}
\toprule\noalign{}
ΔZ & ΔN & Coefficient \\
\midrule\noalign{}
\endhead
\bottomrule\noalign{}
\endlastfoot
-1 & 0 & 1 \\
0 & 0 & -1 \\
\end{longtable}

Light-particle terms

\begin{longtable}[]{@{}lr@{}}
\toprule\noalign{}
Particle & Coefficient \\
\midrule\noalign{}
\endhead
\bottomrule\noalign{}
\endlastfoot
H & 1 \\
\end{longtable}

Constant offset:

\[
C=0
\]

Convention: Neutral atomic mass excess, ground states, fixed Z/N
coordinate convention.

Source index: \texttt{ame2020}, \texttt{filters-implementation}.

\subsubsection{Checkable example}\label{checkable-example-5}

\textbf{116Sn} · Calculated from retained source data

9278.593064000004 ± 0.09674709385299385 keV

Default evaluated reference: 9278.5934 ± 0.0953 keV

\href{https://www-nds.iaea.org/amdc/ame2020/rct2_1.mas20.txt}{Original
source table}

Recalculated and evaluated values/errors are distinct; source rounding
is retained. Full primitive inputs, flags and hashes are in this entry's
worked\_example in inventory-delivery.json.

Uncertainty: use the primitive-identity combination rule above. Keep
directly evaluated reference errors separate from independently
recalculated errors; model uncertainty is unavailable.

\subsubsection{Program mapping}\label{program-mapping-5}

\begin{itemize}
\tightlist
\item
  energies.py ·
  \texttt{Surface.get\ /\ Surface.\_linear\ /\ Reference.published} ·
  line 35
\item
  mass-surface.js · \texttt{create.derive} · line 28
\end{itemize}

\begin{Shaded}
\begin{Highlighting}[]
\ExtensionTok{./Run\_private.command} \AttributeTok{{-}m}\NormalTok{ traplab\_nuclear quantity 116Sn sp }\AttributeTok{{-}{-}format}\NormalTok{ csv}
\end{Highlighting}
\end{Shaded}

\subsection{7 s2p Two-proton separation
S2p}\label{s2p-two-proton-separation-s2p}

Atomic-convention two-proton separation energy

Use and what it reveals: Inspect two-proton thresholds and isotonic
trends.

\[
S_{2p}=M(Z-2,N)+2H-M(Z,N)
\]

Output unit: keV. Inputs are all masses in the state-consistent stencil
below; missing terms make the result unavailable.

Menu or code aliases: \texttt{S2p}.

\begin{longtable}[]{@{}rrr@{}}
\toprule\noalign{}
ΔZ & ΔN & Coefficient \\
\midrule\noalign{}
\endhead
\bottomrule\noalign{}
\endlastfoot
-2 & 0 & 1 \\
0 & 0 & -1 \\
\end{longtable}

Light-particle terms

\begin{longtable}[]{@{}lr@{}}
\toprule\noalign{}
Particle & Coefficient \\
\midrule\noalign{}
\endhead
\bottomrule\noalign{}
\endlastfoot
H & 2 \\
\end{longtable}

Constant offset:

\[
C=0
\]

Convention: Neutral atomic mass excess, ground states, fixed Z/N
coordinate convention.

Source index: \texttt{ame2020}, \texttt{filters-implementation}.

\subsubsection{Checkable example}\label{checkable-example-6}

\textbf{116Sn} · Calculated from retained source data

16088.987128000013 ± 0.29221909704192844 keV

Default evaluated reference: 16088.987 ± 0.2899 keV

\href{https://www-nds.iaea.org/amdc/ame2020/rct1.mas20.txt}{Original
source table}

Recalculated and evaluated values/errors are distinct; source rounding
is retained. Full primitive inputs, flags and hashes are in this entry's
worked\_example in inventory-delivery.json.

Uncertainty: use the primitive-identity combination rule above. Keep
directly evaluated reference errors separate from independently
recalculated errors; model uncertainty is unavailable.

\subsubsection{Program mapping}\label{program-mapping-6}

\begin{itemize}
\tightlist
\item
  energies.py ·
  \texttt{Surface.get\ /\ Surface.\_linear\ /\ Reference.published} ·
  line 35
\item
  mass-surface.js · \texttt{create.derive} · line 28
\end{itemize}

\begin{Shaded}
\begin{Highlighting}[]
\ExtensionTok{./Run\_private.command} \AttributeTok{{-}m}\NormalTok{ traplab\_nuclear quantity 116Sn s2p }\AttributeTok{{-}{-}format}\NormalTok{ csv}
\end{Highlighting}
\end{Shaded}

\subsection{8 qbm Q(beta-)}\label{qbm-qbeta-}

Total atomic beta-minus energy allowance

Use and what it reveals: Check energetics; positive Q does not predict a
branch, endpoint to an excited level, or rate.

\[
Q_{\beta^-}=M(Z,N)-M(Z+1,N-1)
\]

Output unit: keV. Inputs are all masses in the state-consistent stencil
below; missing terms make the result unavailable.

\begin{longtable}[]{@{}rrr@{}}
\toprule\noalign{}
ΔZ & ΔN & Coefficient \\
\midrule\noalign{}
\endhead
\bottomrule\noalign{}
\endlastfoot
0 & 0 & 1 \\
1 & -1 & -1 \\
\end{longtable}

Constant offset:

\[
C=0
\]

Convention: Neutral atomic mass excess, ground states, fixed Z/N
coordinate convention.

Source index: \texttt{ame2020}, \texttt{filters-implementation}.

\subsubsection{Checkable example}\label{checkable-example-7}

\textbf{116Sn} · Calculated from retained source data

-4703.959000000003 ± 5.154893985330833 keV

Default evaluated reference: -4703.9591 ± 5.154 keV

\href{https://www-nds.iaea.org/amdc/ame2020/mass_1.mas20.txt}{Original
source table}

Recalculated and evaluated values/errors are distinct; source rounding
is retained. Full primitive inputs, flags and hashes are in this entry's
worked\_example in inventory-delivery.json.

Uncertainty: use the primitive-identity combination rule above. Keep
directly evaluated reference errors separate from independently
recalculated errors; model uncertainty is unavailable.

\subsubsection{Program mapping}\label{program-mapping-7}

\begin{itemize}
\tightlist
\item
  energies.py ·
  \texttt{Surface.get\ /\ Surface.\_linear\ /\ Reference.published} ·
  line 35
\item
  mass-surface.js · \texttt{create.derive} · line 29
\end{itemize}

\begin{Shaded}
\begin{Highlighting}[]
\ExtensionTok{./Run\_private.command} \AttributeTok{{-}m}\NormalTok{ traplab\_nuclear quantity 116Sn qbm }\AttributeTok{{-}{-}format}\NormalTok{ csv}
\end{Highlighting}
\end{Shaded}

\subsection{9 qec Q(EC)}\label{qec-qec}

Total atomic electron-capture energy allowance

Use and what it reveals: Subtract shell-hole and daughter-excitation
energies separately for a capture-specific neutrino energy.

\[
Q_{\mathrm{EC}}=M(Z,N)-M(Z-1,N+1)
\]

Output unit: keV. Inputs are all masses in the state-consistent stencil
below; missing terms make the result unavailable.

\begin{longtable}[]{@{}rrr@{}}
\toprule\noalign{}
ΔZ & ΔN & Coefficient \\
\midrule\noalign{}
\endhead
\bottomrule\noalign{}
\endlastfoot
0 & 0 & 1 \\
-1 & 1 & -1 \\
\end{longtable}

Constant offset:

\[
C=0
\]

Convention: Neutral atomic mass excess, ground states, fixed Z/N
coordinate convention.

Source index: \texttt{ame2020}, \texttt{filters-implementation}.

\subsubsection{Checkable example}\label{checkable-example-8}

\textbf{116Sn} · Calculated from retained source data

-3276.221000000005 ± 0.24003333101883997 keV

Default evaluated reference: -3276.2204 ± 0.2397 keV

\href{https://www-nds.iaea.org/amdc/ame2020/mass_1.mas20.txt}{Original
source table}

Recalculated and evaluated values/errors are distinct; source rounding
is retained. Full primitive inputs, flags and hashes are in this entry's
worked\_example in inventory-delivery.json.

Uncertainty: use the primitive-identity combination rule above. Keep
directly evaluated reference errors separate from independently
recalculated errors; model uncertainty is unavailable.

\subsubsection{Program mapping}\label{program-mapping-8}

\begin{itemize}
\tightlist
\item
  energies.py ·
  \texttt{Surface.get\ /\ Surface.\_linear\ /\ Reference.published} ·
  line 35
\item
  mass-surface.js · \texttt{create.derive} · line 29
\end{itemize}

\begin{Shaded}
\begin{Highlighting}[]
\ExtensionTok{./Run\_private.command} \AttributeTok{{-}m}\NormalTok{ traplab\_nuclear quantity 116Sn qec }\AttributeTok{{-}{-}format}\NormalTok{ csv}
\end{Highlighting}
\end{Shaded}

\subsection{10 qa Q(alpha)}\label{qa-qalpha}

Neutral-atom alpha-decay energy allowance

Use and what it reveals: Compare alpha energetics, excluding recoil
partition and daughter excitation.

\[
Q_\alpha=M(Z,N)-M(Z-2,N-2)-\alpha
\]

Output unit: keV. Inputs are all masses in the state-consistent stencil
below; missing terms make the result unavailable.

\begin{longtable}[]{@{}rrr@{}}
\toprule\noalign{}
ΔZ & ΔN & Coefficient \\
\midrule\noalign{}
\endhead
\bottomrule\noalign{}
\endlastfoot
0 & 0 & 1 \\
-2 & -2 & -1 \\
\end{longtable}

Light-particle terms

\begin{longtable}[]{@{}lr@{}}
\toprule\noalign{}
Particle & Coefficient \\
\midrule\noalign{}
\endhead
\bottomrule\noalign{}
\endlastfoot
alpha & -1 \\
\end{longtable}

Constant offset:

\[
C=0
\]

Convention: Neutral atomic mass excess, ground states, fixed Z/N
coordinate convention.

Source index: \texttt{ame2020}, \texttt{filters-implementation}.

\subsubsection{Checkable example}\label{checkable-example-9}

\textbf{116Sn} · Calculated from retained source data

-3376.037870000003 ± 0.26779847367003423 keV

Default evaluated reference: -3376.0374 ± 0.2664 keV

\href{https://www-nds.iaea.org/amdc/ame2020/rct1.mas20.txt}{Original
source table}

Recalculated and evaluated values/errors are distinct; source rounding
is retained. Full primitive inputs, flags and hashes are in this entry's
worked\_example in inventory-delivery.json.

Uncertainty: use the primitive-identity combination rule above. Keep
directly evaluated reference errors separate from independently
recalculated errors; model uncertainty is unavailable.

\subsubsection{Program mapping}\label{program-mapping-9}

\begin{itemize}
\tightlist
\item
  energies.py ·
  \texttt{Surface.get\ /\ Surface.\_linear\ /\ Reference.published} ·
  line 35
\item
  mass-surface.js · \texttt{create.derive} · line 29
\end{itemize}

\begin{Shaded}
\begin{Highlighting}[]
\ExtensionTok{./Run\_private.command} \AttributeTok{{-}m}\NormalTok{ traplab\_nuclear quantity 116Sn qa }\AttributeTok{{-}{-}format}\NormalTok{ csv}
\end{Highlighting}
\end{Shaded}

\subsection{11 q2bm Q(2beta-)}\label{q2bm-q2beta-}

Total double-beta-minus energy allowance

Use and what it reveals: Compare double-beta candidate ground-state
differences; no matrix element or rate is calculated.

\[
Q_{2\beta^-}=M(Z,N)-M(Z+2,N-2)
\]

Output unit: keV. Inputs are all masses in the state-consistent stencil
below; missing terms make the result unavailable.

\begin{longtable}[]{@{}rrr@{}}
\toprule\noalign{}
ΔZ & ΔN & Coefficient \\
\midrule\noalign{}
\endhead
\bottomrule\noalign{}
\endlastfoot
0 & 0 & 1 \\
2 & -2 & -1 \\
\end{longtable}

Constant offset:

\[
C=0
\]

Convention: Neutral atomic mass excess, ground states, fixed Z/N
coordinate convention.

Source index: \texttt{ame2020}, \texttt{filters-implementation}.

\subsubsection{Checkable example}\label{checkable-example-10}

\textbf{116Sn} · Calculated from retained source data

-6262.1860000000015 ± 24.2061903652764 keV

Default evaluated reference: -6262.1863 ± 24.206 keV

\href{https://www-nds.iaea.org/amdc/ame2020/rct1.mas20.txt}{Original
source table}

Recalculated and evaluated values/errors are distinct; source rounding
is retained. Full primitive inputs, flags and hashes are in this entry's
worked\_example in inventory-delivery.json.

Uncertainty: use the primitive-identity combination rule above. Keep
directly evaluated reference errors separate from independently
recalculated errors; model uncertainty is unavailable.

\subsubsection{Program mapping}\label{program-mapping-10}

\begin{itemize}
\tightlist
\item
  energies.py ·
  \texttt{Surface.get\ /\ Surface.\_linear\ /\ Reference.published} ·
  line 36
\item
  mass-surface.js · \texttt{create.derive} · line 35
\end{itemize}

\begin{Shaded}
\begin{Highlighting}[]
\ExtensionTok{./Run\_private.command} \AttributeTok{{-}m}\NormalTok{ traplab\_nuclear quantity 116Sn q2bm }\AttributeTok{{-}{-}format}\NormalTok{ csv}
\end{Highlighting}
\end{Shaded}

\subsection{12 q2ec Q(2EC)}\label{q2ec-q2ec}

Total double-electron-capture atomic energy allowance

Use and what it reveals: Capture shells, vacancies and final excitation
are additional physics.

\[
Q_{2{\mathrm{EC}}}=M(Z,N)-M(Z-2,N+2)
\]

Output unit: keV. Inputs are all masses in the state-consistent stencil
below; missing terms make the result unavailable.

\begin{longtable}[]{@{}rrr@{}}
\toprule\noalign{}
ΔZ & ΔN & Coefficient \\
\midrule\noalign{}
\endhead
\bottomrule\noalign{}
\endlastfoot
0 & 0 & 1 \\
-2 & 2 & -1 \\
\end{longtable}

Constant offset:

\[
C=0
\]

Convention: Neutral atomic mass excess, ground states, fixed Z/N
coordinate convention.

Source index: \texttt{ame2020}, \texttt{filters-implementation}.

\subsubsection{Checkable example}\label{checkable-example-11}

\textbf{116Sn} · Calculated from retained source data

-2813.4900000000052 ± 0.1865904606350496 keV

Default evaluated reference: -2813.4899 ± 0.1295 keV

\href{https://www-nds.iaea.org/amdc/ame2020/rct1.mas20.txt}{Original
source table}

Recalculated and evaluated values/errors are distinct; source rounding
is retained. Full primitive inputs, flags and hashes are in this entry's
worked\_example in inventory-delivery.json.

Uncertainty: use the primitive-identity combination rule above. Keep
directly evaluated reference errors separate from independently
recalculated errors; model uncertainty is unavailable.

\subsubsection{Program mapping}\label{program-mapping-11}

\begin{itemize}
\tightlist
\item
  energies.py ·
  \texttt{Surface.get\ /\ Surface.\_linear\ /\ Reference.published} ·
  line 36
\item
  mass-surface.js · \texttt{create.derive} · line 36
\end{itemize}

\begin{Shaded}
\begin{Highlighting}[]
\ExtensionTok{./Run\_private.command} \AttributeTok{{-}m}\NormalTok{ traplab\_nuclear quantity 116Sn q2ec }\AttributeTok{{-}{-}format}\NormalTok{ csv}
\end{Highlighting}
\end{Shaded}

\subsection{13 qecbp Q(EC beta+)}\label{qecbp-qec-beta}

Mixed electron-capture/positron threshold

Use and what it reveals: Includes two fixed electron rest energies; this
is not single-beta-plus Q.

\[
Q_{{\mathrm{EC}}\beta^+}=M(Z,N)-M(Z-2,N+2)-2e
\]

Output unit: keV. Inputs are all masses in the state-consistent stencil
below; missing terms make the result unavailable.

\begin{longtable}[]{@{}rrr@{}}
\toprule\noalign{}
ΔZ & ΔN & Coefficient \\
\midrule\noalign{}
\endhead
\bottomrule\noalign{}
\endlastfoot
0 & 0 & 1 \\
-2 & 2 & -1 \\
\end{longtable}

Constant offset:

\[
C=-2e
\]

Convention: Neutral atomic mass excess, ground states, fixed Z/N
coordinate convention.

Source index: \texttt{ame2020}, \texttt{filters-implementation}.

\subsubsection{Checkable example}\label{checkable-example-12}

\textbf{116Sn} · Calculated from retained source data

-3835.4878999613634 ± 0.1865904606350496 keV

Default evaluated reference: -3835.487799961358 ± 0.1295 keV

\href{https://www-nds.iaea.org/amdc/ame2020/rct1.mas20.txt}{Original
source table}

Recalculated and evaluated values/errors are distinct; source rounding
is retained. Full primitive inputs, flags and hashes are in this entry's
worked\_example in inventory-delivery.json.

Uncertainty: use the primitive-identity combination rule above. Keep
directly evaluated reference errors separate from independently
recalculated errors; model uncertainty is unavailable.

\subsubsection{Program mapping}\label{program-mapping-12}

\begin{itemize}
\tightlist
\item
  energies.py ·
  \texttt{Surface.get\ /\ Surface.\_linear\ /\ Reference.published} ·
  line 36
\item
  mass-surface.js · \texttt{create.derive} · line 37
\end{itemize}

\begin{Shaded}
\begin{Highlighting}[]
\ExtensionTok{./Run\_private.command} \AttributeTok{{-}m}\NormalTok{ traplab\_nuclear quantity 116Sn qecbp }\AttributeTok{{-}{-}format}\NormalTok{ csv}
\end{Highlighting}
\end{Shaded}

\subsection{14 q2bp Q(2beta+)}\label{q2bp-q2beta}

Double-positron threshold

Use and what it reveals: Includes four fixed electron rest energies.

\[
Q_{2\beta^+}=M(Z,N)-M(Z-2,N+2)-4e
\]

Output unit: keV. Inputs are all masses in the state-consistent stencil
below; missing terms make the result unavailable.

\begin{longtable}[]{@{}rrr@{}}
\toprule\noalign{}
ΔZ & ΔN & Coefficient \\
\midrule\noalign{}
\endhead
\bottomrule\noalign{}
\endlastfoot
0 & 0 & 1 \\
-2 & 2 & -1 \\
\end{longtable}

Constant offset:

\[
C=-4e
\]

Convention: Neutral atomic mass excess, ground states, fixed Z/N
coordinate convention.

Source index: \texttt{ame2020}, \texttt{filters-implementation}.

\subsubsection{Checkable example}\label{checkable-example-13}

\textbf{116Sn} · Calculated from retained source data

-4857.4857999227215 ± 0.1865904606350496 keV

Default evaluated reference: -4857.485699922717 ± 0.1295 keV

\href{https://www-nds.iaea.org/amdc/ame2020/rct1.mas20.txt}{Original
source table}

Recalculated and evaluated values/errors are distinct; source rounding
is retained. Full primitive inputs, flags and hashes are in this entry's
worked\_example in inventory-delivery.json.

Uncertainty: use the primitive-identity combination rule above. Keep
directly evaluated reference errors separate from independently
recalculated errors; model uncertainty is unavailable.

\subsubsection{Program mapping}\label{program-mapping-13}

\begin{itemize}
\tightlist
\item
  energies.py ·
  \texttt{Surface.get\ /\ Surface.\_linear\ /\ Reference.published} ·
  line 36
\item
  mass-surface.js · \texttt{create.derive} · line 38
\end{itemize}

\begin{Shaded}
\begin{Highlighting}[]
\ExtensionTok{./Run\_private.command} \AttributeTok{{-}m}\NormalTok{ traplab\_nuclear quantity 116Sn q2bp }\AttributeTok{{-}{-}format}\NormalTok{ csv}
\end{Highlighting}
\end{Shaded}

\subsection{15 q4bm Q(4beta-)}\label{q4bm-q4beta-}

Formal quadruple-beta-minus atomic mass difference

Use and what it reveals: Check mass allowance only; not an observed-mode
or probability claim.

\[
Q_{4\beta^-}=M(Z,N)-M(Z+4,N-4)
\]

Output unit: keV. Inputs are all masses in the state-consistent stencil
below; missing terms make the result unavailable.

\begin{longtable}[]{@{}rrr@{}}
\toprule\noalign{}
ΔZ & ΔN & Coefficient \\
\midrule\noalign{}
\endhead
\bottomrule\noalign{}
\endlastfoot
0 & 0 & 1 \\
4 & -4 & -1 \\
\end{longtable}

Constant offset:

\[
C=0
\]

Convention: Neutral atomic mass excess, ground states, fixed Z/N
coordinate convention.

Source index: \texttt{ame2020}, \texttt{filters-implementation}.

\subsubsection{Checkable example}\label{checkable-example-14}

\textbf{116Sn} · Calculated from retained source data

-18479.102000000014 ± 13.017353993803809 keV

Default evaluated reference: -18479.1015 ± 13.0175 keV

\href{https://www-nds.iaea.org/amdc/ame2020/rct2_1.mas20.txt}{Original
source table}

Recalculated and evaluated values/errors are distinct; source rounding
is retained. Full primitive inputs, flags and hashes are in this entry's
worked\_example in inventory-delivery.json.

Uncertainty: use the primitive-identity combination rule above. Keep
directly evaluated reference errors separate from independently
recalculated errors; model uncertainty is unavailable.

\subsubsection{Program mapping}\label{program-mapping-14}

\begin{itemize}
\tightlist
\item
  energies.py ·
  \texttt{Surface.get\ /\ Surface.\_linear\ /\ Reference.published} ·
  line 36
\item
  mass-surface.js · \texttt{create.derive} · line 39
\end{itemize}

\begin{Shaded}
\begin{Highlighting}[]
\ExtensionTok{./Run\_private.command} \AttributeTok{{-}m}\NormalTok{ traplab\_nuclear quantity 116Sn q4bm }\AttributeTok{{-}{-}format}\NormalTok{ csv}
\end{Highlighting}
\end{Shaded}

\subsection{16 qep Q(EC p)}\label{qep-qec-p}

Electron-capture plus proton-emission atomic mass difference

Use and what it reveals: Check the combined final-state energy
allowance; no branching or sequential threshold model.

\[
Q_{{\mathrm{EC}}p}=M(Z,N)-M(Z-2,N+1)-H
\]

Output unit: keV. Inputs are all masses in the state-consistent stencil
below; missing terms make the result unavailable.

\begin{longtable}[]{@{}rrr@{}}
\toprule\noalign{}
ΔZ & ΔN & Coefficient \\
\midrule\noalign{}
\endhead
\bottomrule\noalign{}
\endlastfoot
0 & 0 & 1 \\
-2 & 1 & -1 \\
\end{longtable}

Light-particle terms

\begin{longtable}[]{@{}lr@{}}
\toprule\noalign{}
Particle & Coefficient \\
\midrule\noalign{}
\endhead
\bottomrule\noalign{}
\endlastfoot
H & -1 \\
\end{longtable}

Constant offset:

\[
C=0
\]

Convention: Neutral atomic mass excess, ground states, fixed Z/N
coordinate convention.

Source index: \texttt{ame2020}, \texttt{filters-implementation}.

\subsubsection{Checkable example}\label{checkable-example-15}

\textbf{116Sn} · Calculated from retained source data

-10730.470064000012 ± 0.6580402724522262 keV

Default evaluated reference: -10730.4702 ± 0.6573 keV

\href{https://www-nds.iaea.org/amdc/ame2020/rct1.mas20.txt}{Original
source table}

Recalculated and evaluated values/errors are distinct; source rounding
is retained. Full primitive inputs, flags and hashes are in this entry's
worked\_example in inventory-delivery.json.

Uncertainty: use the primitive-identity combination rule above. Keep
directly evaluated reference errors separate from independently
recalculated errors; model uncertainty is unavailable.

\subsubsection{Program mapping}\label{program-mapping-15}

\begin{itemize}
\tightlist
\item
  energies.py ·
  \texttt{Surface.get\ /\ Surface.\_linear\ /\ Reference.published} ·
  line 36
\item
  mass-surface.js · \texttt{create.derive} · line 40
\end{itemize}

\begin{Shaded}
\begin{Highlighting}[]
\ExtensionTok{./Run\_private.command} \AttributeTok{{-}m}\NormalTok{ traplab\_nuclear quantity 116Sn qep }\AttributeTok{{-}{-}format}\NormalTok{ csv}
\end{Highlighting}
\end{Shaded}

\subsection{17 qbmn Q(beta- n)}\label{qbmn-qbeta--n}

Beta-minus plus neutron-emission atomic mass difference

Use and what it reveals: Check a beta-delayed neutron energy allowance;
no feeding model.

\[
Q_{\beta^-n}=M(Z,N)-M(Z+1,N-2)-n
\]

Output unit: keV. Inputs are all masses in the state-consistent stencil
below; missing terms make the result unavailable.

\begin{longtable}[]{@{}rrr@{}}
\toprule\noalign{}
ΔZ & ΔN & Coefficient \\
\midrule\noalign{}
\endhead
\bottomrule\noalign{}
\endlastfoot
0 & 0 & 1 \\
1 & -2 & -1 \\
\end{longtable}

Light-particle terms

\begin{longtable}[]{@{}lr@{}}
\toprule\noalign{}
Particle & Coefficient \\
\midrule\noalign{}
\endhead
\bottomrule\noalign{}
\endlastfoot
n & -1 \\
\end{longtable}

Constant offset:

\[
C=0
\]

Convention: Neutral atomic mass excess, ground states, fixed Z/N
coordinate convention.

Source index: \texttt{ame2020}, \texttt{filters-implementation}.

\subsubsection{Checkable example}\label{checkable-example-16}

\textbf{116Sn} · Calculated from retained source data

-12593.88506000001 ± 16.025287554162638 keV

Default evaluated reference: -12593.8847 ± 16.0256 keV

\href{https://www-nds.iaea.org/amdc/ame2020/rct1.mas20.txt}{Original
source table}

Recalculated and evaluated values/errors are distinct; source rounding
is retained. Full primitive inputs, flags and hashes are in this entry's
worked\_example in inventory-delivery.json.

Uncertainty: use the primitive-identity combination rule above. Keep
directly evaluated reference errors separate from independently
recalculated errors; model uncertainty is unavailable.

\subsubsection{Program mapping}\label{program-mapping-16}

\begin{itemize}
\tightlist
\item
  energies.py ·
  \texttt{Surface.get\ /\ Surface.\_linear\ /\ Reference.published} ·
  line 36
\item
  mass-surface.js · \texttt{create.derive} · line 41
\end{itemize}

\begin{Shaded}
\begin{Highlighting}[]
\ExtensionTok{./Run\_private.command} \AttributeTok{{-}m}\NormalTok{ traplab\_nuclear quantity 116Sn qbmn }\AttributeTok{{-}{-}format}\NormalTok{ csv}
\end{Highlighting}
\end{Shaded}

\subsection{18 qda Q(d,alpha)}\label{qda-qdalpha}

Target(d,alpha) reaction Q

Use and what it reveals: Neutral-atom bookkeeping conserves mass number
and charge.

\[
Q_{(d,\alpha)}=M(Z,N)+d-M(Z-1,N-1)-\alpha
\]

Output unit: keV. Inputs are all masses in the state-consistent stencil
below; missing terms make the result unavailable.

\begin{longtable}[]{@{}rrr@{}}
\toprule\noalign{}
ΔZ & ΔN & Coefficient \\
\midrule\noalign{}
\endhead
\bottomrule\noalign{}
\endlastfoot
0 & 0 & 1 \\
-1 & -1 & -1 \\
\end{longtable}

Light-particle terms

\begin{longtable}[]{@{}lr@{}}
\toprule\noalign{}
Particle & Coefficient \\
\midrule\noalign{}
\endhead
\bottomrule\noalign{}
\endlastfoot
d & 1 \\
alpha & -1 \\
\end{longtable}

Constant offset:

\[
C=0
\]

Convention: Neutral atomic mass excess, ground states, fixed Z/N
coordinate convention.

Source index: \texttt{ame2020}, \texttt{filters-implementation}.

\subsubsection{Checkable example}\label{checkable-example-17}

\textbf{116Sn} · Calculated from retained source data

7754.635024999995 ± 0.3159383210770735 keV

Default evaluated reference: 7754.6353 ± 0.3154 keV

\href{https://www-nds.iaea.org/amdc/ame2020/rct2_1.mas20.txt}{Original
source table}

Recalculated and evaluated values/errors are distinct; source rounding
is retained. Full primitive inputs, flags and hashes are in this entry's
worked\_example in inventory-delivery.json.

Uncertainty: use the primitive-identity combination rule above. Keep
directly evaluated reference errors separate from independently
recalculated errors; model uncertainty is unavailable.

\subsubsection{Program mapping}\label{program-mapping-17}

\begin{itemize}
\tightlist
\item
  energies.py ·
  \texttt{Surface.get\ /\ Surface.\_linear\ /\ Reference.published} ·
  line 36
\item
  mass-surface.js · \texttt{create.derive} · line 45
\end{itemize}

\begin{Shaded}
\begin{Highlighting}[]
\ExtensionTok{./Run\_private.command} \AttributeTok{{-}m}\NormalTok{ traplab\_nuclear quantity 116Sn qda }\AttributeTok{{-}{-}format}\NormalTok{ csv}
\end{Highlighting}
\end{Shaded}

\subsection{19 qpa Q(p,alpha)}\label{qpa-qpalpha}

Target(p,alpha) reaction Q

Use and what it reveals: The projectile hydrogen atom and alpha atom
provide the consistent electron bookkeeping.

\[
Q_{(p,\alpha)}=M(Z,N)+H-M(Z-1,N-2)-\alpha
\]

Output unit: keV. Inputs are all masses in the state-consistent stencil
below; missing terms make the result unavailable.

\begin{longtable}[]{@{}rrr@{}}
\toprule\noalign{}
ΔZ & ΔN & Coefficient \\
\midrule\noalign{}
\endhead
\bottomrule\noalign{}
\endlastfoot
0 & 0 & 1 \\
-1 & -2 & -1 \\
\end{longtable}

Light-particle terms

\begin{longtable}[]{@{}lr@{}}
\toprule\noalign{}
Particle & Coefficient \\
\midrule\noalign{}
\endhead
\bottomrule\noalign{}
\endlastfoot
H & 1 \\
alpha & -1 \\
\end{longtable}

Constant offset:

\[
C=0
\]

Convention: Neutral atomic mass excess, ground states, fixed Z/N
coordinate convention.

Source index: \texttt{ame2020}, \texttt{filters-implementation}.

\subsubsection{Checkable example}\label{checkable-example-18}

\textbf{116Sn} · Calculated from retained source data

2705.1991940000007 ± 0.21109245052583003 keV

Default evaluated reference: 2705.1998 ± 0.2103 keV

\href{https://www-nds.iaea.org/amdc/ame2020/rct2_1.mas20.txt}{Original
source table}

Recalculated and evaluated values/errors are distinct; source rounding
is retained. Full primitive inputs, flags and hashes are in this entry's
worked\_example in inventory-delivery.json.

Uncertainty: use the primitive-identity combination rule above. Keep
directly evaluated reference errors separate from independently
recalculated errors; model uncertainty is unavailable.

\subsubsection{Program mapping}\label{program-mapping-18}

\begin{itemize}
\tightlist
\item
  energies.py ·
  \texttt{Surface.get\ /\ Surface.\_linear\ /\ Reference.published} ·
  line 36
\item
  mass-surface.js · \texttt{create.derive} · line 46
\end{itemize}

\begin{Shaded}
\begin{Highlighting}[]
\ExtensionTok{./Run\_private.command} \AttributeTok{{-}m}\NormalTok{ traplab\_nuclear quantity 116Sn qpa }\AttributeTok{{-}{-}format}\NormalTok{ csv}
\end{Highlighting}
\end{Shaded}

\subsection{20 qna Q(n,alpha)}\label{qna-qnalpha}

Target(n,alpha) reaction Q

Use and what it reveals: The neutron and alpha mass excesses are
explicit light-particle inputs.

\[
Q_{(n,\alpha)}=M(Z,N)+n-M(Z-2,N-1)-\alpha
\]

Output unit: keV. Inputs are all masses in the state-consistent stencil
below; missing terms make the result unavailable.

\begin{longtable}[]{@{}rrr@{}}
\toprule\noalign{}
ΔZ & ΔN & Coefficient \\
\midrule\noalign{}
\endhead
\bottomrule\noalign{}
\endlastfoot
0 & 0 & 1 \\
-2 & -1 & -1 \\
\end{longtable}

Light-particle terms

\begin{longtable}[]{@{}lr@{}}
\toprule\noalign{}
Particle & Coefficient \\
\midrule\noalign{}
\endhead
\bottomrule\noalign{}
\endlastfoot
n & 1 \\
alpha & -1 \\
\end{longtable}

Constant offset:

\[
C=0
\]

Convention: Neutral atomic mass excess, ground states, fixed Z/N
coordinate convention.

Source index: \texttt{ame2020}, \texttt{filters-implementation}.

\subsubsection{Checkable example}\label{checkable-example-19}

\textbf{116Sn} · Calculated from retained source data

3163.709189999986 ± 0.263137257149192 keV

Default evaluated reference: 3163.7098 ± 0.2602 keV

\href{https://www-nds.iaea.org/amdc/ame2020/rct2_1.mas20.txt}{Original
source table}

Recalculated and evaluated values/errors are distinct; source rounding
is retained. Full primitive inputs, flags and hashes are in this entry's
worked\_example in inventory-delivery.json.

Uncertainty: use the primitive-identity combination rule above. Keep
directly evaluated reference errors separate from independently
recalculated errors; model uncertainty is unavailable.

\subsubsection{Program mapping}\label{program-mapping-19}

\begin{itemize}
\tightlist
\item
  energies.py ·
  \texttt{Surface.get\ /\ Surface.\_linear\ /\ Reference.published} ·
  line 36
\item
  mass-surface.js · \texttt{create.derive} · line 47
\end{itemize}

\begin{Shaded}
\begin{Highlighting}[]
\ExtensionTok{./Run\_private.command} \AttributeTok{{-}m}\NormalTok{ traplab\_nuclear quantity 116Sn qna }\AttributeTok{{-}{-}format}\NormalTok{ csv}
\end{Highlighting}
\end{Shaded}

\subsection{21 d2n Neutron shell gap
delta2n}\label{d2n-neutron-shell-gap-delta2n}

Empirical two-neutron shell gap

Use and what it reveals: A local two-step curvature; interpretation as a
shell closure needs other structure evidence.

\[
\delta_{2n}=M(Z,N-2)-2M(Z,N)+M(Z,N+2)=S_{2n}(Z,N)-S_{2n}(Z,N+2)
\]

Output unit: keV. Inputs are all masses in the state-consistent stencil
below; missing terms make the result unavailable.

Menu or code aliases: \texttt{dS2n}.

\begin{longtable}[]{@{}rrr@{}}
\toprule\noalign{}
ΔZ & ΔN & Coefficient \\
\midrule\noalign{}
\endhead
\bottomrule\noalign{}
\endlastfoot
0 & -2 & 1 \\
0 & 0 & -2 \\
0 & 2 & 1 \\
\end{longtable}

Constant offset:

\[
C=0
\]

Convention: Neutral atomic mass excess, ground states, fixed Z/N
coordinate convention.

Source index: \texttt{ame2020}, \texttt{filters-implementation}.

\subsubsection{Checkable example}\label{checkable-example-20}

\textbf{116Sn} · Calculated from retained source data

839.3800000000192 ± 0.5354493440093097 keV

Default evaluated reference: 839.3800000000192 ± 0.5354493440093097 keV

Recalculated and evaluated values/errors are distinct; source rounding
is retained. Full primitive inputs, flags and hashes are in this entry's
worked\_example in inventory-delivery.json.

Uncertainty: use the primitive-identity combination rule above. Keep
directly evaluated reference errors separate from independently
recalculated errors; model uncertainty is unavailable.

\subsubsection{Program mapping}\label{program-mapping-20}

\begin{itemize}
\tightlist
\item
  energies.py ·
  \texttt{Surface.get\ /\ Surface.\_linear\ /\ Reference.published} ·
  line 37
\item
  mass-surface.js · \texttt{create.derive} · line 49
\end{itemize}

\begin{Shaded}
\begin{Highlighting}[]
\ExtensionTok{./Run\_private.command} \AttributeTok{{-}m}\NormalTok{ traplab\_nuclear quantity 116Sn d2n }\AttributeTok{{-}{-}format}\NormalTok{ csv}
\end{Highlighting}
\end{Shaded}

\subsection{22 d2p Proton shell gap
delta2p}\label{d2p-proton-shell-gap-delta2p}

Empirical two-proton shell gap

Use and what it reveals: A local proton two-step curvature.

\[
\delta_{2p}=M(Z-2,N)-2M(Z,N)+M(Z+2,N)=S_{2p}(Z,N)-S_{2p}(Z+2,N)
\]

Output unit: keV. Inputs are all masses in the state-consistent stencil
below; missing terms make the result unavailable.

Menu or code aliases: \texttt{dS2p}.

\begin{longtable}[]{@{}rrr@{}}
\toprule\noalign{}
ΔZ & ΔN & Coefficient \\
\midrule\noalign{}
\endhead
\bottomrule\noalign{}
\endlastfoot
-2 & 0 & 1 \\
0 & 0 & -2 \\
2 & 0 & 1 \\
\end{longtable}

Constant offset:

\[
C=0
\]

Convention: Neutral atomic mass excess, ground states, fixed Z/N
coordinate convention.

Source index: \texttt{ame2020}, \texttt{filters-implementation}.

\subsubsection{Checkable example}\label{checkable-example-21}

\textbf{116Sn} · Calculated from retained source data

5346.266000000018 ± 18.309087251963163 keV

Default evaluated reference: 5346.266000000018 ± 18.309087251963163 keV

Recalculated and evaluated values/errors are distinct; source rounding
is retained. Full primitive inputs, flags and hashes are in this entry's
worked\_example in inventory-delivery.json.

Uncertainty: use the primitive-identity combination rule above. Keep
directly evaluated reference errors separate from independently
recalculated errors; model uncertainty is unavailable.

\subsubsection{Program mapping}\label{program-mapping-21}

\begin{itemize}
\tightlist
\item
  energies.py ·
  \texttt{Surface.get\ /\ Surface.\_linear\ /\ Reference.published} ·
  line 37
\item
  mass-surface.js · \texttt{create.derive} · line 55
\end{itemize}

\begin{Shaded}
\begin{Highlighting}[]
\ExtensionTok{./Run\_private.command} \AttributeTok{{-}m}\NormalTok{ traplab\_nuclear quantity 116Sn d2p }\AttributeTok{{-}{-}format}\NormalTok{ csv}
\end{Highlighting}
\end{Shaded}

\subsection{23 d3n Neutron OES delta3n}\label{d3n-neutron-oes-delta3n}

Signed three-point neutron odd-even staggering

Use and what it reveals: Apply the implemented (-1)\^{}N convention; no
absolute-value clamp.

\[
\Delta_n^{(3)}=\frac{(-1)^N}{2}[M(Z,N-1)-2M(Z,N)+M(Z,N+1)]
\]

Output unit: keV. Inputs are all masses in the state-consistent stencil
below; missing terms make the result unavailable.

Menu or code aliases: \texttt{D3n}.

\begin{longtable}[]{@{}rrr@{}}
\toprule\noalign{}
ΔZ & ΔN & Coefficient \\
\midrule\noalign{}
\endhead
\bottomrule\noalign{}
\endlastfoot
0 & -1 & 0.5 * (-1)\^{}N \\
0 & 0 & -1 * (-1)\^{}N \\
0 & 1 & 0.5 * (-1)\^{}N \\
\end{longtable}

Constant offset:

\[
C=0
\]

Convention: (-1)\^{}N multiplies the full stencil; signs are not
absolute-valued.

Convention: Neutral atomic mass excess, ground states, fixed Z/N
coordinate convention.

Source index: \texttt{ame2020}, \texttt{filters-implementation},
\texttt{pairing-context-satula1998}. Pairing literature is context, not
a claim of direct code attribution.

\subsubsection{Checkable example}\label{checkable-example-22}

\textbf{116Sn} · Calculated from retained source data

1310.185000000005 ± 0.25998942286177723 keV

Default evaluated reference: 1310.185000000005 ± 0.25998942286177723 keV

Recalculated and evaluated values/errors are distinct; source rounding
is retained. Full primitive inputs, flags and hashes are in this entry's
worked\_example in inventory-delivery.json.

Uncertainty: use the primitive-identity combination rule above. Keep
directly evaluated reference errors separate from independently
recalculated errors; model uncertainty is unavailable.

\subsubsection{Program mapping}\label{program-mapping-22}

\begin{itemize}
\tightlist
\item
  energies.py ·
  \texttt{Surface.get\ /\ Surface.\_linear\ /\ Reference.published} ·
  line 37
\item
  mass-surface.js · \texttt{create.derive} · line 51
\end{itemize}

\begin{Shaded}
\begin{Highlighting}[]
\ExtensionTok{./Run\_private.command} \AttributeTok{{-}m}\NormalTok{ traplab\_nuclear quantity 116Sn d3n }\AttributeTok{{-}{-}format}\NormalTok{ csv}
\end{Highlighting}
\end{Shaded}

\subsection{24 d3p Proton OES delta3p}\label{d3p-proton-oes-delta3p}

Signed three-point proton odd-even staggering

Use and what it reveals: Apply (-1)\^{}Z, not (-1)\^{}N.

\[
\Delta_p^{(3)}=\frac{(-1)^Z}{2}[M(Z-1,N)-2M(Z,N)+M(Z+1,N)]
\]

Output unit: keV. Inputs are all masses in the state-consistent stencil
below; missing terms make the result unavailable.

Menu or code aliases: \texttt{D3p}.

\begin{longtable}[]{@{}rrr@{}}
\toprule\noalign{}
ΔZ & ΔN & Coefficient \\
\midrule\noalign{}
\endhead
\bottomrule\noalign{}
\endlastfoot
-1 & 0 & 0.5 * (-1)\^{}Z \\
0 & 0 & -1 * (-1)\^{}Z \\
1 & 0 & 0.5 * (-1)\^{}Z \\
\end{longtable}

Constant offset:

\[
C=0
\]

Convention: (-1)\^{}Z multiplies the full stencil; signs are not
absolute-valued.

Convention: Neutral atomic mass excess, ground states, fixed Z/N
coordinate convention.

Source index: \texttt{ame2020}, \texttt{filters-implementation},
\texttt{pairing-context-satula1998}. Pairing literature is context, not
a claim of direct code attribution.

\subsubsection{Checkable example}\label{checkable-example-23}

\textbf{116Sn} · Calculated from retained source data

2438.0185000000056 ± 4.219596455823708 keV

Default evaluated reference: 2438.0185000000056 ± 4.219596455823708 keV

Recalculated and evaluated values/errors are distinct; source rounding
is retained. Full primitive inputs, flags and hashes are in this entry's
worked\_example in inventory-delivery.json.

Uncertainty: use the primitive-identity combination rule above. Keep
directly evaluated reference errors separate from independently
recalculated errors; model uncertainty is unavailable.

\subsubsection{Program mapping}\label{program-mapping-23}

\begin{itemize}
\tightlist
\item
  energies.py ·
  \texttt{Surface.get\ /\ Surface.\_linear\ /\ Reference.published} ·
  line 37
\item
  mass-surface.js · \texttt{create.derive} · line 52
\end{itemize}

\begin{Shaded}
\begin{Highlighting}[]
\ExtensionTok{./Run\_private.command} \AttributeTok{{-}m}\NormalTok{ traplab\_nuclear quantity 116Sn d3p }\AttributeTok{{-}{-}format}\NormalTok{ csv}
\end{Highlighting}
\end{Shaded}

\subsection{25 d5n Neutron OES delta5n}\label{d5n-neutron-oes-delta5n}

Signed five-point neutron odd-even staggering

Use and what it reveals: Suppress smooth low-order trends; the ME sign
is opposite to the BE stencil sign.

\[
\Delta_n^{(5)}=\frac{(-1)^N}{8}[-M(Z,N-2)+4M(Z,N-1)-6M(Z,N)+4M(Z,N+1)-M(Z,N+2)]
\]

Output unit: keV. Inputs are all masses in the state-consistent stencil
below; missing terms make the result unavailable.

\begin{longtable}[]{@{}rrr@{}}
\toprule\noalign{}
ΔZ & ΔN & Coefficient \\
\midrule\noalign{}
\endhead
\bottomrule\noalign{}
\endlastfoot
0 & -2 & -0.125 * (-1)\^{}N \\
0 & -1 & 0.5 * (-1)\^{}N \\
0 & 0 & -0.75 * (-1)\^{}N \\
0 & 1 & 0.5 * (-1)\^{}N \\
0 & 2 & -0.125 * (-1)\^{}N \\
\end{longtable}

Constant offset:

\[
C=0
\]

Convention: (-1)\^{}N multiplies the full stencil; signs are not
absolute-valued.

Convention: Neutral atomic mass excess, ground states, fixed Z/N
coordinate convention.

Source index: \texttt{ame2020}, \texttt{filters-implementation},
\texttt{pairing-context-satula1998}. Pairing literature is context, not
a claim of direct code attribution.

\subsubsection{Checkable example}\label{checkable-example-24}

\textbf{116Sn} · Calculated from retained source data

1205.2625000000062 ± 0.2597427212647161 keV

Default evaluated reference: 1205.2625000000062 ± 0.2597427212647161 keV

Recalculated and evaluated values/errors are distinct; source rounding
is retained. Full primitive inputs, flags and hashes are in this entry's
worked\_example in inventory-delivery.json.

Uncertainty: use the primitive-identity combination rule above. Keep
directly evaluated reference errors separate from independently
recalculated errors; model uncertainty is unavailable.

\subsubsection{Program mapping}\label{program-mapping-24}

\begin{itemize}
\tightlist
\item
  energies.py ·
  \texttt{Surface.get\ /\ Surface.\_linear\ /\ Reference.published} ·
  line 37
\item
  mass-surface.js · \texttt{create.derive} · line 53
\end{itemize}

\begin{Shaded}
\begin{Highlighting}[]
\ExtensionTok{./Run\_private.command} \AttributeTok{{-}m}\NormalTok{ traplab\_nuclear quantity 116Sn d5n }\AttributeTok{{-}{-}format}\NormalTok{ csv}
\end{Highlighting}
\end{Shaded}

\subsection{26 d5p Proton OES delta5p}\label{d5p-proton-oes-delta5p}

Signed five-point proton odd-even staggering

Use and what it reveals: Uses proton parity and neighbors at fixed N.

\[
\Delta_p^{(5)}=\frac{(-1)^Z}{8}[-M(Z-2,N)+4M(Z-1,N)-6M(Z,N)+4M(Z+1,N)-M(Z+2,N)]
\]

Output unit: keV. Inputs are all masses in the state-consistent stencil
below; missing terms make the result unavailable.

\begin{longtable}[]{@{}rrr@{}}
\toprule\noalign{}
ΔZ & ΔN & Coefficient \\
\midrule\noalign{}
\endhead
\bottomrule\noalign{}
\endlastfoot
-2 & 0 & -0.125 * (-1)\^{}Z \\
-1 & 0 & 0.5 * (-1)\^{}Z \\
0 & 0 & -0.75 * (-1)\^{}Z \\
1 & 0 & 0.5 * (-1)\^{}Z \\
2 & 0 & -0.125 * (-1)\^{}Z \\
\end{longtable}

Constant offset:

\[
C=0
\]

Convention: (-1)\^{}Z multiplies the full stencil; signs are not
absolute-valued.

Convention: Neutral atomic mass excess, ground states, fixed Z/N
coordinate convention.

Source index: \texttt{ame2020}, \texttt{filters-implementation},
\texttt{pairing-context-satula1998}. Pairing literature is context, not
a claim of direct code attribution.

\subsubsection{Checkable example}\label{checkable-example-25}

\textbf{116Sn} · Calculated from retained source data

1769.735250000007 ± 4.7998167217613625 keV

Default evaluated reference: 1769.735250000007 ± 4.7998167217613625 keV

Recalculated and evaluated values/errors are distinct; source rounding
is retained. Full primitive inputs, flags and hashes are in this entry's
worked\_example in inventory-delivery.json.

Uncertainty: use the primitive-identity combination rule above. Keep
directly evaluated reference errors separate from independently
recalculated errors; model uncertainty is unavailable.

\subsubsection{Program mapping}\label{program-mapping-25}

\begin{itemize}
\tightlist
\item
  energies.py ·
  \texttt{Surface.get\ /\ Surface.\_linear\ /\ Reference.published} ·
  line 37
\item
  mass-surface.js · \texttt{create.derive} · line 54
\end{itemize}

\begin{Shaded}
\begin{Highlighting}[]
\ExtensionTok{./Run\_private.command} \AttributeTok{{-}m}\NormalTok{ traplab\_nuclear quantity 116Sn d5p }\AttributeTok{{-}{-}format}\NormalTok{ csv}
\end{Highlighting}
\end{Shaded}

\subsection{27 vpn Parity-dependent
deltaVpn}\label{vpn-parity-dependent-deltavpn}

Parity-dependent average proton-neutron interaction indicator

Use and what it reveals: Select the physical branch using both Z and N
parity. This is an empirical mass difference.

\[
\delta V_{pn}(Z,N)=\sum_i b_i B(Z+\Delta Z_i,N+\Delta N_i)\quad\text{(four parity branches)}
\]

Output unit: keV. Inputs are all masses in the state-consistent stencil
below; missing terms make the result unavailable.

Constant offset:

\[
C=0
\]

Z even, N even

\[
\delta V_{pn}^{ee}=\frac14[-M(Z,N)+M(Z,N-2)+M(Z-2,N)-M(Z-2,N-2)]
\]

\begin{longtable}[]{@{}rrr@{}}
\toprule\noalign{}
ΔZ & ΔN & Coefficient \\
\midrule\noalign{}
\endhead
\bottomrule\noalign{}
\endlastfoot
-2 & -2 & -0.25 \\
-2 & 0 & 0.25 \\
0 & -2 & 0.25 \\
0 & 0 & -0.25 \\
\end{longtable}

Z even, N odd

\[
\delta V_{pn}^{eo}=\frac12[-M(Z,N)+M(Z,N-1)+M(Z-2,N)-M(Z-2,N-1)]
\]

\begin{longtable}[]{@{}rrr@{}}
\toprule\noalign{}
ΔZ & ΔN & Coefficient \\
\midrule\noalign{}
\endhead
\bottomrule\noalign{}
\endlastfoot
-2 & -1 & -0.5 \\
-2 & 0 & 0.5 \\
0 & -1 & 0.5 \\
0 & 0 & -0.5 \\
\end{longtable}

Z odd, N even

\[
\delta V_{pn}^{oe}=\frac12[-M(Z,N)+M(Z,N-2)+M(Z-1,N)-M(Z-1,N-2)]
\]

\begin{longtable}[]{@{}rrr@{}}
\toprule\noalign{}
ΔZ & ΔN & Coefficient \\
\midrule\noalign{}
\endhead
\bottomrule\noalign{}
\endlastfoot
-1 & -2 & -0.5 \\
-1 & 0 & 0.5 \\
0 & -2 & 0.5 \\
0 & 0 & -0.5 \\
\end{longtable}

Z odd, N odd

\[
\begin{aligned}
& \delta V_{pn}^{oo}=-M(Z,N)+M(Z,N-1) \\
& +M(Z-1,N)-M(Z-1,N-1)
\end{aligned}
\]

\begin{longtable}[]{@{}rrr@{}}
\toprule\noalign{}
ΔZ & ΔN & Coefficient \\
\midrule\noalign{}
\endhead
\bottomrule\noalign{}
\endlastfoot
-1 & -1 & -1 \\
-1 & 0 & 1 \\
0 & -1 & 1 \\
0 & 0 & -1 \\
\end{longtable}

Convention: Neutral atomic mass excess, ground states, fixed Z/N
coordinate convention.

Source index: \texttt{ame2020}, \texttt{filters-implementation},
\texttt{delta-vpn-zhang1989}, \texttt{delta-vpn-cakirli2005},
\texttt{delta-vpn-cakirli2005-erratum}, \texttt{delta-vpn-zhang1989},
\texttt{delta-vpn-cakirli2005}, \texttt{delta-vpn-cakirli2005-erratum}.

\subsubsection{Checkable example}\label{checkable-example-26}

\textbf{116Sn} · Calculated from retained source data

381.54175000000396 ± 0.09641479398930436 keV

Default evaluated reference: 381.54175000000396 ± 0.09641479398930436
keV

Recalculated and evaluated values/errors are distinct; source rounding
is retained. Full primitive inputs, flags and hashes are in this entry's
worked\_example in inventory-delivery.json.

Uncertainty: use the primitive-identity combination rule above. Keep
directly evaluated reference errors separate from independently
recalculated errors; model uncertainty is unavailable.

\subsubsection{Program mapping}\label{program-mapping-26}

\begin{itemize}
\tightlist
\item
  energies.py ·
  \texttt{Surface.get\ /\ Surface.\_linear\ /\ Reference.published} ·
  line 37
\item
  mass-surface.js · \texttt{create.derive} · line 60
\end{itemize}

\begin{Shaded}
\begin{Highlighting}[]
\ExtensionTok{./Run\_private.command} \AttributeTok{{-}m}\NormalTok{ traplab\_nuclear quantity 116Sn vpn }\AttributeTok{{-}{-}format}\NormalTok{ csv}
\end{Highlighting}
\end{Shaded}

\subsection{28 wig Wigner N-chain
indicator}\label{wig-wigner-n-chain-indicator}

Fixed-Z neutron-chain Wigner cusp indicator

Use and what it reveals: Compares vpn at N and N±2. It is not the native
legacy isobar W or R W1/W2/W3.

\[
W_N(Z,N)=\delta V_{pn}(Z,N)-\frac12[\delta V_{pn}(Z,N-2)+\delta V_{pn}(Z,N+2)]
\]

Output unit: keV. Inputs are all masses in the state-consistent stencil
below; missing terms make the result unavailable.

Constant offset:

\[
C=0
\]

Z even, N even

\[
\begin{aligned}
& 0.125M(Z-2,N-4)-0.375M(Z-2,N-2)+0.375M(Z-2,N) \\
& -0.125M(Z-2,N+2)-0.125M(Z,N-4)+0.375M(Z,N-2) \\
& -0.375M(Z,N)+0.125M(Z,N+2)
\end{aligned}
\]

\begin{longtable}[]{@{}rrr@{}}
\toprule\noalign{}
ΔZ & ΔN & Coefficient \\
\midrule\noalign{}
\endhead
\bottomrule\noalign{}
\endlastfoot
-2 & -4 & 0.125 \\
-2 & -2 & -0.375 \\
-2 & 0 & 0.375 \\
-2 & 2 & -0.125 \\
0 & -4 & -0.125 \\
0 & -2 & 0.375 \\
0 & 0 & -0.375 \\
0 & 2 & 0.125 \\
\end{longtable}

Z even, N odd

\[
\begin{aligned}
& 0.25M(Z-2,N-3)-0.25M(Z-2,N-2)-0.5M(Z-2,N-1) \\
& +0.5M(Z-2,N)+0.25M(Z-2,N+1)-0.25M(Z-2,N+2) \\
& -0.25M(Z,N-3)+0.25M(Z,N-2)+0.5M(Z,N-1) \\
& -0.5M(Z,N)-0.25M(Z,N+1)+0.25M(Z,N+2)
\end{aligned}
\]

\begin{longtable}[]{@{}rrr@{}}
\toprule\noalign{}
ΔZ & ΔN & Coefficient \\
\midrule\noalign{}
\endhead
\bottomrule\noalign{}
\endlastfoot
-2 & -3 & 0.25 \\
-2 & -2 & -0.25 \\
-2 & -1 & -0.5 \\
-2 & 0 & 0.5 \\
-2 & 1 & 0.25 \\
-2 & 2 & -0.25 \\
0 & -3 & -0.25 \\
0 & -2 & 0.25 \\
0 & -1 & 0.5 \\
0 & 0 & -0.5 \\
0 & 1 & -0.25 \\
0 & 2 & 0.25 \\
\end{longtable}

Z odd, N even

\[
\begin{aligned}
& 0.25M(Z-1,N-4)-0.75M(Z-1,N-2)+0.75M(Z-1,N) \\
& -0.25M(Z-1,N+2)-0.25M(Z,N-4)+0.75M(Z,N-2) \\
& -0.75M(Z,N)+0.25M(Z,N+2)
\end{aligned}
\]

\begin{longtable}[]{@{}rrr@{}}
\toprule\noalign{}
ΔZ & ΔN & Coefficient \\
\midrule\noalign{}
\endhead
\bottomrule\noalign{}
\endlastfoot
-1 & -4 & 0.25 \\
-1 & -2 & -0.75 \\
-1 & 0 & 0.75 \\
-1 & 2 & -0.25 \\
0 & -4 & -0.25 \\
0 & -2 & 0.75 \\
0 & 0 & -0.75 \\
0 & 2 & 0.25 \\
\end{longtable}

Z odd, N odd

\[
\begin{aligned}
& 0.5M(Z-1,N-3)-0.5M(Z-1,N-2)-M(Z-1,N-1) \\
& +M(Z-1,N)+0.5M(Z-1,N+1)-0.5M(Z-1,N+2) \\
& -0.5M(Z,N-3)+0.5M(Z,N-2)+M(Z,N-1) \\
& -M(Z,N)-0.5M(Z,N+1)+0.5M(Z,N+2)
\end{aligned}
\]

\begin{longtable}[]{@{}rrr@{}}
\toprule\noalign{}
ΔZ & ΔN & Coefficient \\
\midrule\noalign{}
\endhead
\bottomrule\noalign{}
\endlastfoot
-1 & -3 & 0.5 \\
-1 & -2 & -0.5 \\
-1 & -1 & -1 \\
-1 & 0 & 1 \\
-1 & 1 & 0.5 \\
-1 & 2 & -0.5 \\
0 & -3 & -0.5 \\
0 & -2 & 0.5 \\
0 & -1 & 1 \\
0 & 0 & -1 \\
0 & 1 & -0.5 \\
0 & 2 & 0.5 \\
\end{longtable}

Convention: Neutral atomic mass excess, ground states, fixed Z/N
coordinate convention.

Attribution: this entry follows the current implementation. Its name
does not prove identity with a similarly named quantity in a paper.

Source index: \texttt{ame2020}, \texttt{filters-implementation},
\texttt{delta-vpn-zhang1989}, \texttt{delta-vpn-cakirli2005},
\texttt{delta-vpn-cakirli2005-erratum}.

\subsubsection{Checkable example}\label{checkable-example-27}

\textbf{116Sn} · Calculated from retained source data

-6.846624999992855 ± 0.16974401683122736 keV

Default evaluated reference: -6.846624999992855 ± 0.16974401683122736
keV

Recalculated and evaluated values/errors are distinct; source rounding
is retained. Full primitive inputs, flags and hashes are in this entry's
worked\_example in inventory-delivery.json.

Uncertainty: use the primitive-identity combination rule above. Keep
directly evaluated reference errors separate from independently
recalculated errors; model uncertainty is unavailable.

\subsubsection{Program mapping}\label{program-mapping-27}

\begin{itemize}
\tightlist
\item
  energies.py ·
  \texttt{Surface.get\ /\ Surface.\_linear\ /\ Reference.published} ·
  line 37
\item
  mass-surface.js · \texttt{create.derive} · line 26
\end{itemize}

\begin{Shaded}
\begin{Highlighting}[]
\ExtensionTok{./Run\_private.command} \AttributeTok{{-}m}\NormalTok{ traplab\_nuclear quantity 116Sn wig }\AttributeTok{{-}{-}format}\NormalTok{ csv}
\end{Highlighting}
\end{Shaded}

\subsection{29 D1nS1n D1nS1n mass
filter}\label{d1ns1n-d1ns1n-mass-filter}

One-neutron finite-difference gap

Use and what it reveals: A one-step second difference; without parity
factor it is not the pairing indicator.

\[
S_n(Z,N)-S_n(Z,N+1)
\]

Output unit: keV. Inputs are all masses in the state-consistent stencil
below; missing terms make the result unavailable.

\begin{longtable}[]{@{}rrr@{}}
\toprule\noalign{}
ΔZ & ΔN & Coefficient \\
\midrule\noalign{}
\endhead
\bottomrule\noalign{}
\endlastfoot
0 & -1 & 1 \\
0 & 0 & -2 \\
0 & 1 & 1 \\
\end{longtable}

Constant offset:

\[
C=0
\]

Convention: Neutral atomic mass excess, ground states, fixed Z/N
coordinate convention.

Source index: \texttt{ame2020}, \texttt{filters-implementation}.

\subsubsection{Checkable example}\label{checkable-example-28}

\textbf{116Sn} · Calculated from retained source data

2620.37000000001 ± 0.5199788457235545 keV

Default evaluated reference: 2620.37000000001 ± 0.5199788457235545 keV

Recalculated and evaluated values/errors are distinct; source rounding
is retained. Full primitive inputs, flags and hashes are in this entry's
worked\_example in inventory-delivery.json.

Uncertainty: use the primitive-identity combination rule above. Keep
directly evaluated reference errors separate from independently
recalculated errors; model uncertainty is unavailable.

\subsubsection{Program mapping}\label{program-mapping-28}

\begin{itemize}
\tightlist
\item
  energies.py ·
  \texttt{Surface.get\ /\ Surface.\_linear\ /\ Reference.published} ·
  line 24
\item
  mass-surface.js · \texttt{create.derive} · line 8
\end{itemize}

\begin{Shaded}
\begin{Highlighting}[]
\ExtensionTok{./Run\_private.command} \AttributeTok{{-}m}\NormalTok{ traplab\_nuclear quantity 116Sn D1nS1n }\AttributeTok{{-}{-}format}\NormalTok{ csv}
\end{Highlighting}
\end{Shaded}

\subsection{30 D1pS1p D1pS1p mass
filter}\label{d1ps1p-d1ps1p-mass-filter}

One-proton finite-difference gap

Use and what it reveals: A one-step second difference along Z.

\[
S_p(Z,N)-S_p(Z+1,N)
\]

Output unit: keV. Inputs are all masses in the state-consistent stencil
below; missing terms make the result unavailable.

\begin{longtable}[]{@{}rrr@{}}
\toprule\noalign{}
ΔZ & ΔN & Coefficient \\
\midrule\noalign{}
\endhead
\bottomrule\noalign{}
\endlastfoot
-1 & 0 & 1 \\
0 & 0 & -2 \\
1 & 0 & 1 \\
\end{longtable}

Constant offset:

\[
C=0
\]

Convention: Neutral atomic mass excess, ground states, fixed Z/N
coordinate convention.

Source index: \texttt{ame2020}, \texttt{filters-implementation}.

\subsubsection{Checkable example}\label{checkable-example-29}

\textbf{116Sn} · Calculated from retained source data

4876.037000000011 ± 8.439192911647416 keV

Default evaluated reference: 4876.037000000011 ± 8.439192911647416 keV

Recalculated and evaluated values/errors are distinct; source rounding
is retained. Full primitive inputs, flags and hashes are in this entry's
worked\_example in inventory-delivery.json.

Uncertainty: use the primitive-identity combination rule above. Keep
directly evaluated reference errors separate from independently
recalculated errors; model uncertainty is unavailable.

\subsubsection{Program mapping}\label{program-mapping-29}

\begin{itemize}
\tightlist
\item
  energies.py ·
  \texttt{Surface.get\ /\ Surface.\_linear\ /\ Reference.published} ·
  line 25
\item
  mass-surface.js · \texttt{create.derive} · line 9
\end{itemize}

\begin{Shaded}
\begin{Highlighting}[]
\ExtensionTok{./Run\_private.command} \AttributeTok{{-}m}\NormalTok{ traplab\_nuclear quantity 116Sn D1pS1p }\AttributeTok{{-}{-}format}\NormalTok{ csv}
\end{Highlighting}
\end{Shaded}

\subsection{31 D1nS2n D1nS2n mass
filter}\label{d1ns2n-d1ns2n-mass-filter}

Forward slope of two-neutron separation energy

Use and what it reveals: Sign is next minus current, not current minus
next.

\[
S_{2n}(Z,N+1)-S_{2n}(Z,N)
\]

Output unit: keV. Inputs are all masses in the state-consistent stencil
below; missing terms make the result unavailable.

\begin{longtable}[]{@{}rrr@{}}
\toprule\noalign{}
ΔZ & ΔN & Coefficient \\
\midrule\noalign{}
\endhead
\bottomrule\noalign{}
\endlastfoot
0 & -2 & -1 \\
0 & -1 & 1 \\
0 & 0 & 1 \\
0 & 1 & -1 \\
\end{longtable}

Constant offset:

\[
C=0
\]

Convention: Neutral atomic mass excess, ground states, fixed Z/N
coordinate convention.

Source index: \texttt{ame2020}, \texttt{filters-implementation}.

\subsubsection{Checkable example}\label{checkable-example-30}

\textbf{116Sn} · Calculated from retained source data

-602.3480000000127 ± 0.49352912781314134 keV

Default evaluated reference: -602.3480000000127 ± 0.49352912781314134
keV

Recalculated and evaluated values/errors are distinct; source rounding
is retained. Full primitive inputs, flags and hashes are in this entry's
worked\_example in inventory-delivery.json.

Uncertainty: use the primitive-identity combination rule above. Keep
directly evaluated reference errors separate from independently
recalculated errors; model uncertainty is unavailable.

\subsubsection{Program mapping}\label{program-mapping-30}

\begin{itemize}
\tightlist
\item
  energies.py ·
  \texttt{Surface.get\ /\ Surface.\_linear\ /\ Reference.published} ·
  line 26
\item
  mass-surface.js · \texttt{create.derive} · line 10
\end{itemize}

\begin{Shaded}
\begin{Highlighting}[]
\ExtensionTok{./Run\_private.command} \AttributeTok{{-}m}\NormalTok{ traplab\_nuclear quantity 116Sn D1nS2n }\AttributeTok{{-}{-}format}\NormalTok{ csv}
\end{Highlighting}
\end{Shaded}

\subsection{32 D1pS2p D1pS2p mass
filter}\label{d1ps2p-d1ps2p-mass-filter}

Forward slope of two-proton separation energy

Use and what it reveals: Proton analogue of the forward S2n slope.

\[
S_{2p}(Z+1,N)-S_{2p}(Z,N)
\]

Output unit: keV. Inputs are all masses in the state-consistent stencil
below; missing terms make the result unavailable.

\begin{longtable}[]{@{}rrr@{}}
\toprule\noalign{}
ΔZ & ΔN & Coefficient \\
\midrule\noalign{}
\endhead
\bottomrule\noalign{}
\endlastfoot
-2 & 0 & -1 \\
-1 & 0 & 1 \\
0 & 0 & 1 \\
1 & 0 & -1 \\
\end{longtable}

Constant offset:

\[
C=0
\]

Convention: Neutral atomic mass excess, ground states, fixed Z/N
coordinate convention.

Source index: \texttt{ame2020}, \texttt{filters-implementation}.

\subsubsection{Checkable example}\label{checkable-example-31}

\textbf{116Sn} · Calculated from retained source data

-2407.838000000018 ± 8.442067578502318 keV

Default evaluated reference: -2407.838000000018 ± 8.442067578502318 keV

Recalculated and evaluated values/errors are distinct; source rounding
is retained. Full primitive inputs, flags and hashes are in this entry's
worked\_example in inventory-delivery.json.

Uncertainty: use the primitive-identity combination rule above. Keep
directly evaluated reference errors separate from independently
recalculated errors; model uncertainty is unavailable.

\subsubsection{Program mapping}\label{program-mapping-31}

\begin{itemize}
\tightlist
\item
  energies.py ·
  \texttt{Surface.get\ /\ Surface.\_linear\ /\ Reference.published} ·
  line 27
\item
  mass-surface.js · \texttt{create.derive} · line 11
\end{itemize}

\begin{Shaded}
\begin{Highlighting}[]
\ExtensionTok{./Run\_private.command} \AttributeTok{{-}m}\NormalTok{ traplab\_nuclear quantity 116Sn D1pS2p }\AttributeTok{{-}{-}format}\NormalTok{ csv}
\end{Highlighting}
\end{Shaded}

\subsection{33 D2pS2n D2pS2n mass
filter}\label{d2ps2n-d2ps2n-mass-filter}

Quarter-normalized two-proton/two-neutron mixed difference

Use and what it reveals: Equals the automatic vpn branch only at even Z
and even N.

\[
\frac14[S_{2n}(Z,N)-S_{2n}(Z-2,N)]
\]

Output unit: keV. Inputs are all masses in the state-consistent stencil
below; missing terms make the result unavailable.

Menu or code aliases: \texttt{dVpn}.

\begin{longtable}[]{@{}rrr@{}}
\toprule\noalign{}
ΔZ & ΔN & Coefficient \\
\midrule\noalign{}
\endhead
\bottomrule\noalign{}
\endlastfoot
0 & 0 & -0.25 \\
0 & -2 & 0.25 \\
-2 & 0 & 0.25 \\
-2 & -2 & -0.25 \\
\end{longtable}

Constant offset:

\[
C=0
\]

Convention: Neutral atomic mass excess, ground states, fixed Z/N
coordinate convention.

Source index: \texttt{ame2020}, \texttt{filters-implementation}.

\subsubsection{Checkable example}\label{checkable-example-32}

\textbf{116Sn} · Calculated from retained source data

381.54175000000396 ± 0.09641479398930436 keV

Default evaluated reference: 381.54175000000396 ± 0.09641479398930436
keV

Recalculated and evaluated values/errors are distinct; source rounding
is retained. Full primitive inputs, flags and hashes are in this entry's
worked\_example in inventory-delivery.json.

Uncertainty: use the primitive-identity combination rule above. Keep
directly evaluated reference errors separate from independently
recalculated errors; model uncertainty is unavailable.

\subsubsection{Program mapping}\label{program-mapping-32}

\begin{itemize}
\tightlist
\item
  energies.py ·
  \texttt{Surface.get\ /\ Surface.\_linear\ /\ Reference.published} ·
  line 28
\item
  mass-surface.js · \texttt{create.derive} · line 12
\end{itemize}

\begin{Shaded}
\begin{Highlighting}[]
\ExtensionTok{./Run\_private.command} \AttributeTok{{-}m}\NormalTok{ traplab\_nuclear quantity 116Sn D2pS2n }\AttributeTok{{-}{-}format}\NormalTok{ csv}
\end{Highlighting}
\end{Shaded}

\subsection{34 D1pS1n D1pS1n mass
filter}\label{d1ps1n-d1ps1n-mass-filter}

One-proton/one-neutron mixed difference

Use and what it reveals: Equals automatic vpn at odd Z and odd N.

\[
S_n(Z,N)-S_n(Z-1,N)
\]

Output unit: keV. Inputs are all masses in the state-consistent stencil
below; missing terms make the result unavailable.

\begin{longtable}[]{@{}rrr@{}}
\toprule\noalign{}
ΔZ & ΔN & Coefficient \\
\midrule\noalign{}
\endhead
\bottomrule\noalign{}
\endlastfoot
0 & 0 & -1 \\
0 & -1 & 1 \\
-1 & 0 & 1 \\
-1 & -1 & -1 \\
\end{longtable}

Constant offset:

\[
C=0
\]

Convention: Neutral atomic mass excess, ground states, fixed Z/N
coordinate convention.

Source index: \texttt{ame2020}, \texttt{filters-implementation}.

\subsubsection{Checkable example}\label{checkable-example-33}

\textbf{116Sn} · Calculated from retained source data

525.5829999999987 ± 0.31652172121356853 keV

Default evaluated reference: 525.5829999999987 ± 0.31652172121356853 keV

Recalculated and evaluated values/errors are distinct; source rounding
is retained. Full primitive inputs, flags and hashes are in this entry's
worked\_example in inventory-delivery.json.

Uncertainty: use the primitive-identity combination rule above. Keep
directly evaluated reference errors separate from independently
recalculated errors; model uncertainty is unavailable.

\subsubsection{Program mapping}\label{program-mapping-33}

\begin{itemize}
\tightlist
\item
  energies.py ·
  \texttt{Surface.get\ /\ Surface.\_linear\ /\ Reference.published} ·
  line 29
\item
  mass-surface.js · \texttt{create.derive} · line 13
\end{itemize}

\begin{Shaded}
\begin{Highlighting}[]
\ExtensionTok{./Run\_private.command} \AttributeTok{{-}m}\NormalTok{ traplab\_nuclear quantity 116Sn D1pS1n }\AttributeTok{{-}{-}format}\NormalTok{ csv}
\end{Highlighting}
\end{Shaded}

\subsection{35 D1pS2n D1pS2n mass
filter}\label{d1ps2n-d1ps2n-mass-filter}

Half-normalized one-proton/two-neutron mixed difference

Use and what it reveals: Equals automatic vpn at odd Z and even N.

\[
\frac12[S_{2n}(Z,N)-S_{2n}(Z-1,N)]
\]

Output unit: keV. Inputs are all masses in the state-consistent stencil
below; missing terms make the result unavailable.

\begin{longtable}[]{@{}rrr@{}}
\toprule\noalign{}
ΔZ & ΔN & Coefficient \\
\midrule\noalign{}
\endhead
\bottomrule\noalign{}
\endlastfoot
0 & 0 & -0.5 \\
0 & -2 & 0.5 \\
-1 & 0 & 0.5 \\
-1 & -2 & -0.5 \\
\end{longtable}

Constant offset:

\[
C=0
\]

Convention: Neutral atomic mass excess, ground states, fixed Z/N
coordinate convention.

Source index: \texttt{ame2020}, \texttt{filters-implementation}.

\subsubsection{Checkable example}\label{checkable-example-34}

\textbf{116Sn} · Calculated from retained source data

398.50500000000466 ± 0.10670637281812179 keV

Default evaluated reference: 398.50500000000466 ± 0.10670637281812179
keV

Recalculated and evaluated values/errors are distinct; source rounding
is retained. Full primitive inputs, flags and hashes are in this entry's
worked\_example in inventory-delivery.json.

Uncertainty: use the primitive-identity combination rule above. Keep
directly evaluated reference errors separate from independently
recalculated errors; model uncertainty is unavailable.

\subsubsection{Program mapping}\label{program-mapping-34}

\begin{itemize}
\tightlist
\item
  energies.py ·
  \texttt{Surface.get\ /\ Surface.\_linear\ /\ Reference.published} ·
  line 30
\item
  mass-surface.js · \texttt{create.derive} · line 14
\end{itemize}

\begin{Shaded}
\begin{Highlighting}[]
\ExtensionTok{./Run\_private.command} \AttributeTok{{-}m}\NormalTok{ traplab\_nuclear quantity 116Sn D1pS2n }\AttributeTok{{-}{-}format}\NormalTok{ csv}
\end{Highlighting}
\end{Shaded}

\subsection{36 D2pS1n D2pS1n mass
filter}\label{d2ps1n-d2ps1n-mass-filter}

Half-normalized two-proton/one-neutron mixed difference

Use and what it reveals: Equals automatic vpn at even Z and odd N; dZ is
-2.

\[
\frac12[S_n(Z,N)-S_n(Z-2,N)]
\]

Output unit: keV. Inputs are all masses in the state-consistent stencil
below; missing terms make the result unavailable.

\begin{longtable}[]{@{}rrr@{}}
\toprule\noalign{}
ΔZ & ΔN & Coefficient \\
\midrule\noalign{}
\endhead
\bottomrule\noalign{}
\endlastfoot
0 & 0 & -0.5 \\
0 & -1 & 0.5 \\
-2 & 0 & 0.5 \\
-2 & -1 & -0.5 \\
\end{longtable}

Constant offset:

\[
C=0
\]

Convention: Neutral atomic mass excess, ground states, fixed Z/N
coordinate convention.

Source index: \texttt{ame2020}, \texttt{filters-implementation}.

\subsubsection{Checkable example}\label{checkable-example-35}

\textbf{116Sn} · Calculated from retained source data

260.2425000000003 ± 0.19081535577620581 keV

Default evaluated reference: 260.2425000000003 ± 0.19081535577620581 keV

Recalculated and evaluated values/errors are distinct; source rounding
is retained. Full primitive inputs, flags and hashes are in this entry's
worked\_example in inventory-delivery.json.

Uncertainty: use the primitive-identity combination rule above. Keep
directly evaluated reference errors separate from independently
recalculated errors; model uncertainty is unavailable.

\subsubsection{Program mapping}\label{program-mapping-35}

\begin{itemize}
\tightlist
\item
  energies.py ·
  \texttt{Surface.get\ /\ Surface.\_linear\ /\ Reference.published} ·
  line 31
\item
  mass-surface.js · \texttt{create.derive} · line 15
\end{itemize}

\begin{Shaded}
\begin{Highlighting}[]
\ExtensionTok{./Run\_private.command} \AttributeTok{{-}m}\NormalTok{ traplab\_nuclear quantity 116Sn D2pS1n }\AttributeTok{{-}{-}format}\NormalTok{ csv}
\end{Highlighting}
\end{Shaded}

\subsection{37 Gplus Gplus mass filter}\label{gplus-gplus-mass-filter}

One-sided third mass difference on the two-neutron grid

Use and what it reveals: Removes a constant shell-gap background; needs
N−2 through N+4.

\[
\delta_{2n}(Z,N)-\delta_{2n}(Z,N+2)
\]

Output unit: keV. Inputs are all masses in the state-consistent stencil
below; missing terms make the result unavailable.

\begin{longtable}[]{@{}rrr@{}}
\toprule\noalign{}
ΔZ & ΔN & Coefficient \\
\midrule\noalign{}
\endhead
\bottomrule\noalign{}
\endlastfoot
0 & -2 & 1 \\
0 & 0 & -3 \\
0 & 2 & 3 \\
0 & 4 & -1 \\
\end{longtable}

Constant offset:

\[
C=0
\]

Convention: Neutral atomic mass excess, ground states, fixed Z/N
coordinate convention.

Attribution: this entry follows the current implementation. Its name
does not prove identity with a similarly named quantity in a paper.

Source index: \texttt{ame2020}, \texttt{filters-implementation}.

\subsubsection{Checkable example}\label{checkable-example-36}

\textbf{116Sn} · Calculated from retained source data

157.41400000004796 ± 1.7807846585143303 keV

Default evaluated reference: 157.41400000004796 ± 1.7807846585143303 keV

Recalculated and evaluated values/errors are distinct; source rounding
is retained. Full primitive inputs, flags and hashes are in this entry's
worked\_example in inventory-delivery.json.

Uncertainty: use the primitive-identity combination rule above. Keep
directly evaluated reference errors separate from independently
recalculated errors; model uncertainty is unavailable.

\subsubsection{Program mapping}\label{program-mapping-36}

\begin{itemize}
\tightlist
\item
  energies.py ·
  \texttt{Surface.get\ /\ Surface.\_linear\ /\ Reference.published} ·
  line 32
\item
  mass-surface.js · \texttt{create.derive} · line 16
\end{itemize}

\begin{Shaded}
\begin{Highlighting}[]
\ExtensionTok{./Run\_private.command} \AttributeTok{{-}m}\NormalTok{ traplab\_nuclear quantity 116Sn Gplus }\AttributeTok{{-}{-}format}\NormalTok{ csv}
\end{Highlighting}
\end{Shaded}

\subsection{38 Gsym Gsym mass filter}\label{gsym-gsym-mass-filter}

Half-normalized symmetric fourth mass difference

Use and what it reveals: Uses the executable half-normalization; needs
N−4 through N+4.

\[
\frac12[2\delta_{2n}(Z,N)-\delta_{2n}(Z,N-2)-\delta_{2n}(Z,N+2)]
\]

Output unit: keV. Inputs are all masses in the state-consistent stencil
below; missing terms make the result unavailable.

\begin{longtable}[]{@{}rrr@{}}
\toprule\noalign{}
ΔZ & ΔN & Coefficient \\
\midrule\noalign{}
\endhead
\bottomrule\noalign{}
\endlastfoot
0 & -4 & -0.5 \\
0 & -2 & 2 \\
0 & 0 & -3 \\
0 & 2 & 2 \\
0 & 4 & -0.5 \\
\end{longtable}

Constant offset:

\[
C=0
\]

Convention: Neutral atomic mass excess, ground states, fixed Z/N
coordinate convention.

Attribution: this entry follows the current implementation. Its name
does not prove identity with a similarly named quantity in a paper.

Source index: \texttt{ame2020}, \texttt{filters-implementation}.

\subsubsection{Checkable example}\label{checkable-example-37}

\textbf{116Sn} · Calculated from retained source data

29.176000000043132 ± 1.1469616384169088 keV

Default evaluated reference: 29.176000000043132 ± 1.1469616384169088 keV

Recalculated and evaluated values/errors are distinct; source rounding
is retained. Full primitive inputs, flags and hashes are in this entry's
worked\_example in inventory-delivery.json.

Uncertainty: use the primitive-identity combination rule above. Keep
directly evaluated reference errors separate from independently
recalculated errors; model uncertainty is unavailable.

\subsubsection{Program mapping}\label{program-mapping-37}

\begin{itemize}
\tightlist
\item
  energies.py ·
  \texttt{Surface.get\ /\ Surface.\_linear\ /\ Reference.published} ·
  line 33
\item
  mass-surface.js · \texttt{create.derive} · line 17
\end{itemize}

\begin{Shaded}
\begin{Highlighting}[]
\ExtensionTok{./Run\_private.command} \AttributeTok{{-}m}\NormalTok{ traplab\_nuclear quantity 116Sn Gsym }\AttributeTok{{-}{-}format}\NormalTok{ csv}
\end{Highlighting}
\end{Shaded}

\subsection{39 W Legacy isobar Wigner
indicator}\label{w-legacy-isobar-wigner-indicator}

Preserved native isobar Wigner cusp indicator

Use and what it reveals: Uses the fixed quarter-normalized V0 at three
isobaric centers; do not replace by parity-dependent vpn.

\[
W_{\mathrm{isobar}}(Z,N)=V_0(Z,N)-\frac12[V_0(Z-1,N+1)+V_0(Z+1,N-1)],\quad V_0={\mathrm{D2pS2n}}
\]

Output unit: keV. Inputs are all masses in the state-consistent stencil
below; missing terms make the result unavailable.

\begin{longtable}[]{@{}rrr@{}}
\toprule\noalign{}
ΔZ & ΔN & Coefficient \\
\midrule\noalign{}
\endhead
\bottomrule\noalign{}
\endlastfoot
-3 & -1 & 0.125 \\
-3 & 1 & -0.125 \\
-2 & -2 & -0.25 \\
-2 & 0 & 0.25 \\
-1 & -3 & 0.125 \\
-1 & -1 & -0.25 \\
-1 & 1 & 0.125 \\
0 & -2 & 0.25 \\
0 & 0 & -0.25 \\
1 & -3 & -0.125 \\
1 & -1 & 0.125 \\
\end{longtable}

Constant offset:

\[
C=0
\]

Convention: Neutral atomic mass excess, ground states, fixed Z/N
coordinate convention.

Source index: \texttt{ame2020}, \texttt{filters-implementation}.

\subsubsection{Checkable example}\label{checkable-example-38}

\textbf{116Sn} · Calculated from retained source data

-3.282374999997046 ± 2.6904594664527095 keV

Default evaluated reference: -3.282374999997046 ± 2.6904594664527095 keV

Recalculated and evaluated values/errors are distinct; source rounding
is retained. Full primitive inputs, flags and hashes are in this entry's
worked\_example in inventory-delivery.json.

Uncertainty: use the primitive-identity combination rule above. Keep
directly evaluated reference errors separate from independently
recalculated errors; model uncertainty is unavailable.

\subsubsection{Program mapping}\label{program-mapping-38}

\begin{itemize}
\tightlist
\item
  energies.py ·
  \texttt{Surface.get\ /\ Surface.\_linear\ /\ Reference.published} ·
  line 230
\end{itemize}

\begin{Shaded}
\begin{Highlighting}[]
\ExtensionTok{./Run\_private.command} \AttributeTok{{-}m}\NormalTok{ traplab\_nuclear quantity 116Sn W }\AttributeTok{{-}{-}format}\NormalTok{ csv}
\end{Highlighting}
\end{Shaded}

\section{Separate legacy recipe
appendix}\label{separate-legacy-recipe-appendix}

This appendix preserves the R-derived binding-energy recipes and their
individual code entry points. The gsi namespace prevents similarly named
recipes from being silently merged with the shared mass-excess kernel.
Recipes commonly output MeV while the shared kernel uses keV;
normalization, shifts and parity filters can also differ. Compare exact
stencils before comparing values.

\subsection{1 gsi:identity\_op Identity operator; binding
energy}\label{gsiidentity_op-identity-operator-binding-energy}

Identity operator; binding energy

Use: Does not do a thing (except maybe unit conversion): gives back the
binding energy

Binding-energy form

\[
10^{-3}[1B(Z,N)]
\]

Equivalent mass-excess form

\[
10^{-3}[Z H+N n-M(Z,N)]
\]

Output unit: MeV

Exact stencil \texttt{terms}

\begin{longtable}[]{@{}rrr@{}}
\toprule\noalign{}
ΔZ & ΔN & Coefficient \\
\midrule\noalign{}
\endhead
\bottomrule\noalign{}
\endlastfoot
0 & 0 & -0.001 \\
\end{longtable}

Exact stencil \texttt{be\_terms}

\begin{longtable}[]{@{}rrr@{}}
\toprule\noalign{}
ΔZ & ΔN & Coefficient \\
\midrule\noalign{}
\endhead
\bottomrule\noalign{}
\endlastfoot
0 & 0 & 1.0 \\
\end{longtable}

Exact stencil \texttt{light\_terms}

\begin{longtable}[]{@{}lr@{}}
\toprule\noalign{}
Particle & Coefficient \\
\midrule\noalign{}
\endhead
\bottomrule\noalign{}
\endlastfoot
H & 0.05 \\
n & 0.066 \\
\end{longtable}

offset\_latex:

\[
0
\]

\begin{itemize}
\item
  \texttt{N\_filter}: oe
\item
  \texttt{Z\_filter}: oe
\item
  \texttt{meaning}: e=even, o=odd, oe=both; applied by comparison()
  after calculation, not apply\_energy\_operator().
\item
  nuclear.py ·
  \texttt{operator\_catalog\ /\ apply\_energy\_operator\ /\ comparison}
  · line 147
\item
  nuclear\_operator\_weights.csv · `` · line 2
\item
  nuclear\_operator\_metadata.csv · `` · line 2
\end{itemize}

\textbf{Z=50, N=66} · Calculated from retained source data

988.6815266 ± 9.6e-05 MeV

Recalculated and evaluated values/errors are distinct; source rounding
is retained. Full primitive inputs, flags and hashes are in this entry's
worked\_example in inventory-delivery.json.

Uncertainty: combine repeated primitive identities and use the stated
covariance approximation. Marginal mass errors alone do not recover full
AME evaluated covariance.

\subsection{2 gsi:S1n One-neutron separation
energy}\label{gsis1n-one-neutron-separation-energy}

One-neutron separation energy

Use: Calculated by S1n(N,Z) = ME(N-1,Z)+ME\_n-ME(N,Z) =
BE(N,Z)-BE(N-1,Z).

Binding-energy form

\[
10^{-3}[-1B(Z,N-1)+1B(Z,N)]
\]

Equivalent mass-excess form

\[
0.001M(Z,N-1)-0.001M(Z,N) + 0.001 n
\]

Output unit: MeV

Exact stencil \texttt{terms}

\begin{longtable}[]{@{}rrr@{}}
\toprule\noalign{}
ΔZ & ΔN & Coefficient \\
\midrule\noalign{}
\endhead
\bottomrule\noalign{}
\endlastfoot
0 & -1 & 0.001 \\
0 & 0 & -0.001 \\
\end{longtable}

Exact stencil \texttt{be\_terms}

\begin{longtable}[]{@{}rrr@{}}
\toprule\noalign{}
ΔZ & ΔN & Coefficient \\
\midrule\noalign{}
\endhead
\bottomrule\noalign{}
\endlastfoot
0 & -1 & -1.0 \\
0 & 0 & 1.0 \\
\end{longtable}

Exact stencil \texttt{light\_terms}

\begin{longtable}[]{@{}lr@{}}
\toprule\noalign{}
Particle & Coefficient \\
\midrule\noalign{}
\endhead
\bottomrule\noalign{}
\endlastfoot
H & 0.0 \\
n & 0.001 \\
\end{longtable}

offset\_latex:

\[
0
\]

\begin{itemize}
\item
  \texttt{N\_filter}: oe
\item
  \texttt{Z\_filter}: oe
\item
  \texttt{meaning}: e=even, o=odd, oe=both; applied by comparison()
  after calculation, not apply\_energy\_operator().
\item
  nuclear.py ·
  \texttt{operator\_catalog\ /\ apply\_energy\_operator\ /\ comparison}
  · line 147
\item
  nuclear\_operator\_weights.csv · `` · line 3
\item
  nuclear\_operator\_metadata.csv · `` · line 3
\end{itemize}

\textbf{Z=50, N=66} · Calculated from retained source data

9.563451099999838 ± 9.716480844420989e-05 MeV

Recalculated and evaluated values/errors are distinct; source rounding
is retained. Full primitive inputs, flags and hashes are in this entry's
worked\_example in inventory-delivery.json.

Uncertainty: combine repeated primitive identities and use the stated
covariance approximation. Marginal mass errors alone do not recover full
AME evaluated covariance.

\subsection{3 gsi:S1p One-proton separation
energy}\label{gsis1p-one-proton-separation-energy}

One-proton separation energy

Use: Calculated by S1p(N,Z) = ME(N,Z-1)+ME\_p-ME(N,Z) =
BE(N,Z)-BE(N,Z-1).

Binding-energy form

\[
10^{-3}[1B(Z,N)-1B(Z-1,N)]
\]

Equivalent mass-excess form

\[
-0.001M(Z,N)+0.001M(Z-1,N) + 0.001 H
\]

Output unit: MeV

Exact stencil \texttt{terms}

\begin{longtable}[]{@{}rrr@{}}
\toprule\noalign{}
ΔZ & ΔN & Coefficient \\
\midrule\noalign{}
\endhead
\bottomrule\noalign{}
\endlastfoot
0 & 0 & -0.001 \\
-1 & 0 & 0.001 \\
\end{longtable}

Exact stencil \texttt{be\_terms}

\begin{longtable}[]{@{}rrr@{}}
\toprule\noalign{}
ΔZ & ΔN & Coefficient \\
\midrule\noalign{}
\endhead
\bottomrule\noalign{}
\endlastfoot
0 & 0 & 1.0 \\
-1 & 0 & -1.0 \\
\end{longtable}

Exact stencil \texttt{light\_terms}

\begin{longtable}[]{@{}lr@{}}
\toprule\noalign{}
Particle & Coefficient \\
\midrule\noalign{}
\endhead
\bottomrule\noalign{}
\endlastfoot
H & 0.001 \\
n & 0.0 \\
\end{longtable}

offset\_latex:

\[
0
\]

\begin{itemize}
\item
  \texttt{N\_filter}: oe
\item
  \texttt{Z\_filter}: oe
\item
  \texttt{meaning}: e=even, o=odd, oe=both; applied by comparison()
  after calculation, not apply\_energy\_operator().
\item
  nuclear.py ·
  \texttt{operator\_catalog\ /\ apply\_energy\_operator\ /\ comparison}
  · line 147
\item
  nuclear\_operator\_weights.csv · `` · line 5
\item
  nuclear\_operator\_metadata.csv · `` · line 4
\end{itemize}

\textbf{Z=50, N=66} · Calculated from retained source data

9.278593059999935 ± 9.67470929795826e-05 MeV

Recalculated and evaluated values/errors are distinct; source rounding
is retained. Full primitive inputs, flags and hashes are in this entry's
worked\_example in inventory-delivery.json.

Uncertainty: combine repeated primitive identities and use the stated
covariance approximation. Marginal mass errors alone do not recover full
AME evaluated covariance.

\subsection{4 gsi:S2n Two-neutron separation
energy}\label{gsis2n-two-neutron-separation-energy}

Two-neutron separation energy

Use: Calculated by S2n(N,Z) = ME(N-2,Z)+2*ME\_n-ME(N,Z) =
BE(N,Z)-BE(N-2,Z).

Binding-energy form

\[
10^{-3}[-1B(Z,N-2)+1B(Z,N)]
\]

Equivalent mass-excess form

\[
0.001M(Z,N-2)-0.001M(Z,N) + 0.002 n
\]

Output unit: MeV

Exact stencil \texttt{terms}

\begin{longtable}[]{@{}rrr@{}}
\toprule\noalign{}
ΔZ & ΔN & Coefficient \\
\midrule\noalign{}
\endhead
\bottomrule\noalign{}
\endlastfoot
0 & -2 & 0.001 \\
0 & 0 & -0.001 \\
\end{longtable}

Exact stencil \texttt{be\_terms}

\begin{longtable}[]{@{}rrr@{}}
\toprule\noalign{}
ΔZ & ΔN & Coefficient \\
\midrule\noalign{}
\endhead
\bottomrule\noalign{}
\endlastfoot
0 & -2 & -1.0 \\
0 & 0 & 1.0 \\
\end{longtable}

Exact stencil \texttt{light\_terms}

\begin{longtable}[]{@{}lr@{}}
\toprule\noalign{}
Particle & Coefficient \\
\midrule\noalign{}
\endhead
\bottomrule\noalign{}
\endlastfoot
H & 0.0 \\
n & 0.002 \\
\end{longtable}

offset\_latex:

\[
0
\]

\begin{itemize}
\item
  \texttt{N\_filter}: oe
\item
  \texttt{Z\_filter}: oe
\item
  \texttt{meaning}: e=even, o=odd, oe=both; applied by comparison()
  after calculation, not apply\_energy\_operator().
\item
  nuclear.py ·
  \texttt{operator\_catalog\ /\ apply\_energy\_operator\ /\ comparison}
  · line 147
\item
  nuclear\_operator\_weights.csv · `` · line 7
\item
  nuclear\_operator\_metadata.csv · `` · line 5
\end{itemize}

\textbf{Z=50, N=66} · Calculated from retained source data

17.10888019999993 ± 0.0001002845950283492 MeV

Recalculated and evaluated values/errors are distinct; source rounding
is retained. Full primitive inputs, flags and hashes are in this entry's
worked\_example in inventory-delivery.json.

Uncertainty: combine repeated primitive identities and use the stated
covariance approximation. Marginal mass errors alone do not recover full
AME evaluated covariance.

\subsection{5 gsi:S2p Two-proton separation
energy}\label{gsis2p-two-proton-separation-energy}

Two-proton separation energy

Use: Calculated by S2p(N,Z) = ME(N,Z-2)+2*ME\_p-ME(N,Z) =
BE(N,Z)-BE(N,Z-2).

Binding-energy form

\[
10^{-3}[1B(Z,N)-1B(Z-2,N)]
\]

Equivalent mass-excess form

\[
-0.001M(Z,N)+0.001M(Z-2,N) + 0.002 H
\]

Output unit: MeV

Exact stencil \texttt{terms}

\begin{longtable}[]{@{}rrr@{}}
\toprule\noalign{}
ΔZ & ΔN & Coefficient \\
\midrule\noalign{}
\endhead
\bottomrule\noalign{}
\endlastfoot
0 & 0 & -0.001 \\
-2 & 0 & 0.001 \\
\end{longtable}

Exact stencil \texttt{be\_terms}

\begin{longtable}[]{@{}rrr@{}}
\toprule\noalign{}
ΔZ & ΔN & Coefficient \\
\midrule\noalign{}
\endhead
\bottomrule\noalign{}
\endlastfoot
0 & 0 & 1.0 \\
-2 & 0 & -1.0 \\
\end{longtable}

Exact stencil \texttt{light\_terms}

\begin{longtable}[]{@{}lr@{}}
\toprule\noalign{}
Particle & Coefficient \\
\midrule\noalign{}
\endhead
\bottomrule\noalign{}
\endlastfoot
H & 0.002 \\
n & 0.0 \\
\end{longtable}

offset\_latex:

\[
0
\]

\begin{itemize}
\item
  \texttt{N\_filter}: oe
\item
  \texttt{Z\_filter}: oe
\item
  \texttt{meaning}: e=even, o=odd, oe=both; applied by comparison()
  after calculation, not apply\_energy\_operator().
\item
  nuclear.py ·
  \texttt{operator\_catalog\ /\ apply\_energy\_operator\ /\ comparison}
  · line 147
\item
  nuclear\_operator\_weights.csv · `` · line 9
\item
  nuclear\_operator\_metadata.csv · `` · line 6
\end{itemize}

\textbf{Z=50, N=66} · Calculated from retained source data

16.08898712000007 ± 0.00029221909588526216 MeV

Recalculated and evaluated values/errors are distinct; source rounding
is retained. Full primitive inputs, flags and hashes are in this entry's
worked\_example in inventory-delivery.json.

Uncertainty: combine repeated primitive identities and use the stated
covariance approximation. Marginal mass errors alone do not recover full
AME evaluated covariance.

\subsection{6 gsi:D1nS1n Empirical one-neutron shell
gap}\label{gsid1ns1n-empirical-one-neutron-shell-gap}

Empirical one-neutron shell gap

Use: Calculated by D1nS1n(N,Z) = S1n(N,Z)-S1n(N+1,Z) =
\{BE(N,Z)-BE(N-1,Z)\} - \{BE(N+1,Z)-BE(N,Z)\}.

Binding-energy form

\[
10^{-3}[-1B(Z,N-1)+2B(Z,N)-1B(Z,N+1)]
\]

Equivalent mass-excess form

\[
0.001M(Z,N-1)-0.002M(Z,N)+0.001M(Z,N+1)
\]

Output unit: MeV

Exact stencil \texttt{terms}

\begin{longtable}[]{@{}rrr@{}}
\toprule\noalign{}
ΔZ & ΔN & Coefficient \\
\midrule\noalign{}
\endhead
\bottomrule\noalign{}
\endlastfoot
0 & -1 & 0.001 \\
0 & 0 & -0.002 \\
0 & 1 & 0.001 \\
\end{longtable}

Exact stencil \texttt{be\_terms}

\begin{longtable}[]{@{}rrr@{}}
\toprule\noalign{}
ΔZ & ΔN & Coefficient \\
\midrule\noalign{}
\endhead
\bottomrule\noalign{}
\endlastfoot
0 & -1 & -1.0 \\
0 & 0 & 2.0 \\
0 & 1 & -1.0 \\
\end{longtable}

Exact stencil \texttt{light\_terms}

\begin{longtable}[]{@{}lr@{}}
\toprule\noalign{}
Particle & Coefficient \\
\midrule\noalign{}
\endhead
\bottomrule\noalign{}
\endlastfoot
H & 0.0 \\
n & 0.0 \\
\end{longtable}

offset\_latex:

\[
0
\]

\begin{itemize}
\item
  \texttt{N\_filter}: oe
\item
  \texttt{Z\_filter}: oe
\item
  \texttt{meaning}: e=even, o=odd, oe=both; applied by comparison()
  after calculation, not apply\_energy\_operator().
\item
  nuclear.py ·
  \texttt{operator\_catalog\ /\ apply\_energy\_operator\ /\ comparison}
  · line 147
\item
  nuclear\_operator\_weights.csv · `` · line 11
\item
  nuclear\_operator\_metadata.csv · `` · line 7
\end{itemize}

\textbf{Z=50, N=66} · Calculated from retained source data

2.620369999999866 ± 0.0005199788457235544 MeV

Recalculated and evaluated values/errors are distinct; source rounding
is retained. Full primitive inputs, flags and hashes are in this entry's
worked\_example in inventory-delivery.json.

Uncertainty: combine repeated primitive identities and use the stated
covariance approximation. Marginal mass errors alone do not recover full
AME evaluated covariance.

\subsection{7 gsi:D1nS2n The slope of the S2n
values}\label{gsid1ns2n-the-slope-of-the-s2n-values}

The slope of the S2n values

Use: Calculated by D1nS2n(N,Z) = (-1)\emph{\{S2n(N,Z)-S2n(N+1,Z)\} =
(-1)}\{{[}BE(N,Z)-BE(N-2,Z){]} - {[}BE(N+1,Z)-BE(N-1,Z){]}\}. See S.
Ayet PRC 99 064313 (2019).

Binding-energy form

\[
10^{-3}[1B(Z,N-2)-1B(Z,N-1)-1B(Z,N)+1B(Z,N+1)]
\]

Equivalent mass-excess form

\[
-0.001M(Z,N-2)+0.001M(Z,N-1)+0.001M(Z,N)-0.001M(Z,N+1)
\]

Output unit: MeV

Exact stencil \texttt{terms}

\begin{longtable}[]{@{}rrr@{}}
\toprule\noalign{}
ΔZ & ΔN & Coefficient \\
\midrule\noalign{}
\endhead
\bottomrule\noalign{}
\endlastfoot
0 & -2 & -0.001 \\
0 & -1 & 0.001 \\
0 & 0 & 0.001 \\
0 & 1 & -0.001 \\
\end{longtable}

Exact stencil \texttt{be\_terms}

\begin{longtable}[]{@{}rrr@{}}
\toprule\noalign{}
ΔZ & ΔN & Coefficient \\
\midrule\noalign{}
\endhead
\bottomrule\noalign{}
\endlastfoot
0 & -2 & 1.0 \\
0 & -1 & -1.0 \\
0 & 0 & -1.0 \\
0 & 1 & 1.0 \\
\end{longtable}

Exact stencil \texttt{light\_terms}

\begin{longtable}[]{@{}lr@{}}
\toprule\noalign{}
Particle & Coefficient \\
\midrule\noalign{}
\endhead
\bottomrule\noalign{}
\endlastfoot
H & 0.0 \\
n & 0.0 \\
\end{longtable}

offset\_latex:

\[
0
\]

\begin{itemize}
\item
  \texttt{N\_filter}: oe
\item
  \texttt{Z\_filter}: oe
\item
  \texttt{meaning}: e=even, o=odd, oe=both; applied by comparison()
  after calculation, not apply\_energy\_operator().
\item
  nuclear.py ·
  \texttt{operator\_catalog\ /\ apply\_energy\_operator\ /\ comparison}
  · line 147
\item
  nuclear\_operator\_weights.csv · `` · line 14
\item
  nuclear\_operator\_metadata.csv · `` · line 8
\end{itemize}

\textbf{Z=50, N=66} · Calculated from retained source data

-0.60234800000012 ± 0.0004935291278131413 MeV

Recalculated and evaluated values/errors are distinct; source rounding
is retained. Full primitive inputs, flags and hashes are in this entry's
worked\_example in inventory-delivery.json.

Uncertainty: combine repeated primitive identities and use the stated
covariance approximation. Marginal mass errors alone do not recover full
AME evaluated covariance.

\subsection{8 gsi:D2nS2n Empirical two-neutron shell
gap}\label{gsid2ns2n-empirical-two-neutron-shell-gap}

Empirical two-neutron shell gap

Use: Calculated by D2nS2n(N,Z) = S2n(N,Z)-S2n(N+2,Z) =
\{BE(N,Z)-BE(N-2,Z)\} - \{BE(N+2,Z)-BE(N,Z)\}.

Binding-energy form

\[
10^{-3}[-1B(Z,N-2)+2B(Z,N)-1B(Z,N+2)]
\]

Equivalent mass-excess form

\[
0.001M(Z,N-2)-0.002M(Z,N)+0.001M(Z,N+2)
\]

Output unit: MeV

Exact stencil \texttt{terms}

\begin{longtable}[]{@{}rrr@{}}
\toprule\noalign{}
ΔZ & ΔN & Coefficient \\
\midrule\noalign{}
\endhead
\bottomrule\noalign{}
\endlastfoot
0 & -2 & 0.001 \\
0 & 0 & -0.002 \\
0 & 2 & 0.001 \\
\end{longtable}

Exact stencil \texttt{be\_terms}

\begin{longtable}[]{@{}rrr@{}}
\toprule\noalign{}
ΔZ & ΔN & Coefficient \\
\midrule\noalign{}
\endhead
\bottomrule\noalign{}
\endlastfoot
0 & -2 & -1.0 \\
0 & 0 & 2.0 \\
0 & 2 & -1.0 \\
\end{longtable}

Exact stencil \texttt{light\_terms}

\begin{longtable}[]{@{}lr@{}}
\toprule\noalign{}
Particle & Coefficient \\
\midrule\noalign{}
\endhead
\bottomrule\noalign{}
\endlastfoot
H & 0.0 \\
n & 0.0 \\
\end{longtable}

offset\_latex:

\[
0
\]

\begin{itemize}
\item
  \texttt{N\_filter}: oe
\item
  \texttt{Z\_filter}: oe
\item
  \texttt{meaning}: e=even, o=odd, oe=both; applied by comparison()
  after calculation, not apply\_energy\_operator().
\item
  nuclear.py ·
  \texttt{operator\_catalog\ /\ apply\_energy\_operator\ /\ comparison}
  · line 147
\item
  nuclear\_operator\_weights.csv · `` · line 18
\item
  nuclear\_operator\_metadata.csv · `` · line 9
\end{itemize}

\textbf{Z=50, N=66} · Calculated from retained source data

0.8393799999998919 ± 0.0005354493440093097 MeV

Recalculated and evaluated values/errors are distinct; source rounding
is retained. Full primitive inputs, flags and hashes are in this entry's
worked\_example in inventory-delivery.json.

Uncertainty: combine repeated primitive identities and use the stated
covariance approximation. Marginal mass errors alone do not recover full
AME evaluated covariance.

\subsection{9 gsi:D1pS1p Empirical one-proton shell
gap}\label{gsid1ps1p-empirical-one-proton-shell-gap}

Empirical one-proton shell gap

Use: Calculated by D1pS1p(N,Z) = S1p(N,Z)-S1p(N,Z+1) =
\{BE(N,Z)-BE(N,Z-1)\} - \{BE(N,Z+1)-BE(N,Z)\}.

Binding-energy form

\[
10^{-3}[-1B(Z+1,N)+2B(Z,N)-1B(Z-1,N)]
\]

Equivalent mass-excess form

\[
0.001M(Z+1,N)-0.002M(Z,N)+0.001M(Z-1,N)
\]

Output unit: MeV

Exact stencil \texttt{terms}

\begin{longtable}[]{@{}rrr@{}}
\toprule\noalign{}
ΔZ & ΔN & Coefficient \\
\midrule\noalign{}
\endhead
\bottomrule\noalign{}
\endlastfoot
1 & 0 & 0.001 \\
0 & 0 & -0.002 \\
-1 & 0 & 0.001 \\
\end{longtable}

Exact stencil \texttt{be\_terms}

\begin{longtable}[]{@{}rrr@{}}
\toprule\noalign{}
ΔZ & ΔN & Coefficient \\
\midrule\noalign{}
\endhead
\bottomrule\noalign{}
\endlastfoot
1 & 0 & -1.0 \\
0 & 0 & 2.0 \\
-1 & 0 & -1.0 \\
\end{longtable}

Exact stencil \texttt{light\_terms}

\begin{longtable}[]{@{}lr@{}}
\toprule\noalign{}
Particle & Coefficient \\
\midrule\noalign{}
\endhead
\bottomrule\noalign{}
\endlastfoot
H & 0.0 \\
n & 0.0 \\
\end{longtable}

offset\_latex:

\[
0
\]

\begin{itemize}
\item
  \texttt{N\_filter}: oe
\item
  \texttt{Z\_filter}: oe
\item
  \texttt{meaning}: e=even, o=odd, oe=both; applied by comparison()
  after calculation, not apply\_energy\_operator().
\item
  nuclear.py ·
  \texttt{operator\_catalog\ /\ apply\_energy\_operator\ /\ comparison}
  · line 147
\item
  nuclear\_operator\_weights.csv · `` · line 21
\item
  nuclear\_operator\_metadata.csv · `` · line 10
\end{itemize}

\textbf{Z=50, N=66} · Calculated from retained source data

4.876036999999997 ± 0.008439192911647417 MeV

Recalculated and evaluated values/errors are distinct; source rounding
is retained. Full primitive inputs, flags and hashes are in this entry's
worked\_example in inventory-delivery.json.

Uncertainty: combine repeated primitive identities and use the stated
covariance approximation. Marginal mass errors alone do not recover full
AME evaluated covariance.

\subsection{10 gsi:D1pS2p The slope of the S2p
values}\label{gsid1ps2p-the-slope-of-the-s2p-values}

The slope of the S2p values

Use: Calculated by D1pS2p(N,Z) = (-1)\emph{\{S2p(N,Z)-S2p(N,Z+1)\} =
(-1)}\{{[}BE(N,Z)-BE(N,Z-2){]} - {[}BE(N,Z+1)-BE(N,Z-1){]}\}. See
analogue for neutrons in S. Ayet PRC 99 064313 (2019).

Binding-energy form

\[
10^{-3}[1B(Z+1,N)-1B(Z,N)-1B(Z-1,N)+1B(Z-2,N)]
\]

Equivalent mass-excess form

\[
-0.001M(Z+1,N)+0.001M(Z,N)+0.001M(Z-1,N)-0.001M(Z-2,N)
\]

Output unit: MeV

Exact stencil \texttt{terms}

\begin{longtable}[]{@{}rrr@{}}
\toprule\noalign{}
ΔZ & ΔN & Coefficient \\
\midrule\noalign{}
\endhead
\bottomrule\noalign{}
\endlastfoot
1 & 0 & -0.001 \\
0 & 0 & 0.001 \\
-1 & 0 & 0.001 \\
-2 & 0 & -0.001 \\
\end{longtable}

Exact stencil \texttt{be\_terms}

\begin{longtable}[]{@{}rrr@{}}
\toprule\noalign{}
ΔZ & ΔN & Coefficient \\
\midrule\noalign{}
\endhead
\bottomrule\noalign{}
\endlastfoot
1 & 0 & 1.0 \\
0 & 0 & -1.0 \\
-1 & 0 & -1.0 \\
-2 & 0 & 1.0 \\
\end{longtable}

Exact stencil \texttt{light\_terms}

\begin{longtable}[]{@{}lr@{}}
\toprule\noalign{}
Particle & Coefficient \\
\midrule\noalign{}
\endhead
\bottomrule\noalign{}
\endlastfoot
H & 0.0 \\
n & 0.0 \\
\end{longtable}

offset\_latex:

\[
0
\]

\begin{itemize}
\item
  \texttt{N\_filter}: oe
\item
  \texttt{Z\_filter}: oe
\item
  \texttt{meaning}: e=even, o=odd, oe=both; applied by comparison()
  after calculation, not apply\_energy\_operator().
\item
  nuclear.py ·
  \texttt{operator\_catalog\ /\ apply\_energy\_operator\ /\ comparison}
  · line 147
\item
  nuclear\_operator\_weights.csv · `` · line 24
\item
  nuclear\_operator\_metadata.csv · `` · line 11
\end{itemize}

\textbf{Z=50, N=66} · Calculated from retained source data

-2.407838000000197 ± 0.008442067578502318 MeV

Recalculated and evaluated values/errors are distinct; source rounding
is retained. Full primitive inputs, flags and hashes are in this entry's
worked\_example in inventory-delivery.json.

Uncertainty: combine repeated primitive identities and use the stated
covariance approximation. Marginal mass errors alone do not recover full
AME evaluated covariance.

\subsection{11 gsi:D2pS2p Empirical two-proton shell
gap}\label{gsid2ps2p-empirical-two-proton-shell-gap}

Empirical two-proton shell gap

Use: Calculated by D2pS2p(N,Z) = S2p(N,Z)-S2p(N,Z+2) =
\{BE(N,Z)-BE(N,Z-2)\} - \{BE(N,Z+2)-BE(N,Z)\}.

Binding-energy form

\[
10^{-3}[-1B(Z+2,N)+2B(Z,N)-1B(Z-2,N)]
\]

Equivalent mass-excess form

\[
0.001M(Z+2,N)-0.002M(Z,N)+0.001M(Z-2,N)
\]

Output unit: MeV

Exact stencil \texttt{terms}

\begin{longtable}[]{@{}rrr@{}}
\toprule\noalign{}
ΔZ & ΔN & Coefficient \\
\midrule\noalign{}
\endhead
\bottomrule\noalign{}
\endlastfoot
2 & 0 & 0.001 \\
0 & 0 & -0.002 \\
-2 & 0 & 0.001 \\
\end{longtable}

Exact stencil \texttt{be\_terms}

\begin{longtable}[]{@{}rrr@{}}
\toprule\noalign{}
ΔZ & ΔN & Coefficient \\
\midrule\noalign{}
\endhead
\bottomrule\noalign{}
\endlastfoot
2 & 0 & -1.0 \\
0 & 0 & 2.0 \\
-2 & 0 & -1.0 \\
\end{longtable}

Exact stencil \texttt{light\_terms}

\begin{longtable}[]{@{}lr@{}}
\toprule\noalign{}
Particle & Coefficient \\
\midrule\noalign{}
\endhead
\bottomrule\noalign{}
\endlastfoot
H & 0.0 \\
n & 0.0 \\
\end{longtable}

offset\_latex:

\[
0
\]

\begin{itemize}
\item
  \texttt{N\_filter}: oe
\item
  \texttt{Z\_filter}: oe
\item
  \texttt{meaning}: e=even, o=odd, oe=both; applied by comparison()
  after calculation, not apply\_energy\_operator().
\item
  nuclear.py ·
  \texttt{operator\_catalog\ /\ apply\_energy\_operator\ /\ comparison}
  · line 147
\item
  nuclear\_operator\_weights.csv · `` · line 28
\item
  nuclear\_operator\_metadata.csv · `` · line 12
\end{itemize}

\textbf{Z=50, N=66} · Calculated from retained source data

5.346266000000014 ± 0.018309087251963167 MeV

Recalculated and evaluated values/errors are distinct; source rounding
is retained. Full primitive inputs, flags and hashes are in this entry's
worked\_example in inventory-delivery.json.

Uncertainty: combine repeated primitive identities and use the stated
covariance approximation. Marginal mass errors alone do not recover full
AME evaluated covariance.

\subsection{12 gsi:D2pS2n dVnp\_ee: Two Proton two neutron interaction
strength}\label{gsid2ps2n-dvnp_ee-two-proton-two-neutron-interaction-strength}

dVnp\_ee: Two Proton two neutron interaction strength

Use: Calculated by D2pS2n(N,Z) = 1/4\emph{{[}S2n(N,Z)-S2n(N,Z-2){]} =
1/4}\{{[}BE(N,Z)-BE(N-2,Z){]} - {[}BE(N,Z-2)-BE(N-2,Z-2){]}\}.

Binding-energy form

\[
-0.25B(Z,N-2)+0.25B(Z,N)+0.25B(Z-2,N-2)-0.25B(Z-2,N)
\]

Equivalent mass-excess form

\[
0.25M(Z,N-2)-0.25M(Z,N)-0.25M(Z-2,N-2)+0.25M(Z-2,N)
\]

Output unit: keV

Exact stencil \texttt{terms}

\begin{longtable}[]{@{}rrr@{}}
\toprule\noalign{}
ΔZ & ΔN & Coefficient \\
\midrule\noalign{}
\endhead
\bottomrule\noalign{}
\endlastfoot
0 & -2 & 0.25 \\
0 & 0 & -0.25 \\
-2 & -2 & -0.25 \\
-2 & 0 & 0.25 \\
\end{longtable}

Exact stencil \texttt{be\_terms}

\begin{longtable}[]{@{}rrr@{}}
\toprule\noalign{}
ΔZ & ΔN & Coefficient \\
\midrule\noalign{}
\endhead
\bottomrule\noalign{}
\endlastfoot
0 & -2 & -0.25 \\
0 & 0 & 0.25 \\
-2 & -2 & 0.25 \\
-2 & 0 & -0.25 \\
\end{longtable}

Exact stencil \texttt{light\_terms}

\begin{longtable}[]{@{}lr@{}}
\toprule\noalign{}
Particle & Coefficient \\
\midrule\noalign{}
\endhead
\bottomrule\noalign{}
\endlastfoot
H & 0.0 \\
n & 0.0 \\
\end{longtable}

offset\_latex:

\[
0
\]

\begin{itemize}
\item
  \texttt{N\_filter}: e
\item
  \texttt{Z\_filter}: e
\item
  \texttt{meaning}: e=even, o=odd, oe=both; applied by comparison()
  after calculation, not apply\_energy\_operator().
\item
  nuclear.py ·
  \texttt{operator\_catalog\ /\ apply\_energy\_operator\ /\ comparison}
  · line 147
\item
  nuclear\_operator\_weights.csv · `` · line 31
\item
  nuclear\_operator\_metadata.csv · `` · line 13
\end{itemize}

\textbf{Z=50, N=66} · Calculated from retained source data

381.54175000000396 ± 0.09641479398930436 keV

Recalculated and evaluated values/errors are distinct; source rounding
is retained. Full primitive inputs, flags and hashes are in this entry's
worked\_example in inventory-delivery.json.

Uncertainty: combine repeated primitive identities and use the stated
covariance approximation. Marginal mass errors alone do not recover full
AME evaluated covariance.

\subsection{13 gsi:D1pS1n dVnp\_oo: One proton one neutron interaction
strength}\label{gsid1ps1n-dvnp_oo-one-proton-one-neutron-interaction-strength}

dVnp\_oo: One proton one neutron interaction strength

Use: Calculated by D1pS1n(N,Z) = {[}S1n(N,Z)-S1n(N,Z-1){]} =
\{{[}BE(N,Z)-BE(N-1,Z){]} - {[}BE(N,Z-1)-BE(N-1,Z-1){]}\}.

Binding-energy form

\[
-1B(Z,N-1)+1B(Z,N)+1B(Z-1,N-1)-1B(Z-1,N)
\]

Equivalent mass-excess form

\[
1M(Z,N-1)-1M(Z,N)-1M(Z-1,N-1)+1M(Z-1,N)
\]

Output unit: keV

Exact stencil \texttt{terms}

\begin{longtable}[]{@{}rrr@{}}
\toprule\noalign{}
ΔZ & ΔN & Coefficient \\
\midrule\noalign{}
\endhead
\bottomrule\noalign{}
\endlastfoot
0 & -1 & 1.0 \\
0 & 0 & -1.0 \\
-1 & -1 & -1.0 \\
-1 & 0 & 1.0 \\
\end{longtable}

Exact stencil \texttt{be\_terms}

\begin{longtable}[]{@{}rrr@{}}
\toprule\noalign{}
ΔZ & ΔN & Coefficient \\
\midrule\noalign{}
\endhead
\bottomrule\noalign{}
\endlastfoot
0 & -1 & -1.0 \\
0 & 0 & 1.0 \\
-1 & -1 & 1.0 \\
-1 & 0 & -1.0 \\
\end{longtable}

Exact stencil \texttt{light\_terms}

\begin{longtable}[]{@{}lr@{}}
\toprule\noalign{}
Particle & Coefficient \\
\midrule\noalign{}
\endhead
\bottomrule\noalign{}
\endlastfoot
H & 0.0 \\
n & 0.0 \\
\end{longtable}

offset\_latex:

\[
0
\]

\begin{itemize}
\item
  \texttt{N\_filter}: o
\item
  \texttt{Z\_filter}: o
\item
  \texttt{meaning}: e=even, o=odd, oe=both; applied by comparison()
  after calculation, not apply\_energy\_operator().
\item
  nuclear.py ·
  \texttt{operator\_catalog\ /\ apply\_energy\_operator\ /\ comparison}
  · line 147
\item
  nuclear\_operator\_weights.csv · `` · line 35
\item
  nuclear\_operator\_metadata.csv · `` · line 14
\end{itemize}

\textbf{Z=51, N=67} · Calculated from retained source data

484.877000000095 ± 8.973390106308763 keV

Recalculated and evaluated values/errors are distinct; source rounding
is retained. Full primitive inputs, flags and hashes are in this entry's
worked\_example in inventory-delivery.json.

Uncertainty: combine repeated primitive identities and use the stated
covariance approximation. Marginal mass errors alone do not recover full
AME evaluated covariance.

\subsection{14 gsi:D1pS2n dVnp\_eo: One proton two neutron interaction
strength}\label{gsid1ps2n-dvnp_eo-one-proton-two-neutron-interaction-strength}

dVnp\_eo: One proton two neutron interaction strength

Use: Calculated by D2pS2n(N,Z) = 1/2\emph{{[}S2n(N,Z)-S2n(N,Z-1){]} =
1/2}\{{[}BE(N,Z)-BE(N-2,Z){]} - {[}BE(N,Z-1)-BE(N-2,Z-1){]}\}.

Binding-energy form

\[
-0.5B(Z,N-2)+0.5B(Z,N)+0.5B(Z-1,N-2)-0.5B(Z-1,N)
\]

Equivalent mass-excess form

\[
0.5M(Z,N-2)-0.5M(Z,N)-0.5M(Z-1,N-2)+0.5M(Z-1,N)
\]

Output unit: keV

Exact stencil \texttt{terms}

\begin{longtable}[]{@{}rrr@{}}
\toprule\noalign{}
ΔZ & ΔN & Coefficient \\
\midrule\noalign{}
\endhead
\bottomrule\noalign{}
\endlastfoot
0 & -2 & 0.5 \\
0 & 0 & -0.5 \\
-1 & -2 & -0.5 \\
-1 & 0 & 0.5 \\
\end{longtable}

Exact stencil \texttt{be\_terms}

\begin{longtable}[]{@{}rrr@{}}
\toprule\noalign{}
ΔZ & ΔN & Coefficient \\
\midrule\noalign{}
\endhead
\bottomrule\noalign{}
\endlastfoot
0 & -2 & -0.5 \\
0 & 0 & 0.5 \\
-1 & -2 & 0.5 \\
-1 & 0 & -0.5 \\
\end{longtable}

Exact stencil \texttt{light\_terms}

\begin{longtable}[]{@{}lr@{}}
\toprule\noalign{}
Particle & Coefficient \\
\midrule\noalign{}
\endhead
\bottomrule\noalign{}
\endlastfoot
H & 0.0 \\
n & 0.0 \\
\end{longtable}

offset\_latex:

\[
0
\]

\begin{itemize}
\item
  \texttt{N\_filter}: e
\item
  \texttt{Z\_filter}: o
\item
  \texttt{meaning}: e=even, o=odd, oe=both; applied by comparison()
  after calculation, not apply\_energy\_operator().
\item
  nuclear.py ·
  \texttt{operator\_catalog\ /\ apply\_energy\_operator\ /\ comparison}
  · line 147
\item
  nuclear\_operator\_weights.csv · `` · line 39
\item
  nuclear\_operator\_metadata.csv · `` · line 15
\end{itemize}

\textbf{Z=51, N=66} · Calculated from retained source data

334.954000000027 ± 9.05529749649342 keV

Recalculated and evaluated values/errors are distinct; source rounding
is retained. Full primitive inputs, flags and hashes are in this entry's
worked\_example in inventory-delivery.json.

Uncertainty: combine repeated primitive identities and use the stated
covariance approximation. Marginal mass errors alone do not recover full
AME evaluated covariance.

\subsection{15 gsi:D2pS1n dVnp\_oe: Two proton one neutron interaction
strength}\label{gsid2ps1n-dvnp_oe-two-proton-one-neutron-interaction-strength}

dVnp\_oe: Two proton one neutron interaction strength

Use: Calculated by D2pS1n(N,Z) = 1/2\emph{{[}S1n(N,Z)-S1n(N,Z-2){]} =
1/2}\{{[}BE(N,Z)-BE(N-1,Z){]} - {[}BE(N,Z-2)-BE(N-1,Z-2){]}\}.

Binding-energy form

\[
-0.5B(Z,N-1)+0.5B(Z,N)+0.5B(Z-2,N-1)-0.5B(Z-2,N)
\]

Equivalent mass-excess form

\[
0.5M(Z,N-1)-0.5M(Z,N)-0.5M(Z-2,N-1)+0.5M(Z-2,N)
\]

Output unit: keV

Exact stencil \texttt{terms}

\begin{longtable}[]{@{}rrr@{}}
\toprule\noalign{}
ΔZ & ΔN & Coefficient \\
\midrule\noalign{}
\endhead
\bottomrule\noalign{}
\endlastfoot
0 & -1 & 0.5 \\
0 & 0 & -0.5 \\
-2 & -1 & -0.5 \\
-2 & 0 & 0.5 \\
\end{longtable}

Exact stencil \texttt{be\_terms}

\begin{longtable}[]{@{}rrr@{}}
\toprule\noalign{}
ΔZ & ΔN & Coefficient \\
\midrule\noalign{}
\endhead
\bottomrule\noalign{}
\endlastfoot
0 & -1 & -0.5 \\
0 & 0 & 0.5 \\
-2 & -1 & 0.5 \\
-2 & 0 & -0.5 \\
\end{longtable}

Exact stencil \texttt{light\_terms}

\begin{longtable}[]{@{}lr@{}}
\toprule\noalign{}
Particle & Coefficient \\
\midrule\noalign{}
\endhead
\bottomrule\noalign{}
\endlastfoot
H & 0.0 \\
n & 0.0 \\
\end{longtable}

offset\_latex:

\[
0
\]

\begin{itemize}
\item
  \texttt{N\_filter}: o
\item
  \texttt{Z\_filter}: e
\item
  \texttt{meaning}: e=even, o=odd, oe=both; applied by comparison()
  after calculation, not apply\_energy\_operator().
\item
  nuclear.py ·
  \texttt{operator\_catalog\ /\ apply\_energy\_operator\ /\ comparison}
  · line 147
\item
  nuclear\_operator\_weights.csv · `` · line 43
\item
  nuclear\_operator\_metadata.csv · `` · line 16
\end{itemize}

\textbf{Z=50, N=67} · Calculated from retained source data

401.108499999973 ± 0.43083697612902266 keV

Recalculated and evaluated values/errors are distinct; source rounding
is retained. Full primitive inputs, flags and hashes are in this entry's
worked\_example in inventory-delivery.json.

Uncertainty: combine repeated primitive identities and use the stated
covariance approximation. Marginal mass errors alone do not recover full
AME evaluated covariance.

\subsection{16 gsi:Gplus\_2n2n2n One-sided (neutron-deficient) 3rd order
mass
difference}\label{gsigplus_2n2n2n-one-sided-neutron-deficient-3rd-order-mass-difference}

One-sided (neutron-deficient) 3rd order mass difference

Use: Calculated by Gplus\_2n2n2n(N,Z) = D2nS2n(N,Z) - D2nS2n(N+2,Z) =
\{{[}BE(N,Z)-BE(N-2,Z){]}-{[}BE(N+2,Z)-BE(N,Z){]}\} -
\{{[}BE(N+2,Z)-BE(N,Z){]}-{[}BE(N+4,Z)-BE(N+2,Z){]}\}. One-sided
(neutron-deficient) 3rd order mass difference: the idea is to correct a
possible constant offset of \texttt{D2nSn}: can be used, where data with
less neutrons is unavailable.

Binding-energy form

\[
10^{-3}[-1B(Z,N-2)+3B(Z,N)-3B(Z,N+2)+1B(Z,N+4)]
\]

Equivalent mass-excess form

\[
0.001M(Z,N-2)-0.003M(Z,N)+0.003M(Z,N+2)-0.001M(Z,N+4)
\]

Output unit: MeV

Exact stencil \texttt{terms}

\begin{longtable}[]{@{}rrr@{}}
\toprule\noalign{}
ΔZ & ΔN & Coefficient \\
\midrule\noalign{}
\endhead
\bottomrule\noalign{}
\endlastfoot
0 & -2 & 0.001 \\
0 & 0 & -0.003 \\
0 & 2 & 0.003 \\
0 & 4 & -0.001 \\
\end{longtable}

Exact stencil \texttt{be\_terms}

\begin{longtable}[]{@{}rrr@{}}
\toprule\noalign{}
ΔZ & ΔN & Coefficient \\
\midrule\noalign{}
\endhead
\bottomrule\noalign{}
\endlastfoot
0 & -2 & -1.0 \\
0 & 0 & 3.0 \\
0 & 2 & -3.0 \\
0 & 4 & 1.0 \\
\end{longtable}

Exact stencil \texttt{light\_terms}

\begin{longtable}[]{@{}lr@{}}
\toprule\noalign{}
Particle & Coefficient \\
\midrule\noalign{}
\endhead
\bottomrule\noalign{}
\endlastfoot
H & 0.0 \\
n & 0.0 \\
\end{longtable}

offset\_latex:

\[
0
\]

\begin{itemize}
\item
  \texttt{N\_filter}: oe
\item
  \texttt{Z\_filter}: oe
\item
  \texttt{meaning}: e=even, o=odd, oe=both; applied by comparison()
  after calculation, not apply\_energy\_operator().
\item
  nuclear.py ·
  \texttt{operator\_catalog\ /\ apply\_energy\_operator\ /\ comparison}
  · line 147
\item
  nuclear\_operator\_weights.csv · `` · line 47
\item
  nuclear\_operator\_metadata.csv · `` · line 17
\end{itemize}

\textbf{Z=50, N=66} · Calculated from retained source data

0.15741399999990335 ± 0.0017807846585143304 MeV

Recalculated and evaluated values/errors are distinct; source rounding
is retained. Full primitive inputs, flags and hashes are in this entry's
worked\_example in inventory-delivery.json.

Uncertainty: combine repeated primitive identities and use the stated
covariance approximation. Marginal mass errors alone do not recover full
AME evaluated covariance.

\subsection{17 gsi:Gsym\_2n2n2n2n Two-sided (symmetrical) 4th order mass
difference}\label{gsigsym_2n2n2n2n-two-sided-symmetrical-4th-order-mass-difference}

Two-sided (symmetrical) 4th order mass difference

Use: Calculated by Gsym\_2n2n2n2n(N,Z) = 2\emph{D2nS2n(N,Z) -
D2nS2n(N+2,Z) - D2nS2n(N-2,Z) = =
2}\{{[}BE(N,Z)-BE(N-2,Z){]}-{[}BE(N+2,Z)-BE(N,Z){]}\} -
\{{[}BE(N+2,Z)-BE(N,Z){]}-{[}BE(N+4,Z)-BE(N+2,Z){]}\} -
\{{[}BE(N-2,Z)-BE(N-4,Z){]}-{[}BE(N,Z)-BE(N-2,Z){]}\}. Two-sided
(symmetrical) 4th order mass difference: the idea is to correct a
possible constant offset of \texttt{D2nSn}: can be used, where data with
less and more neutrons is available.

Binding-energy form

\[
10^{-3}[0.5B(Z,N-4)-2B(Z,N-2)+3B(Z,N)-2B(Z,N+2)+0.5B(Z,N+4)]
\]

Equivalent mass-excess form

\[
\begin{aligned}
& -0.0005M(Z,N-4)+0.002M(Z,N-2)-0.003M(Z,N) \\
& +0.002M(Z,N+2)-0.0005M(Z,N+4)
\end{aligned}
\]

Output unit: MeV

Exact stencil \texttt{terms}

\begin{longtable}[]{@{}rrr@{}}
\toprule\noalign{}
ΔZ & ΔN & Coefficient \\
\midrule\noalign{}
\endhead
\bottomrule\noalign{}
\endlastfoot
0 & -4 & -0.0005 \\
0 & -2 & 0.002 \\
0 & 0 & -0.003 \\
0 & 2 & 0.002 \\
0 & 4 & -0.0005 \\
\end{longtable}

Exact stencil \texttt{be\_terms}

\begin{longtable}[]{@{}rrr@{}}
\toprule\noalign{}
ΔZ & ΔN & Coefficient \\
\midrule\noalign{}
\endhead
\bottomrule\noalign{}
\endlastfoot
0 & -4 & 0.5 \\
0 & -2 & -2.0 \\
0 & 0 & 3.0 \\
0 & 2 & -2.0 \\
0 & 4 & 0.5 \\
\end{longtable}

Exact stencil \texttt{light\_terms}

\begin{longtable}[]{@{}lr@{}}
\toprule\noalign{}
Particle & Coefficient \\
\midrule\noalign{}
\endhead
\bottomrule\noalign{}
\endlastfoot
H & 0.0 \\
n & 0.0 \\
\end{longtable}

offset\_latex:

\[
0
\]

\begin{itemize}
\item
  \texttt{N\_filter}: oe
\item
  \texttt{Z\_filter}: oe
\item
  \texttt{meaning}: e=even, o=odd, oe=both; applied by comparison()
  after calculation, not apply\_energy\_operator().
\item
  nuclear.py ·
  \texttt{operator\_catalog\ /\ apply\_energy\_operator\ /\ comparison}
  · line 147
\item
  nuclear\_operator\_weights.csv · `` · line 51
\item
  nuclear\_operator\_metadata.csv · `` · line 19
\end{itemize}

\textbf{Z=50, N=66} · Calculated from retained source data

0.02917599999994991 ± 0.001146961638416909 MeV

Recalculated and evaluated values/errors are distinct; source rounding
is retained. Full primitive inputs, flags and hashes are in this entry's
worked\_example in inventory-delivery.json.

Uncertainty: combine repeated primitive identities and use the stated
covariance approximation. Marginal mass errors alone do not recover full
AME evaluated covariance.

\subsection{18 gsi:V\_1p1n V\_1p1n: Proton-neutron interaction strength
where N/Z can be both odd and
even}\label{gsiv_1p1n-v_1p1n-proton-neutron-interaction-strength-where-nz-can-be-both-odd-and-even}

V\_1p1n: Proton-neutron interaction strength where N/Z can be both odd
and even

Use: Calculated by V\_1p1n(N,Z) = {[}S1n(N,Z)-S1n(N,Z-1){]} =
\{{[}BE(N,Z)-BE(N-1,Z){]} - {[}BE(N,Z-1)-BE(N-1,Z-1){]}\}.

Binding-energy form

\[
-1B(Z,N-1)+1B(Z,N)+1B(Z-1,N-1)-1B(Z-1,N)
\]

Equivalent mass-excess form

\[
1M(Z,N-1)-1M(Z,N)-1M(Z-1,N-1)+1M(Z-1,N)
\]

Output unit: keV

Exact stencil \texttt{terms}

\begin{longtable}[]{@{}rrr@{}}
\toprule\noalign{}
ΔZ & ΔN & Coefficient \\
\midrule\noalign{}
\endhead
\bottomrule\noalign{}
\endlastfoot
0 & -1 & 1.0 \\
0 & 0 & -1.0 \\
-1 & -1 & -1.0 \\
-1 & 0 & 1.0 \\
\end{longtable}

Exact stencil \texttt{be\_terms}

\begin{longtable}[]{@{}rrr@{}}
\toprule\noalign{}
ΔZ & ΔN & Coefficient \\
\midrule\noalign{}
\endhead
\bottomrule\noalign{}
\endlastfoot
0 & -1 & -1.0 \\
0 & 0 & 1.0 \\
-1 & -1 & 1.0 \\
-1 & 0 & -1.0 \\
\end{longtable}

Exact stencil \texttt{light\_terms}

\begin{longtable}[]{@{}lr@{}}
\toprule\noalign{}
Particle & Coefficient \\
\midrule\noalign{}
\endhead
\bottomrule\noalign{}
\endlastfoot
H & 0.0 \\
n & 0.0 \\
\end{longtable}

offset\_latex:

\[
0
\]

\begin{itemize}
\item
  \texttt{N\_filter}: oe
\item
  \texttt{Z\_filter}: oe
\item
  \texttt{meaning}: e=even, o=odd, oe=both; applied by comparison()
  after calculation, not apply\_energy\_operator().
\item
  nuclear.py ·
  \texttt{operator\_catalog\ /\ apply\_energy\_operator\ /\ comparison}
  · line 147
\item
  nuclear\_operator\_weights.csv · `` · line 56
\item
  nuclear\_operator\_metadata.csv · `` · line 22
\end{itemize}

\textbf{Z=50, N=66} · Calculated from retained source data

525.5829999998678 ± 0.31652172121356853 keV

Recalculated and evaluated values/errors are distinct; source rounding
is retained. Full primitive inputs, flags and hashes are in this entry's
worked\_example in inventory-delivery.json.

Uncertainty: combine repeated primitive identities and use the stated
covariance approximation. Marginal mass errors alone do not recover full
AME evaluated covariance.

\subsection{19 gsi:R1D1n1p dnp1: Residual neutron-proton interaction
1}\label{gsir1d1n1p-dnp1-residual-neutron-proton-interaction-1}

dnp1: Residual neutron-proton interaction 1

Use: Calculated by R1D1n1p(N,Z) =
1/2\emph{{[}V1n-1p(Z,N)-V1n-1p(Z,N-1)){]}=BE(Z-1,N-1)-BE(Z,N-1)+1/2}{[}BE(Z,N)-BE(Z-1,N)-BE(Z-1,N-2)+BE(Z,N-2){]}.
The delta-np(1) formula can be imagined as a generalization of
delta-np(2) also to even-even nuclei.

Binding-energy form

\[
\begin{aligned}
& 0.5B(Z,N-2)-1B(Z,N-1)+0.5B(Z,N) \\
& -0.5B(Z-1,N-2)+1B(Z-1,N-1)-0.5B(Z-1,N)
\end{aligned}
\]

Equivalent mass-excess form

\[
\begin{aligned}
& -0.5M(Z,N-2)+1M(Z,N-1)-0.5M(Z,N) \\
& +0.5M(Z-1,N-2)-1M(Z-1,N-1)+0.5M(Z-1,N)
\end{aligned}
\]

Output unit: keV

Exact stencil \texttt{terms}

\begin{longtable}[]{@{}rrr@{}}
\toprule\noalign{}
ΔZ & ΔN & Coefficient \\
\midrule\noalign{}
\endhead
\bottomrule\noalign{}
\endlastfoot
0 & -2 & -0.5 \\
0 & -1 & 1.0 \\
0 & 0 & -0.5 \\
-1 & -2 & 0.5 \\
-1 & -1 & -1.0 \\
-1 & 0 & 0.5 \\
\end{longtable}

Exact stencil \texttt{be\_terms}

\begin{longtable}[]{@{}rrr@{}}
\toprule\noalign{}
ΔZ & ΔN & Coefficient \\
\midrule\noalign{}
\endhead
\bottomrule\noalign{}
\endlastfoot
0 & -2 & 0.5 \\
0 & -1 & -1.0 \\
0 & 0 & 0.5 \\
-1 & -2 & -0.5 \\
-1 & -1 & 1.0 \\
-1 & 0 & -0.5 \\
\end{longtable}

Exact stencil \texttt{light\_terms}

\begin{longtable}[]{@{}lr@{}}
\toprule\noalign{}
Particle & Coefficient \\
\midrule\noalign{}
\endhead
\bottomrule\noalign{}
\endlastfoot
H & 0.0 \\
n & 0.0 \\
\end{longtable}

offset\_latex:

\[
0
\]

\begin{itemize}
\item
  \texttt{N\_filter}: oe
\item
  \texttt{Z\_filter}: oe
\item
  \texttt{meaning}: e=even, o=odd, oe=both; applied by comparison()
  after calculation, not apply\_energy\_operator().
\item
  nuclear.py ·
  \texttt{operator\_catalog\ /\ apply\_energy\_operator\ /\ comparison}
  · line 147
\item
  nuclear\_operator\_weights.csv · `` · line 60
\item
  nuclear\_operator\_metadata.csv · `` · line 23
\end{itemize}

\textbf{Z=50, N=66} · Calculated from retained source data

127.0779999999213 ± 0.31970650603326795 keV

Recalculated and evaluated values/errors are distinct; source rounding
is retained. Full primitive inputs, flags and hashes are in this entry's
worked\_example in inventory-delivery.json.

Uncertainty: combine repeated primitive identities and use the stated
covariance approximation. Marginal mass errors alone do not recover full
AME evaluated covariance.

\subsection{20 gsi:R2D1n1p dnp2: Residual neutron-proton interaction
2}\label{gsir2d1n1p-dnp2-residual-neutron-proton-interaction-2}

dnp2: Residual neutron-proton interaction 2

Use: Calculated by R2D1n1p(N,Z) =
D1pS1n(N,Z)-D2pS2n(N-1,Z-1)={[}S1n(N,Z)-S1n(N,Z-1){]}-1/4\emph{{[}S2n(N-1,Z-1)-S2n(N-1,Z-3){]}=\{{[}BE(N,Z)-BE(N-1,Z){]}
- {[}BE(N,Z-1)-BE(N-1,Z-1){]}\}-1/4}\{{[}BE(N-1,Z-1)-BE(N-3,Z-1){]} -
{[}BE(N-1,Z-3)-BE(N-3,Z-3){]}\}. the residual interaction between the
last two (unpaired) neutron and proton, which is only relevant for
odd-odd nuclei.

Binding-energy form

\[
\begin{aligned}
& -1B(Z,N-1)+1B(Z,N)+0.25B(Z-1,N-3) \\
& +0.75B(Z-1,N-1)-1B(Z-1,N)-0.25B(Z-3,N-3) \\
& +0.25B(Z-3,N-1)
\end{aligned}
\]

Equivalent mass-excess form

\[
\begin{aligned}
& 1M(Z,N-1)-1M(Z,N)-0.25M(Z-1,N-3) \\
& -0.75M(Z-1,N-1)+1M(Z-1,N)+0.25M(Z-3,N-3) \\
& -0.25M(Z-3,N-1)
\end{aligned}
\]

Output unit: keV

Exact stencil \texttt{terms}

\begin{longtable}[]{@{}rrr@{}}
\toprule\noalign{}
ΔZ & ΔN & Coefficient \\
\midrule\noalign{}
\endhead
\bottomrule\noalign{}
\endlastfoot
0 & -1 & 1.0 \\
0 & 0 & -1.0 \\
-1 & -3 & -0.25 \\
-1 & -1 & -0.75 \\
-1 & 0 & 1.0 \\
-3 & -3 & 0.25 \\
-3 & -1 & -0.25 \\
\end{longtable}

Exact stencil \texttt{be\_terms}

\begin{longtable}[]{@{}rrr@{}}
\toprule\noalign{}
ΔZ & ΔN & Coefficient \\
\midrule\noalign{}
\endhead
\bottomrule\noalign{}
\endlastfoot
0 & -1 & -1.0 \\
0 & 0 & 1.0 \\
-1 & -3 & 0.25 \\
-1 & -1 & 0.75 \\
-1 & 0 & -1.0 \\
-3 & -3 & -0.25 \\
-3 & -1 & 0.25 \\
\end{longtable}

Exact stencil \texttt{light\_terms}

\begin{longtable}[]{@{}lr@{}}
\toprule\noalign{}
Particle & Coefficient \\
\midrule\noalign{}
\endhead
\bottomrule\noalign{}
\endlastfoot
H & 0.0 \\
n & 0.0 \\
\end{longtable}

offset\_latex:

\[
0
\]

\begin{itemize}
\item
  \texttt{N\_filter}: o
\item
  \texttt{Z\_filter}: o
\item
  \texttt{meaning}: e=even, o=odd, oe=both; applied by comparison()
  after calculation, not apply\_energy\_operator().
\item
  nuclear.py ·
  \texttt{operator\_catalog\ /\ apply\_energy\_operator\ /\ comparison}
  · line 147
\item
  nuclear\_operator\_weights.csv · `` · line 66
\item
  nuclear\_operator\_metadata.csv · `` · line 25
\end{itemize}

\textbf{Z=51, N=67} · Calculated from retained source data

103.33525000009104 ± 8.973651308831872 keV

Recalculated and evaluated values/errors are distinct; source rounding
is retained. Full primitive inputs, flags and hashes are in this entry's
worked\_example in inventory-delivery.json.

Uncertainty: combine repeated primitive identities and use the stated
covariance approximation. Marginal mass errors alone do not recover full
AME evaluated covariance.

\subsection{21 gsi:Pgap3N\_even Three-point formula of the
neutron-pairing gap for even
N}\label{gsipgap3n_even-three-point-formula-of-the-neutron-pairing-gap-for-even-n}

Three-point formula of the neutron-pairing gap for even N

Use: Calculated by Pgap3N\_even(N,Z) =
1/2\emph{{[}-BE(Z,N-1)-BE(Z,N+1)+2*BE(Z,N){]} =
(-1)\^{}N}1/2*{[}S1n(N,Z)-S1n(N+1,Z){]}. Three-point formula of the
neutron-pairing gap is a measure of the odd-even neutron staggering
between the neighboring isotones with a fixed Z.

Binding-energy form

\[
-0.5B(Z,N-1)+1B(Z,N)-0.5B(Z,N+1)
\]

Equivalent mass-excess form

\[
0.5M(Z,N-1)-1M(Z,N)+0.5M(Z,N+1)
\]

Output unit: keV

Exact stencil \texttt{terms}

\begin{longtable}[]{@{}rrr@{}}
\toprule\noalign{}
ΔZ & ΔN & Coefficient \\
\midrule\noalign{}
\endhead
\bottomrule\noalign{}
\endlastfoot
0 & -1 & 0.5 \\
0 & 0 & -1.0 \\
0 & 1 & 0.5 \\
\end{longtable}

Exact stencil \texttt{be\_terms}

\begin{longtable}[]{@{}rrr@{}}
\toprule\noalign{}
ΔZ & ΔN & Coefficient \\
\midrule\noalign{}
\endhead
\bottomrule\noalign{}
\endlastfoot
0 & -1 & -0.5 \\
0 & 0 & 1.0 \\
0 & 1 & -0.5 \\
\end{longtable}

Exact stencil \texttt{light\_terms}

\begin{longtable}[]{@{}lr@{}}
\toprule\noalign{}
Particle & Coefficient \\
\midrule\noalign{}
\endhead
\bottomrule\noalign{}
\endlastfoot
H & 0.0 \\
n & 0.0 \\
\end{longtable}

offset\_latex:

\[
0
\]

\begin{itemize}
\item
  \texttt{N\_filter}: e
\item
  \texttt{Z\_filter}: e
\item
  \texttt{meaning}: e=even, o=odd, oe=both; applied by comparison()
  after calculation, not apply\_energy\_operator().
\item
  nuclear.py ·
  \texttt{operator\_catalog\ /\ apply\_energy\_operator\ /\ comparison}
  · line 147
\item
  nuclear\_operator\_weights.csv · `` · line 73
\item
  nuclear\_operator\_metadata.csv · `` · line 27
\end{itemize}

\textbf{Z=50, N=66} · Calculated from retained source data

1310.1849999999395 ± 0.25998942286177723 keV

Recalculated and evaluated values/errors are distinct; source rounding
is retained. Full primitive inputs, flags and hashes are in this entry's
worked\_example in inventory-delivery.json.

Uncertainty: combine repeated primitive identities and use the stated
covariance approximation. Marginal mass errors alone do not recover full
AME evaluated covariance.

\subsection{22 gsi:Pgap3N\_odd Three-point formula of the
neutron-pairing gap for odd
N}\label{gsipgap3n_odd-three-point-formula-of-the-neutron-pairing-gap-for-odd-n}

Three-point formula of the neutron-pairing gap for odd N

Use: Calculated by Pgap3N\_odd(N,Z) =
-1/2\emph{{[}-BE(Z,N-1)-BE(Z,N+1)+2*BE(Z,N){]} =
(-1)\^{}N}1/2*{[}S1n(N,Z)-S1n(N+1,Z){]}. Three-point formula of the
neutron-pairing gap is a measure of the odd-even neutron staggering
between the neighboring isotones with a fixed Z.

Binding-energy form

\[
0.5B(Z,N-1)-1B(Z,N)+0.5B(Z,N+1)
\]

Equivalent mass-excess form

\[
-0.5M(Z,N-1)+1M(Z,N)-0.5M(Z,N+1)
\]

Output unit: keV

Exact stencil \texttt{terms}

\begin{longtable}[]{@{}rrr@{}}
\toprule\noalign{}
ΔZ & ΔN & Coefficient \\
\midrule\noalign{}
\endhead
\bottomrule\noalign{}
\endlastfoot
0 & -1 & -0.5 \\
0 & 0 & 1.0 \\
0 & 1 & -0.5 \\
\end{longtable}

Exact stencil \texttt{be\_terms}

\begin{longtable}[]{@{}rrr@{}}
\toprule\noalign{}
ΔZ & ΔN & Coefficient \\
\midrule\noalign{}
\endhead
\bottomrule\noalign{}
\endlastfoot
0 & -1 & 0.5 \\
0 & 0 & -1.0 \\
0 & 1 & 0.5 \\
\end{longtable}

Exact stencil \texttt{light\_terms}

\begin{longtable}[]{@{}lr@{}}
\toprule\noalign{}
Particle & Coefficient \\
\midrule\noalign{}
\endhead
\bottomrule\noalign{}
\endlastfoot
H & 0.0 \\
n & 0.0 \\
\end{longtable}

offset\_latex:

\[
0
\]

\begin{itemize}
\item
  \texttt{N\_filter}: o
\item
  \texttt{Z\_filter}: o
\item
  \texttt{meaning}: e=even, o=odd, oe=both; applied by comparison()
  after calculation, not apply\_energy\_operator().
\item
  nuclear.py ·
  \texttt{operator\_catalog\ /\ apply\_energy\_operator\ /\ comparison}
  · line 147
\item
  nuclear\_operator\_weights.csv · `` · line 76
\item
  nuclear\_operator\_metadata.csv · `` · line 29
\end{itemize}

\textbf{Z=51, N=67} · Calculated from retained source data

1061.3474999999744 ± 6.255797251350143 keV

Recalculated and evaluated values/errors are distinct; source rounding
is retained. Full primitive inputs, flags and hashes are in this entry's
worked\_example in inventory-delivery.json.

Uncertainty: combine repeated primitive identities and use the stated
covariance approximation. Marginal mass errors alone do not recover full
AME evaluated covariance.

\subsection{23 gsi:Pgap3Z\_even Three-point formula of the
proton-pairing gap for even
Z}\label{gsipgap3z_even-three-point-formula-of-the-proton-pairing-gap-for-even-z}

Three-point formula of the proton-pairing gap for even Z

Use: Calculated by Pgap3Z\_even(N,Z) =
1/2\emph{{[}-BE(Z-1,N)-BE(Z+1,N)+2*BE(Z,N){]} =
(-1)\^{}N}1/2*{[}S1p(N,Z)-S1p(N,Z+1){]}. Three-point formula of the
proton-pairing gap is a measure of the odd-even proton staggering
between the neighboring isotones with a fixed N.

Binding-energy form

\[
-0.5B(Z+1,N)+1B(Z,N)-0.5B(Z-1,N)
\]

Equivalent mass-excess form

\[
0.5M(Z+1,N)-1M(Z,N)+0.5M(Z-1,N)
\]

Output unit: keV

Exact stencil \texttt{terms}

\begin{longtable}[]{@{}rrr@{}}
\toprule\noalign{}
ΔZ & ΔN & Coefficient \\
\midrule\noalign{}
\endhead
\bottomrule\noalign{}
\endlastfoot
1 & 0 & 0.5 \\
0 & 0 & -1.0 \\
-1 & 0 & 0.5 \\
\end{longtable}

Exact stencil \texttt{be\_terms}

\begin{longtable}[]{@{}rrr@{}}
\toprule\noalign{}
ΔZ & ΔN & Coefficient \\
\midrule\noalign{}
\endhead
\bottomrule\noalign{}
\endlastfoot
1 & 0 & -0.5 \\
0 & 0 & 1.0 \\
-1 & 0 & -0.5 \\
\end{longtable}

Exact stencil \texttt{light\_terms}

\begin{longtable}[]{@{}lr@{}}
\toprule\noalign{}
Particle & Coefficient \\
\midrule\noalign{}
\endhead
\bottomrule\noalign{}
\endlastfoot
H & 0.0 \\
n & 0.0 \\
\end{longtable}

offset\_latex:

\[
0
\]

\begin{itemize}
\item
  \texttt{N\_filter}: e
\item
  \texttt{Z\_filter}: e
\item
  \texttt{meaning}: e=even, o=odd, oe=both; applied by comparison()
  after calculation, not apply\_energy\_operator().
\item
  nuclear.py ·
  \texttt{operator\_catalog\ /\ apply\_energy\_operator\ /\ comparison}
  · line 147
\item
  nuclear\_operator\_weights.csv · `` · line 79
\item
  nuclear\_operator\_metadata.csv · `` · line 31
\end{itemize}

\textbf{Z=50, N=66} · Calculated from retained source data

2438.0185000000056 ± 4.219596455823708 keV

Recalculated and evaluated values/errors are distinct; source rounding
is retained. Full primitive inputs, flags and hashes are in this entry's
worked\_example in inventory-delivery.json.

Uncertainty: combine repeated primitive identities and use the stated
covariance approximation. Marginal mass errors alone do not recover full
AME evaluated covariance.

\subsection{24 gsi:Pgap3Z\_odd Three-point formula of the proton-pairing
gap for odd
Z}\label{gsipgap3z_odd-three-point-formula-of-the-proton-pairing-gap-for-odd-z}

Three-point formula of the proton-pairing gap for odd Z

Use: Calculated by Pgap3Z\_odd(N,Z) =
-1/2\emph{{[}-BE(Z-1,N)-BE(Z+1,N)+2*BE(Z,N){]} =
(-1)\^{}N}1/2*{[}S1p(N,Z)-S1p(N,Z+1){]}. Three-point formula of the
proton-pairing gap is a measure of the odd-even proton staggering
between the neighboring isotones with a fixed N.

Binding-energy form

\[
0.5B(Z+1,N)-1B(Z,N)+0.5B(Z-1,N)
\]

Equivalent mass-excess form

\[
-0.5M(Z+1,N)+1M(Z,N)-0.5M(Z-1,N)
\]

Output unit: keV

Exact stencil \texttt{terms}

\begin{longtable}[]{@{}rrr@{}}
\toprule\noalign{}
ΔZ & ΔN & Coefficient \\
\midrule\noalign{}
\endhead
\bottomrule\noalign{}
\endlastfoot
1 & 0 & -0.5 \\
0 & 0 & 1.0 \\
-1 & 0 & -0.5 \\
\end{longtable}

Exact stencil \texttt{be\_terms}

\begin{longtable}[]{@{}rrr@{}}
\toprule\noalign{}
ΔZ & ΔN & Coefficient \\
\midrule\noalign{}
\endhead
\bottomrule\noalign{}
\endlastfoot
1 & 0 & 0.5 \\
0 & 0 & -1.0 \\
-1 & 0 & 0.5 \\
\end{longtable}

Exact stencil \texttt{light\_terms}

\begin{longtable}[]{@{}lr@{}}
\toprule\noalign{}
Particle & Coefficient \\
\midrule\noalign{}
\endhead
\bottomrule\noalign{}
\endlastfoot
H & 0.0 \\
n & 0.0 \\
\end{longtable}

offset\_latex:

\[
0
\]

\begin{itemize}
\item
  \texttt{N\_filter}: o
\item
  \texttt{Z\_filter}: o
\item
  \texttt{meaning}: e=even, o=odd, oe=both; applied by comparison()
  after calculation, not apply\_energy\_operator().
\item
  nuclear.py ·
  \texttt{operator\_catalog\ /\ apply\_energy\_operator\ /\ comparison}
  · line 147
\item
  nuclear\_operator\_weights.csv · `` · line 82
\item
  nuclear\_operator\_metadata.csv · `` · line 33
\end{itemize}

\textbf{Z=51, N=67} · Calculated from retained source data

793.9364999999525 ± 4.732536238635685 keV

Recalculated and evaluated values/errors are distinct; source rounding
is retained. Full primitive inputs, flags and hashes are in this entry's
worked\_example in inventory-delivery.json.

Uncertainty: combine repeated primitive identities and use the stated
covariance approximation. Marginal mass errors alone do not recover full
AME evaluated covariance.

\subsection{25 gsi:W1 Wigner Energy 1}\label{gsiw1-wigner-energy-1}

Wigner Energy 1

Use: Calculated by Wigner1(N,Z)
=D2pS2n(Z,N)-1/2\emph{{[}D2pS2n(Z-2,N)+D2pS2n(Z,N+2){]}=1/8}{[}3BE(Z,N)-4BE(Z-2,N)+3BE(Z-2,N-2)-2BE(Z,N-2)+BE(Z-4,N)-BE(Z-4,N-2)-BE(Z,N+2)+BE(Z-2,N+2){]}
/n

Binding-energy form

\[
\begin{aligned}
& -0.25B(Z,N-2)+0.375B(Z,N)-0.125B(Z,N+2) \\
& +0.375B(Z-2,N-2)-0.5B(Z-2,N)+0.125B(Z-2,N+2) \\
& -0.125B(Z-4,N-2)+0.125B(Z-4,N)
\end{aligned}
\]

Equivalent mass-excess form

\[
\begin{aligned}
& 0.25M(Z,N-2)-0.375M(Z,N)+0.125M(Z,N+2) \\
& -0.375M(Z-2,N-2)+0.5M(Z-2,N)-0.125M(Z-2,N+2) \\
& +0.125M(Z-4,N-2)-0.125M(Z-4,N)
\end{aligned}
\]

Output unit: keV

Exact stencil \texttt{terms}

\begin{longtable}[]{@{}rrr@{}}
\toprule\noalign{}
ΔZ & ΔN & Coefficient \\
\midrule\noalign{}
\endhead
\bottomrule\noalign{}
\endlastfoot
0 & -2 & 0.25 \\
0 & 0 & -0.375 \\
0 & 2 & 0.125 \\
-2 & -2 & -0.375 \\
-2 & 0 & 0.5 \\
-2 & 2 & -0.125 \\
-4 & -2 & 0.125 \\
-4 & 0 & -0.125 \\
\end{longtable}

Exact stencil \texttt{be\_terms}

\begin{longtable}[]{@{}rrr@{}}
\toprule\noalign{}
ΔZ & ΔN & Coefficient \\
\midrule\noalign{}
\endhead
\bottomrule\noalign{}
\endlastfoot
0 & -2 & -0.25 \\
0 & 0 & 0.375 \\
0 & 2 & -0.125 \\
-2 & -2 & 0.375 \\
-2 & 0 & -0.5 \\
-2 & 2 & 0.125 \\
-4 & -2 & -0.125 \\
-4 & 0 & 0.125 \\
\end{longtable}

Exact stencil \texttt{light\_terms}

\begin{longtable}[]{@{}lr@{}}
\toprule\noalign{}
Particle & Coefficient \\
\midrule\noalign{}
\endhead
\bottomrule\noalign{}
\endlastfoot
H & 0.0 \\
n & 0.0 \\
\end{longtable}

offset\_latex:

\[
0
\]

\begin{itemize}
\item
  \texttt{N\_filter}: oe
\item
  \texttt{Z\_filter}: oe
\item
  \texttt{meaning}: e=even, o=odd, oe=both; applied by comparison()
  after calculation, not apply\_energy\_operator().
\item
  nuclear.py ·
  \texttt{operator\_catalog\ /\ apply\_energy\_operator\ /\ comparison}
  · line 147
\item
  nuclear\_operator\_weights.csv · `` · line 85
\item
  nuclear\_operator\_metadata.csv · `` · line 35
\end{itemize}

\textbf{Z=50, N=66} · Calculated from retained source data

21.63537499999802 ± 0.841936801740487 keV

Recalculated and evaluated values/errors are distinct; source rounding
is retained. Full primitive inputs, flags and hashes are in this entry's
worked\_example in inventory-delivery.json.

Uncertainty: combine repeated primitive identities and use the stated
covariance approximation. Marginal mass errors alone do not recover full
AME evaluated covariance.

\subsection{26 gsi:W2 Wigner Energy 2}\label{gsiw2-wigner-energy-2}

Wigner Energy 2

Use: Calculated by W2(N,Z)
=1/2\emph{{[}D2pS2n(Z-1,N-1)+D2pS2n(Z+1,N+1){]}-D2pS2n(Z-1,N+1)=1/8}{[}4BE(Z-1,N-1)-BE(Z-1,N-3)+BE(Z-3,N-3)+BE(Z+1,N+1)-BE(Z+1,N-1)-3BE(Z-3,N-1)-3BE(Z-1,N+1)+2BE(Z-3,N+1){]}
/n

Binding-energy form

\[
\begin{aligned}
& -0.125B(Z+1,N-1)+0.125B(Z+1,N+1)-0.125B(Z-1,N-3) \\
& +0.5B(Z-1,N-1)-0.375B(Z-1,N+1)+0.125B(Z-3,N-3) \\
& -0.375B(Z-3,N-1)+0.25B(Z-3,N+1)
\end{aligned}
\]

Equivalent mass-excess form

\[
\begin{aligned}
& 0.125M(Z+1,N-1)-0.125M(Z+1,N+1)+0.125M(Z-1,N-3) \\
& -0.5M(Z-1,N-1)+0.375M(Z-1,N+1)-0.125M(Z-3,N-3) \\
& +0.375M(Z-3,N-1)-0.25M(Z-3,N+1)
\end{aligned}
\]

Output unit: keV

Exact stencil \texttt{terms}

\begin{longtable}[]{@{}rrr@{}}
\toprule\noalign{}
ΔZ & ΔN & Coefficient \\
\midrule\noalign{}
\endhead
\bottomrule\noalign{}
\endlastfoot
1 & -1 & 0.125 \\
1 & 1 & -0.125 \\
-1 & -3 & 0.125 \\
-1 & -1 & -0.5 \\
-1 & 1 & 0.375 \\
-3 & -3 & -0.125 \\
-3 & -1 & 0.375 \\
-3 & 1 & -0.25 \\
\end{longtable}

Exact stencil \texttt{be\_terms}

\begin{longtable}[]{@{}rrr@{}}
\toprule\noalign{}
ΔZ & ΔN & Coefficient \\
\midrule\noalign{}
\endhead
\bottomrule\noalign{}
\endlastfoot
1 & -1 & -0.125 \\
1 & 1 & 0.125 \\
-1 & -3 & -0.125 \\
-1 & -1 & 0.5 \\
-1 & 1 & -0.375 \\
-3 & -3 & 0.125 \\
-3 & -1 & -0.375 \\
-3 & 1 & 0.25 \\
\end{longtable}

Exact stencil \texttt{light\_terms}

\begin{longtable}[]{@{}lr@{}}
\toprule\noalign{}
Particle & Coefficient \\
\midrule\noalign{}
\endhead
\bottomrule\noalign{}
\endlastfoot
H & 0.0 \\
n & 0.0 \\
\end{longtable}

offset\_latex:

\[
0
\]

\begin{itemize}
\item
  \texttt{N\_filter}: oe
\item
  \texttt{Z\_filter}: oe
\item
  \texttt{meaning}: e=even, o=odd, oe=both; applied by comparison()
  after calculation, not apply\_energy\_operator().
\item
  nuclear.py ·
  \texttt{operator\_catalog\ /\ apply\_energy\_operator\ /\ comparison}
  · line 147
\item
  nuclear\_operator\_weights.csv · `` · line 93
\item
  nuclear\_operator\_metadata.csv · `` · line 36
\end{itemize}

\textbf{Z=50, N=66} · Calculated from retained source data

35.21899999995367 ± 1.7382872542031136 keV

Recalculated and evaluated values/errors are distinct; source rounding
is retained. Full primitive inputs, flags and hashes are in this entry's
worked\_example in inventory-delivery.json.

Uncertainty: combine repeated primitive identities and use the stated
covariance approximation. Marginal mass errors alone do not recover full
AME evaluated covariance.

\subsection{27 gsi:W3 Wigner Energy 3}\label{gsiw3-wigner-energy-3}

Wigner Energy 3

Use: Calculated by W3(N,Z)
={[}D2pS2n(Z-1,N)+D2pS2n(Z,N+1){]}-2D2pS2n(Z-1,N+1)=1/4*{[}BE(Z-1,N)-BE(Z-1,N-2)-BE(Z-3,N)+BE(Z-3,N-2)+BE(Z,N+1)-BE(Z,N-1)-BE(Z-2,N+1)+BE(Z-2,N-1)-2BE(Z-1,N+1)+2BE(Z-1,N-1)+2BE(Z-3,N+1)-2BE(Z-3,N-1){]}
/n

Binding-energy form

\[
\begin{aligned}
& -0.25B(Z,N-1)+0.25B(Z,N+1)-0.25B(Z-1,N-2) \\
& -0.5B(Z-1,N-1)+0.25B(Z-1,N)+0.5B(Z-1,N+1) \\
& +0.25B(Z-2,N-1)-0.25B(Z-2,N+1)+0.25B(Z-3,N-2) \\
& +0.5B(Z-3,N-1)-0.25B(Z-3,N)-0.5B(Z-3,N+1)
\end{aligned}
\]

Equivalent mass-excess form

\[
\begin{aligned}
& 0.25M(Z,N-1)-0.25M(Z,N+1)+0.25M(Z-1,N-2) \\
& +0.5M(Z-1,N-1)-0.25M(Z-1,N)-0.5M(Z-1,N+1) \\
& -0.25M(Z-2,N-1)+0.25M(Z-2,N+1)-0.25M(Z-3,N-2) \\
& -0.5M(Z-3,N-1)+0.25M(Z-3,N)+0.5M(Z-3,N+1)
\end{aligned}
\]

Output unit: keV

Exact stencil \texttt{terms}

\begin{longtable}[]{@{}rrr@{}}
\toprule\noalign{}
ΔZ & ΔN & Coefficient \\
\midrule\noalign{}
\endhead
\bottomrule\noalign{}
\endlastfoot
0 & -1 & 0.25 \\
0 & 1 & -0.25 \\
-1 & -2 & 0.25 \\
-1 & -1 & 0.5 \\
-1 & 0 & -0.25 \\
-1 & 1 & -0.5 \\
-2 & -1 & -0.25 \\
-2 & 1 & 0.25 \\
-3 & -2 & -0.25 \\
-3 & -1 & -0.5 \\
-3 & 0 & 0.25 \\
-3 & 1 & 0.5 \\
\end{longtable}

Exact stencil \texttt{be\_terms}

\begin{longtable}[]{@{}rrr@{}}
\toprule\noalign{}
ΔZ & ΔN & Coefficient \\
\midrule\noalign{}
\endhead
\bottomrule\noalign{}
\endlastfoot
0 & -1 & -0.25 \\
0 & 1 & 0.25 \\
-1 & -2 & -0.25 \\
-1 & -1 & -0.5 \\
-1 & 0 & 0.25 \\
-1 & 1 & 0.5 \\
-2 & -1 & 0.25 \\
-2 & 1 & -0.25 \\
-3 & -2 & 0.25 \\
-3 & -1 & 0.5 \\
-3 & 0 & -0.25 \\
-3 & 1 & -0.5 \\
\end{longtable}

Exact stencil \texttt{light\_terms}

\begin{longtable}[]{@{}lr@{}}
\toprule\noalign{}
Particle & Coefficient \\
\midrule\noalign{}
\endhead
\bottomrule\noalign{}
\endlastfoot
H & 0.0 \\
n & 0.0 \\
\end{longtable}

offset\_latex:

\[
0
\]

\begin{itemize}
\item
  \texttt{N\_filter}: oe
\item
  \texttt{Z\_filter}: oe
\item
  \texttt{meaning}: e=even, o=odd, oe=both; applied by comparison()
  after calculation, not apply\_energy\_operator().
\item
  nuclear.py ·
  \texttt{operator\_catalog\ /\ apply\_energy\_operator\ /\ comparison}
  · line 147
\item
  nuclear\_operator\_weights.csv · `` · line 101
\item
  nuclear\_operator\_metadata.csv · `` · line 37
\end{itemize}

\textbf{Z=50, N=66} · Calculated from retained source data

1336.5740000000224 ± 4.919423759445816 keV

Recalculated and evaluated values/errors are distinct; source rounding
is retained. Full primitive inputs, flags and hashes are in this entry's
worked\_example in inventory-delivery.json.

Uncertainty: combine repeated primitive identities and use the stated
covariance approximation. Marginal mass errors alone do not recover full
AME evaluated covariance.

\subsection{28 gsi:Qec Qec value of electron
capture}\label{gsiqec-qec-value-of-electron-capture}

Qec value of electron capture

Use: Calculated by Qec(N,Z) = ME(N,Z)-ME(N+1,Z-1) =
BE(N+1,Z-1)-BE(N,Z)-ME\_n+ME\_p /n

Binding-energy form

\[
-1B(Z,N)+1B(Z-1,N+1)-782.34704
\]

Equivalent mass-excess form

\[
1M(Z,N)-1M(Z-1,N+1) + -1 H + 1 n -782.34704
\]

Output unit: keV

Exact stencil \texttt{terms}

\begin{longtable}[]{@{}rrr@{}}
\toprule\noalign{}
ΔZ & ΔN & Coefficient \\
\midrule\noalign{}
\endhead
\bottomrule\noalign{}
\endlastfoot
0 & 0 & 1.0 \\
-1 & 1 & -1.0 \\
\end{longtable}

Exact stencil \texttt{be\_terms}

\begin{longtable}[]{@{}rrr@{}}
\toprule\noalign{}
ΔZ & ΔN & Coefficient \\
\midrule\noalign{}
\endhead
\bottomrule\noalign{}
\endlastfoot
0 & 0 & -1.0 \\
-1 & 1 & 1.0 \\
\end{longtable}

Exact stencil \texttt{light\_terms}

\begin{longtable}[]{@{}lr@{}}
\toprule\noalign{}
Particle & Coefficient \\
\midrule\noalign{}
\endhead
\bottomrule\noalign{}
\endlastfoot
H & -1.0 \\
n & 1.0 \\
\end{longtable}

offset\_latex:

\[
-782.34704
\]

\begin{itemize}
\item
  \texttt{N\_filter}: oe
\item
  \texttt{Z\_filter}: oe
\item
  \texttt{meaning}: e=even, o=odd, oe=both; applied by comparison()
  after calculation, not apply\_energy\_operator().
\item
  nuclear.py ·
  \texttt{operator\_catalog\ /\ apply\_energy\_operator\ /\ comparison}
  · line 147
\item
  nuclear\_operator\_weights.csv · `` · line 113
\item
  nuclear\_operator\_metadata.csv · `` · line 38
\end{itemize}

\textbf{Z=50, N=66} · Calculated from retained source data

-3276.220999999881 ± 0.24003333101883997 keV

Recalculated and evaluated values/errors are distinct; source rounding
is retained. Full primitive inputs, flags and hashes are in this entry's
worked\_example in inventory-delivery.json.

Uncertainty: combine repeated primitive identities and use the stated
covariance approximation. Marginal mass errors alone do not recover full
AME evaluated covariance.

\section{Data links and citations}\label{data-links-and-citations}

These records distinguish original tables, paper definitions, historical
evaluations and code conventions. Retrieval date is not evaluation date,
and a current online edition need not match a pinned local snapshot.
Complete URLs, hashes, retrieval status and limits are in sources.json
and inventory-delivery.json beside this guide.

\subsection{ame2020}\label{ame2020}

W. J. Huang et al., Chinese Physics C 45, 030002 (2021); M. Wang et al.,
Chinese Physics C 45, 030003 (2021).

Evaluation or version date: 2020

Original numerical retrieval date: not recorded; link-check date is not
numerical retrieval date

\begin{itemize}
\tightlist
\item
  \href{https://www-nds.iaea.org/amdc/}{primary metadata/reference}
\item
  \href{https://cpc.ihep.ac.cn/article/doi/10.1088/1674-1137/abddaf}{primary
  metadata/reference}
\item
  DOI
  \href{https://doi.org/10.1088/1674-1137/abddb0}{10.1088/1674-1137/abddb0}
\item
  DOI
  \href{https://doi.org/10.1088/1674-1137/abddaf}{10.1088/1674-1137/abddaf}
\item
  \href{https://www-nds.iaea.org/amdc/ame2020/mass_1.mas20.txt}{mass\_1.mas20.txt}
\item
  \href{https://www-nds.iaea.org/amdc/ame2020/rct1.mas20.txt}{rct1.mas20.txt}
\item
  \href{https://www-nds.iaea.org/amdc/ame2020/rct2_1.mas20.txt}{rct2\_1.mas20.txt}
\end{itemize}

Published evaluated quantity errors retain AME correlations. Rederived
arbitrary stencils do not recover the full AME covariance from marginal
mass errors. \# value and uncertainty flags remain separate; missing *,
bounds and symbolic values are not zero.

Current link verification is not a new numerical-data retrieval; source
files were not downloaded by this audit.

\subsection{ame2016}\label{ame2016}

W. J. Huang et al., Chinese Physics C 41, 030002 (2017); M. Wang et al.,
Chinese Physics C 41, 030003 (2017).

Evaluation or version date: 2016

Original numerical retrieval date: 2026-10-23T02:53:19.900902Z

\begin{itemize}
\tightlist
\item
  \href{https://cpc.ihep.ac.cn/article/doi/10.1088/1674-1137/41/3/030002}{primary
  metadata/reference}
\item
  \href{https://cpc.ihep.ac.cn/article/doi/10.1088/1674-1137/41/3/030003}{primary
  metadata/reference}
\item
  \href{https://www-nds.iaea.org/amdc/ame2016/mass16.txt}{actual
  recorded numerical source}
\item
  DOI
  \href{https://doi.org/10.1088/1674-1137/41/3/030002}{10.1088/1674-1137/41/3/030002}
\item
  DOI
  \href{https://doi.org/10.1088/1674-1137/41/3/030003}{10.1088/1674-1137/41/3/030003}
\end{itemize}

Source \# estimates and missing tokens retained. Printed zero marginal
uncertainties may be rounded; no full covariance reconstructed. Editions
are separate evaluations, not independent repeated measurements.

Current link verification is not a new numerical-data retrieval; source
files were not downloaded by this audit.

\subsection{ame2012}\label{ame2012}

G. Audi et al., Chinese Physics C 36, 1287--1602 (2012); M. Wang et al.,
Chinese Physics C 36, 1603--2014 (2012).

Evaluation or version date: 2012

Original numerical retrieval date: not recorded; link-check date is not
numerical retrieval date

\begin{itemize}
\tightlist
\item
  \href{https://cpc.ihep.ac.cn/article/doi/10.1088/1674-1137/36/12/002}{primary
  metadata/reference}
\item
  \href{https://cpc.ihep.ac.cn/article/doi/10.1088/1674-1137/36/12/003}{primary
  metadata/reference}
\item
  \href{https://amdc.impcas.ac.cn/masstables/Ame2012/mass.mas12}{actual
  recorded numerical source}
\item
  DOI
  \href{https://doi.org/10.1088/1674-1137/36/12/002}{10.1088/1674-1137/36/12/002}
\item
  DOI
  \href{https://doi.org/10.1088/1674-1137/36/12/003}{10.1088/1674-1137/36/12/003}
\end{itemize}

Source \# estimates and missing tokens retained. Printed zero marginal
uncertainties may be rounded; no full covariance reconstructed. Editions
are separate evaluations, not independent repeated measurements.

Current link verification is not a new numerical-data retrieval; source
files were not downloaded by this audit.

\subsection{ame2003}\label{ame2003}

A. H. Wapstra, G. Audi and C. Thibault, Nuclear Physics A 729, 129--336
(2003); G. Audi, A. H. Wapstra and C. Thibault, Nuclear Physics A 729,
337--676 (2003).

Evaluation or version date: 2003

Original numerical retrieval date: 2026-10-23T02:53:22.411179Z

\begin{itemize}
\tightlist
\item
  \href{https://www-nds.iaea.org/amdc/}{primary metadata/reference}
\item
  \href{https://amdc.impcas.ac.cn/masstables/Ame2003/mass.mas03}{actual
  recorded numerical source}
\item
  DOI
  \href{https://doi.org/10.1016/j.nuclphysa.2003.11.002}{10.1016/j.nuclphysa.2003.11.002}
\item
  DOI
  \href{https://doi.org/10.1016/j.nuclphysa.2003.11.003}{10.1016/j.nuclphysa.2003.11.003}
\end{itemize}

Source \# estimates and missing tokens retained. Printed zero marginal
uncertainties may be rounded; no full covariance reconstructed. Editions
are separate evaluations, not independent repeated measurements.

Current link verification is not a new numerical-data retrieval; source
files were not downloaded by this audit.

\subsection{ame1995}\label{ame1995}

G. Audi and A. H. Wapstra, Nuclear Physics A 595, 409--480 (1995).

Evaluation or version date: 1995

Original numerical retrieval date: 2026-10-23T02:53:27.079438Z

\begin{itemize}
\tightlist
\item
  \href{https://www-nds.iaea.org/amdc/}{primary metadata/reference}
\item
  \href{https://amdc.impcas.ac.cn/masstables/Ame1995/mass_rmd.mas95}{actual
  recorded numerical source}
\item
  DOI
  \href{https://doi.org/10.1016/0375-9474(95)00445-9}{10.1016/0375-9474(95)00445-9}
\end{itemize}

Source \# estimates and missing tokens retained. Printed zero marginal
uncertainties may be rounded; no full covariance reconstructed. Editions
are separate evaluations, not independent repeated measurements.

Current link verification is not a new numerical-data retrieval; source
files were not downloaded by this audit.

\subsection{nubase2020}\label{nubase2020}

F. G. Kondev, M. Wang, W. J. Huang, S. Naimi and G. Audi, Chinese
Physics C 45, 030001 (2021).

Evaluation or version date: 2020

Original numerical retrieval date: not recorded; link-check date is not
numerical retrieval date

\begin{itemize}
\tightlist
\item
  \href{https://cpc.ihep.ac.cn/article/doi/10.1088/1674-1137/abddae}{primary
  metadata/reference}
\item
  \href{https://www-nds.iaea.org/amdc/}{primary metadata/reference}
\item
  DOI
  \href{https://doi.org/10.1088/1674-1137/abddae}{10.1088/1674-1137/abddae}
\end{itemize}

Raw flags, limits, symbolic excitation offsets and source state index
are retained. Not every state is an isomer and not every numeric field
is a measured Gaussian value.

Current link verification is not a new numerical-data retrieval; source
files were not downloaded by this audit.

\subsection{nubase2016}\label{nubase2016}

G. Audi, F. G. Kondev, M. Wang, W. J. Huang and S. Naimi, Chinese
Physics C 41, 030001 (2017).

Evaluation or version date: 2016

Original numerical retrieval date: not recorded; link-check date is not
numerical retrieval date

\begin{itemize}
\tightlist
\item
  \href{https://cpc.ihep.ac.cn/article/doi/10.1088/1674-1137/41/3/030001}{primary
  metadata/reference}
\item
  DOI
  \href{https://doi.org/10.1088/1674-1137/41/3/030001}{10.1088/1674-1137/41/3/030001}
\end{itemize}

Evaluation includes measured and systematics-estimated fields.
Ground-state and isomer counts use retained classifier definitions;
other excited/analogue records are separate. Historical statistics are
not complete state payloads.

Current link verification is not a new numerical-data retrieval; source
files were not downloaded by this audit.

\subsection{nubase2012}\label{nubase2012}

G. Audi et al., Chinese Physics C 36, 1157--1286 (2012).

Evaluation or version date: 2012

Original numerical retrieval date: not recorded; link-check date is not
numerical retrieval date

\begin{itemize}
\tightlist
\item
  \href{https://cpc.ihep.ac.cn/article/doi/10.1088/1674-1137/36/12/001}{primary
  metadata/reference}
\item
  DOI
  \href{https://doi.org/10.1088/1674-1137/36/12/001}{10.1088/1674-1137/36/12/001}
\end{itemize}

Evaluation includes measured and systematics-estimated fields.
Ground-state and isomer counts use retained classifier definitions;
other excited/analogue records are separate. Historical statistics are
not complete state payloads.

Current link verification is not a new numerical-data retrieval; source
files were not downloaded by this audit.

\subsection{nubase2003}\label{nubase2003}

G. Audi, O. Bersillon, J. Blachot and A. H. Wapstra, Nuclear Physics A
729, 3--128 (2003).

Evaluation or version date: 2003

Original numerical retrieval date: not recorded; link-check date is not
numerical retrieval date

\begin{itemize}
\tightlist
\item
  \href{https://amdc.impcas.ac.cn/nubase/nubase2003/nubtab03.asc}{primary
  metadata/reference}
\item
  DOI
  \href{https://doi.org/10.1016/j.nuclphysa.2003.11.001}{10.1016/j.nuclphysa.2003.11.001}
\end{itemize}

Evaluation includes measured and systematics-estimated fields.
Ground-state and isomer counts use retained classifier definitions;
other excited/analogue records are separate. Historical statistics are
not complete state payloads.

Current link verification is not a new numerical-data retrieval; source
files were not downloaded by this audit.

\subsection{nubase1997}\label{nubase1997}

G. Audi, O. Bersillon, J. Blachot and A. H. Wapstra, Nuclear Physics A
624, 1--124 (1997).

Evaluation or version date: 1997

Original numerical retrieval date: not recorded; link-check date is not
numerical retrieval date

\begin{itemize}
\tightlist
\item
  \href{https://www.sciencedirect.com/science/article/pii/S037594749700482X}{primary
  metadata/reference}
\item
  DOI
  \href{https://doi.org/10.1016/S0375-9474(97)00482-X}{10.1016/S0375-9474(97)00482-X}
\end{itemize}

Evaluation includes measured and systematics-estimated fields.
Ground-state and isomer counts use retained classifier definitions;
other excited/analogue records are separate. Historical statistics are
not complete state payloads.

Current link verification is not a new numerical-data retrieval; source
files were not downloaded by this audit.

\subsection{angeli-marinova-2013}\label{angeli-marinova-2013}

I. Angeli and K. P. Marinova, Atomic Data and Nuclear Data Tables 99,
69--95 (2013).

Evaluation or version date: 2013

Original numerical retrieval date: 2026-10-23

\begin{itemize}
\tightlist
\item
  \href{https://www-nds.iaea.org/radii/}{primary metadata/reference}
\item
  DOI
  \href{https://doi.org/10.1016/j.adt.2011.12.006}{10.1016/j.adt.2011.12.006}
\end{itemize}

Published 2013 Ca rms radii only; preliminary update columns excluded.
Independent radius-error propagation does not supply
shared-reference/evaluation correlations.

Current link verification is not a new numerical-data retrieval; source
files were not downloaded by this audit.

\subsection{livechart}\label{livechart}

International Atomic Energy Agency Nuclear Data Section, LiveChart Data
Download API, v1 (v0 alias), official API guide.

Evaluation or version date: Field-dependent; guide identifies ENSDF
April 2022 snapshot plus AME2020/NUBASE2020 and named radii/moment
evaluations

Original numerical retrieval date: 2026-10-23

\begin{itemize}
\tightlist
\item
  \href{https://www-nds.iaea.org/relnsd/vcharthtml/api_v0_guide.html}{primary
  metadata/reference}
\item
  \href{https://nds.iaea.org/relnsd/v1/data}{primary metadata/reference}
\end{itemize}

Blank means unavailable. Preserve operators, cutoff dates, state
identities and extraction dates. Gamma photon energies are not
level-energy differences; uncertainties may be blank and must remain
unknown. API docs identify ENSDF April 2022, not a 2026 reevaluation.

Current link verification is not a new numerical-data retrieval; source
files were not downloaded by this audit.

\subsection{stone-long-lived}\label{stone-long-lived}

N. J. Stone, IAEA INDC(NDS)-0794 (November 2019).

Evaluation or version date: 2019-11

Original numerical retrieval date: not recorded; link-check date is not
numerical retrieval date

\begin{itemize}
\tightlist
\item
  \href{https://nds.iaea.org/records/x2773-kj397}{primary
  metadata/reference}
\item
  DOI
  \href{https://doi.org/10.61092/iaea.yjpc-cns6}{10.61092/iaea.yjpc-cns6}
\end{itemize}

Not a numerical uncertainty source by itself; retain the
quantity-specific convention.

Source report was not downloaded locally by this audit. 2026-10-23 is
the selected LiveChart snapshot retrieval date, not an original report
download or evaluation date.

\subsection{stone-short-lived}\label{stone-short-lived}

N. J. Stone, IAEA INDC(NDS)-0816 (September 2020).

Evaluation or version date: 2020-09

Original numerical retrieval date: not recorded; link-check date is not
numerical retrieval date

\begin{itemize}
\tightlist
\item
  \href{https://nds.iaea.org/records/5g7ec-evw90}{primary
  metadata/reference}
\item
  DOI
  \href{https://doi.org/10.61092/iaea.1p48-p6c6}{10.61092/iaea.1p48-p6c6}
\end{itemize}

Not a numerical uncertainty source by itself; retain the
quantity-specific convention.

Source report was not downloaded locally by this audit. 2026-10-23 is
the selected LiveChart snapshot retrieval date, not an original report
download or evaluation date.

\subsection{stone-quadrupole}\label{stone-quadrupole}

N. J. Stone, IAEA INDC(NDS)-0833.

Evaluation or version date: not recorded

Original numerical retrieval date: not recorded; link-check date is not
numerical retrieval date

\begin{itemize}
\tightlist
\item
  \href{https://nds.iaea.org/records/xk3d8-tw250}{primary
  metadata/reference}
\end{itemize}

Not a numerical uncertainty source by itself; retain the
quantity-specific convention.

Source report was not downloaded locally by this audit. 2026-10-23 is
the selected LiveChart snapshot retrieval date, not an original report
download or evaluation date.

\subsection{be2}\label{be2}

B. Pritychenko, M. Birch, B. Singh and M. Horoi, Atomic Data and Nuclear
Data Tables 107, 1--139 (2016); erratum 114, 371--374 (2017).

Evaluation or version date: 2016

Original numerical retrieval date: 2026-10-23T02:53:18.448505Z

\begin{itemize}
\tightlist
\item
  \href{https://www.nndc.bnl.gov/be2/}{primary metadata/reference}
\item
  \href{https://www.nndc.bnl.gov/be2/help.html}{primary
  metadata/reference}
\item
  \href{https://arxiv.org/abs/1312.5975}{primary metadata/reference}
\item
  \href{https://www.nndc.bnl.gov/be2/data/adopted-entries.json}{actual
  source JSON distribution}
\item
  DOI
  \href{https://doi.org/10.1016/j.adt.2015.10.001}{10.1016/j.adt.2015.10.001}
\item
  DOI
  \href{https://doi.org/10.1016/j.adt.2016.08.002}{10.1016/j.adt.2016.08.002}
\end{itemize}

2016 and fallback 2001 entries stay distinct. Upward 0+ to 2+ B is five
times downward 2+ to 0+ B. Source missing/error operators retained.
Derived beta2 is a collective estimate, not the signed theoretical model
deformation; NNDC help has an inconsistent arrow beside its 0.297
convention. Do not infer branching fractions.

Current link verification is not a new numerical-data retrieval; source
files were not downloaded by this audit.

\subsection{be2-raman2001}\label{be2-raman2001}

S. Raman, C. W. Nestor Jr.~and P. Tikkanen, Atomic Data and Nuclear Data
Tables 78, 1--128 (2001).

Evaluation or version date: 2001

Original numerical retrieval date: not recorded; link-check date is not
numerical retrieval date

\begin{itemize}
\tightlist
\item
  \href{https://www.nndc.bnl.gov/be2/}{primary metadata/reference}
\item
  DOI
  \href{https://doi.org/10.1006/adnd.2001.0858}{10.1006/adnd.2001.0858}
\end{itemize}

Not a numerical uncertainty source by itself; retain the
quantity-specific convention.

Current link verification is not a new numerical-data retrieval; source
files were not downloaded by this audit.

\subsection{codata2018}\label{codata2018}

NIST, Fundamental Physical Constants, 2018 CODATA adjustment, extensive
listing.

Evaluation or version date: 2018

Original numerical retrieval date: not recorded; link-check date is not
numerical retrieval date

\begin{itemize}
\tightlist
\item
  \href{https://physics.nist.gov/cuu/pdf/all_2018.pdf}{primary
  metadata/reference}
\end{itemize}

Software treats constant choices and binding corrections as fixed; their
uncertainties/covariances are not added to mass errors. Keep constant
profile in exported metadata.

Current link verification is not a new numerical-data retrieval; source
files were not downloaded by this audit.

\subsection{codata2022}\label{codata2022}

NIST, Fundamental Physical Constants, 2022 CODATA adjustment, complete
listing.

Evaluation or version date: 2022

Original numerical retrieval date: not recorded; link-check date is not
numerical retrieval date

\begin{itemize}
\tightlist
\item
  \href{https://physics.nist.gov/cuu/Constants/Table/allascii.txt}{primary
  metadata/reference}
\end{itemize}

Software treats constant choices and binding corrections as fixed; their
uncertainties/covariances are not added to mass errors. Keep constant
profile in exported metadata.

GUI default compatibility profile mixes AME2020/CODATA2018 u-energy with
CODATA2022 electron and SI-u masses. A separately named full-2022
profile exists. Website mass kernel remains CODATA2018; do not call the
default entirely CODATA2022.

\subsection{delta-vpn-zhang1989}\label{delta-vpn-zhang1989}

J.-Y. Zhang, R. F. Casten and D. S. Brenner, Physics Letters B 227, 1--5
(1989).

Evaluation or version date: 1989

Original numerical retrieval date: not recorded; link-check date is not
numerical retrieval date

\begin{itemize}
\tightlist
\item
  \href{https://www.sciencedirect.com/science/article/pii/0370269389912732/pdf}{primary
  metadata/reference}
\item
  DOI
  \href{https://doi.org/10.1016/0370-2693(89)91273-2}{10.1016/0370-2693(89)91273-2}
\end{itemize}

Literature attribution is not proof of each software stencil/sign/parity
centre; current equation inventory is authoritative about
implementation.

Current link verification is not a new numerical-data retrieval; source
files were not downloaded by this audit.

\subsection{delta-vpn-cakirli2005}\label{delta-vpn-cakirli2005}

R. B. Cakirli, D. S. Brenner, R. F. Casten and E. A. Millman, Physical
Review Letters 94, 092501 (2005).

Evaluation or version date: 2005

Original numerical retrieval date: not recorded; link-check date is not
numerical retrieval date

\begin{itemize}
\tightlist
\item
  \href{https://journals.aps.org/prl/abstract/10.1103/PhysRevLett.94.092501}{primary
  metadata/reference}
\item
  DOI
  \href{https://doi.org/10.1103/PhysRevLett.94.092501}{10.1103/PhysRevLett.94.092501}
\end{itemize}

Literature attribution is not proof of each software stencil/sign/parity
centre; current equation inventory is authoritative about
implementation.

Current link verification is not a new numerical-data retrieval; source
files were not downloaded by this audit.

\subsection{delta-vpn-cakirli2005-erratum}\label{delta-vpn-cakirli2005-erratum}

R. B. Cakirli, D. S. Brenner, R. F. Casten and E. A. Millman, Physical
Review Letters 95, 119903 (2005).

Evaluation or version date: 2005

Original numerical retrieval date: not recorded; link-check date is not
numerical retrieval date

\begin{itemize}
\tightlist
\item
  \href{https://journals.aps.org/prl/abstract/10.1103/PhysRevLett.95.119903}{primary
  metadata/reference}
\item
  DOI
  \href{https://doi.org/10.1103/PhysRevLett.95.119903}{10.1103/PhysRevLett.95.119903}
\end{itemize}

Literature attribution is not proof of each software stencil/sign/parity
centre; current equation inventory is authoritative about
implementation.

Current link verification is not a new numerical-data retrieval; source
files were not downloaded by this audit.

\subsection{pairing-context-satula1998}\label{pairing-context-satula1998}

W. Satuła, J. Dobaczewski and W. Nazarewicz, Physical Review Letters 81,
3599--3602 (1998).

Evaluation or version date: 1998

Original numerical retrieval date: not recorded; link-check date is not
numerical retrieval date

\begin{itemize}
\tightlist
\item
  \href{https://journals.aps.org/prl/abstract/10.1103/PhysRevLett.81.3599}{primary
  metadata/reference}
\item
  DOI
  \href{https://doi.org/10.1103/PhysRevLett.81.3599}{10.1103/PhysRevLett.81.3599}
\end{itemize}

Literature attribution is not proof of each software stencil/sign/parity
centre; current equation inventory is authoritative about
implementation.

Current link verification is not a new numerical-data retrieval; source
files were not downloaded by this audit.

\subsection{frdm1995}\label{frdm1995}

P. Möller, J.R. Nix, W.D. Myers, W.J. Swiatecki, At. Data Nucl. Data
Tables 59, 185 (1995)

Evaluation or version date: 1992 calculation, 1995 publication

Original numerical retrieval date: not recorded; link-check date is not
numerical retrieval date

\begin{itemize}
\tightlist
\item
  \href{https://doi.org/10.1006/adnd.1995.1002}{model paper}
\item
  \href{https://arxiv.org/abs/nucl-th/9308022}{verified primary
  alternative metadata}
\item
  DOI
  \href{https://doi.org/10.1006/adnd.1995.1002}{10.1006/adnd.1995.1002}
\end{itemize}

No model mass standard uncertainty is supplied (null). Beta2 is not
sigma. Model spread and global rms residuals are not per-nuclide
independent errors. Missing neighbours/shape values remain absent.

legacy TALYS1.95 normalized extraction; see citation

\subsection{frdm2012}\label{frdm2012}

P. Möller, A.J. Sierk, T. Ichikawa, H. Sagawa, At. Data Nucl. Data
Tables 109--110, 1--204 (2016)

Evaluation or version date: 2012 model, 2016 publication

Original numerical retrieval date: not recorded; link-check date is not
numerical retrieval date

\begin{itemize}
\tightlist
\item
  \href{https://doi.org/10.1016/j.adt.2015.10.002}{model paper}
\item
  \href{https://arxiv.org/abs/1508.06294}{verified primary alternative
  metadata}
\item
  DOI
  \href{https://doi.org/10.1016/j.adt.2015.10.002}{10.1016/j.adt.2015.10.002}
\end{itemize}

No model mass standard uncertainty is supplied (null). Beta2 is not
sigma. Model spread and global rms residuals are not per-nuclide
independent errors. Missing neighbours/shape values remain absent.

Frozen attributed table, not a refit to current/new masses.

\subsection{hfb17}\label{hfb17}

S. Goriely, N. Chamel, J.M. Pearson, Phys. Rev.~Lett. 102, 152503 (2009)

Evaluation or version date: 2009

Original numerical retrieval date: not recorded; link-check date is not
numerical retrieval date

\begin{itemize}
\tightlist
\item
  \href{https://doi.org/10.1103/PhysRevLett.102.152503}{model paper}
\item
  DOI
  \href{https://doi.org/10.1103/PhysRevLett.102.152503}{10.1103/PhysRevLett.102.152503}
\end{itemize}

No model mass standard uncertainty is supplied (null). Beta2 is not
sigma. Model spread and global rms residuals are not per-nuclide
independent errors. Missing neighbours/shape values remain absent.

legacy TALYS1.95 normalized extraction; see citation

\subsection{hfbd1m}\label{hfbd1m}

S. Goriely, S. Hilaire, M. Girod, S. Péru, Phys. Rev.~Lett. 102, 242501
(2009)

Evaluation or version date: 2009

Original numerical retrieval date: not recorded; link-check date is not
numerical retrieval date

\begin{itemize}
\tightlist
\item
  \href{https://doi.org/10.1103/PhysRevLett.102.242501}{model paper}
\item
  DOI
  \href{https://doi.org/10.1103/PhysRevLett.102.242501}{10.1103/PhysRevLett.102.242501}
\end{itemize}

No model mass standard uncertainty is supplied (null). Beta2 is not
sigma. Model spread and global rms residuals are not per-nuclide
independent errors. Missing neighbours/shape values remain absent.

legacy TALYS1.95 normalized extraction; see citation

\subsection{hfb14}\label{hfb14}

S. Goriely et al., Phys. Rev.~C 75, 064312 (2007)

Evaluation or version date: 2007

Original numerical retrieval date: not recorded; link-check date is not
numerical retrieval date

\begin{itemize}
\tightlist
\item
  \href{https://doi.org/10.1103/PhysRevC.75.064312}{model paper}
\item
  \href{https://www.astro.ulb.ac.be/Nucdata/Masses/hfb14-plain}{recorded
  source table}
\item
  DOI
  \href{https://doi.org/10.1103/PhysRevC.75.064312}{10.1103/PhysRevC.75.064312}
\end{itemize}

No model mass standard uncertainty is supplied (null). Beta2 is not
sigma. Model spread and global rms residuals are not per-nuclide
independent errors. Missing neighbours/shape values remain absent.

Mcal is atomic mass excess, not full atomic mass. Err/Mexp-Mcal is an
experimental residual, never model sigma. BSkG3 web page count/range is
older than this hashed table; actual parsed coverage retained.

\subsection{hfb24}\label{hfb24}

S. Goriely, N. Chamel, J.M. Pearson, Phys. Rev.~C 88, 024308 (2013)

Evaluation or version date: 2013

Original numerical retrieval date: not recorded; link-check date is not
numerical retrieval date

\begin{itemize}
\tightlist
\item
  \href{https://doi.org/10.1103/PhysRevC.88.024308}{model paper}
\item
  \href{https://www.astro.ulb.ac.be/bruslib/nucdata/hfb24-dat}{recorded
  source table}
\item
  DOI
  \href{https://doi.org/10.1103/PhysRevC.88.024308}{10.1103/PhysRevC.88.024308}
\end{itemize}

No model mass standard uncertainty is supplied (null). Beta2 is not
sigma. Model spread and global rms residuals are not per-nuclide
independent errors. Missing neighbours/shape values remain absent.

Mcal is atomic mass excess, not full atomic mass. Err/Mexp-Mcal is an
experimental residual, never model sigma. BSkG3 web page count/range is
older than this hashed table; actual parsed coverage retained.

\subsection{bskg3}\label{bskg3}

G. Grams et al., Eur. Phys. J. A 59, 270 (2023)

Evaluation or version date: 2023

Original numerical retrieval date: not recorded; link-check date is not
numerical retrieval date

\begin{itemize}
\tightlist
\item
  \href{https://doi.org/10.1140/epja/s10050-023-01158-6}{model paper}
\item
  \href{https://arxiv.org/abs/2307.14276}{verified primary alternative
  metadata}
\item
  \href{https://www.astro.ulb.ac.be/bruslib/nucdata/bskg03-dat}{recorded
  source table}
\item
  DOI
  \href{https://doi.org/10.1140/epja/s10050-023-01158-6}{10.1140/epja/s10050-023-01158-6}
\end{itemize}

No model mass standard uncertainty is supplied (null). Beta2 is not
sigma. Model spread and global rms residuals are not per-nuclide
independent errors. Missing neighbours/shape values remain absent.

Mcal is atomic mass excess, not full atomic mass. Err/Mexp-Mcal is an
experimental residual, never model sigma. BSkG3 web page count/range is
older than this hashed table; actual parsed coverage retained.

\subsection{hfb21}\label{hfb21}

S. Goriely, N. Chamel, J.M. Pearson, Phys. Rev.~C 82, 035804 (2010)

Evaluation or version date: 2010

Original numerical retrieval date: not recorded; link-check date is not
numerical retrieval date

\begin{itemize}
\tightlist
\item
  \href{https://doi.org/10.1103/PhysRevC.82.035804}{model paper}
\item
  \href{https://www.astro.ulb.ac.be/bruslib/nucdata/hfb21-dat}{recorded
  source table}
\item
  DOI
  \href{https://doi.org/10.1103/PhysRevC.82.035804}{10.1103/PhysRevC.82.035804}
\end{itemize}

No model mass standard uncertainty is supplied (null). Beta2 is not
sigma. Model spread and global rms residuals are not per-nuclide
independent errors. Missing neighbours/shape values remain absent.

Mcal is atomic mass excess, not full atomic mass. Err/Mexp-Mcal is an
experimental residual, never model sigma. BSkG3 web page count/range is
older than this hashed table; actual parsed coverage retained.

\subsection{hfb25}\label{hfb25}

S. Goriely, N. Chamel, J.M. Pearson, Phys. Rev.~C 88, 024308 (2013)

Evaluation or version date: 2013

Original numerical retrieval date: not recorded; link-check date is not
numerical retrieval date

\begin{itemize}
\tightlist
\item
  \href{https://doi.org/10.1103/PhysRevC.88.024308}{model paper}
\item
  \href{https://www.astro.ulb.ac.be/bruslib/nucdata/hfb25-dat}{recorded
  source table}
\item
  DOI
  \href{https://doi.org/10.1103/PhysRevC.88.024308}{10.1103/PhysRevC.88.024308}
\end{itemize}

No model mass standard uncertainty is supplied (null). Beta2 is not
sigma. Model spread and global rms residuals are not per-nuclide
independent errors. Missing neighbours/shape values remain absent.

Mcal is atomic mass excess, not full atomic mass. Err/Mexp-Mcal is an
experimental residual, never model sigma. BSkG3 web page count/range is
older than this hashed table; actual parsed coverage retained.

\subsection{hfb26}\label{hfb26}

S. Goriely, N. Chamel, J.M. Pearson, Phys. Rev.~C 88, 024308 (2013)

Evaluation or version date: 2013

Original numerical retrieval date: not recorded; link-check date is not
numerical retrieval date

\begin{itemize}
\tightlist
\item
  \href{https://doi.org/10.1103/PhysRevC.88.024308}{model paper}
\item
  \href{https://www.astro.ulb.ac.be/bruslib/nucdata/hfb26-dat}{recorded
  source table}
\item
  DOI
  \href{https://doi.org/10.1103/PhysRevC.88.024308}{10.1103/PhysRevC.88.024308}
\end{itemize}

No model mass standard uncertainty is supplied (null). Beta2 is not
sigma. Model spread and global rms residuals are not per-nuclide
independent errors. Missing neighbours/shape values remain absent.

Mcal is atomic mass excess, not full atomic mass. Err/Mexp-Mcal is an
experimental residual, never model sigma. BSkG3 web page count/range is
older than this hashed table; actual parsed coverage retained.

\subsection{hfb27}\label{hfb27}

S. Goriely, N. Chamel, J.M. Pearson, Phys. Rev.~C 88, 061302(R) (2013)

Evaluation or version date: 2013

Original numerical retrieval date: not recorded; link-check date is not
numerical retrieval date

\begin{itemize}
\tightlist
\item
  \href{https://doi.org/10.1103/PhysRevC.88.061302}{model paper}
\item
  \href{https://www.astro.ulb.ac.be/bruslib/nucdata/hfb27-dat}{recorded
  source table}
\item
  DOI
  \href{https://doi.org/10.1103/PhysRevC.88.061302}{10.1103/PhysRevC.88.061302}
\end{itemize}

No model mass standard uncertainty is supplied (null). Beta2 is not
sigma. Model spread and global rms residuals are not per-nuclide
independent errors. Missing neighbours/shape values remain absent.

Mcal is atomic mass excess, not full atomic mass. Err/Mexp-Mcal is an
experimental residual, never model sigma. BSkG3 web page count/range is
older than this hashed table; actual parsed coverage retained.

\subsection{ktuy05}\label{ktuy05}

H. Koura, T. Tachibana, M. Uno, M. Yamada, Prog. Theor. Phys. 113,
305--325 (2005)

Evaluation or version date: 2005; fitted to AME2003

Original numerical retrieval date: not recorded; link-check date is not
numerical retrieval date

\begin{itemize}
\tightlist
\item
  \href{https://doi.org/10.1143/PTP.113.305}{model paper}
\item
  \href{https://academic.oup.com/ptp/issue/113/2}{verified primary
  alternative metadata}
\item
  \href{https://wwwndc.jaea.go.jp/nucldata/mass/KTUY05_m246.dat}{recorded
  source table}
\item
  DOI \href{https://doi.org/10.1143/PTP.113.305}{10.1143/PTP.113.305}
\end{itemize}

No model mass standard uncertainty is supplied (null). Beta2 is not
sigma. Model spread and global rms residuals are not per-nuclide
independent errors. Missing neighbours/shape values remain absent.

KTUY05 fitted to AME2003; distinct from KTUY04. Both numerical tables
contain 9436 rows; explanatory prose says 9437. Missing neighbours
remain missing. Source alpha2/alpha4/alpha6 are not beta2; beta2
unavailable. Model sigma unavailable.

\subsection{ws4}\label{ws4}

N. Wang, M. Liu, X. Wu, J. Meng, Phys. Lett. B 734, 215--219 (2014)

Evaluation or version date: 2014-June-3 author-labelled table

Original numerical retrieval date: not recorded; link-check date is not
numerical retrieval date

\begin{itemize}
\tightlist
\item
  \href{https://doi.org/10.1016/j.physletb.2014.05.049}{model paper}
\item
  \href{https://arxiv.org/abs/1405.2616}{verified primary alternative
  metadata}
\item
  \href{http://www.imqmd.com/mass/}{recorded source table}
\item
  DOI
  \href{https://doi.org/10.1016/j.physletb.2014.05.049}{10.1016/j.physletb.2014.05.049}
\end{itemize}

No model mass standard uncertainty is supplied (null). Beta2 is not
sigma. Model spread and global rms residuals are not per-nuclide
independent errors. Missing neighbours/shape values remain absent.

Author-labelled 2014 table supplied for verification. Live author-host
download unavailable at review; papers verified. WS4 and WS4+RBF are
related variants, not independent error estimates. No RBF re-fit to
AME2020 or new input masses.

\subsection{ws4rbf}\label{ws4rbf}

WS4: N. Wang et al., Phys. Lett. B 734, 215--219 (2014); RBF method: N.
Wang, M. Liu, Phys. Rev.~C 84, 051303(R) (2011)

Evaluation or version date: 2014-June-3 author-labelled table; RBF
method 2011

Original numerical retrieval date: not recorded; link-check date is not
numerical retrieval date

\begin{itemize}
\tightlist
\item
  \href{https://doi.org/10.1016/j.physletb.2014.05.049}{model paper}
\item
  \href{https://journals.aps.org/prc/abstract/10.1103/PhysRevC.84.051303}{verified
  primary alternative metadata}
\item
  \href{http://www.imqmd.com/mass/}{recorded source table}
\item
  DOI
  \href{https://doi.org/10.1016/j.physletb.2014.05.049}{10.1016/j.physletb.2014.05.049}
\item
  DOI
  \href{https://doi.org/10.1103/PhysRevC.84.051303}{10.1103/PhysRevC.84.051303}
\end{itemize}

No model mass standard uncertainty is supplied (null). Beta2 is not
sigma. Model spread and global rms residuals are not per-nuclide
independent errors. Missing neighbours/shape values remain absent.

Author-labelled 2014 table supplied for verification. Live author-host
download unavailable at review; papers verified. WS4 and WS4+RBF are
related variants, not independent error estimates. No RBF re-fit to
AME2020 or new input masses.

\section{Individual programs and source file
index}\label{individual-programs-and-source-file-index}

Run these entry points from the GUI checkout. Commands are recorded from
the inventory and were not executed by this documentation build. Full
scientific file hashes and line numbers are retained in
inventory-delivery.json; document generation is not new scientific or
native-GUI acceptance.

\begin{Shaded}
\begin{Highlighting}[]
\ExtensionTok{./Run\_private.command} \AttributeTok{{-}m}\NormalTok{ traplab\_nuclear quantity 116Sn s2n }\AttributeTok{{-}{-}format}\NormalTok{ csv}
\end{Highlighting}
\end{Shaded}

\begin{Shaded}
\begin{Highlighting}[]
\ExtensionTok{./Run\_private.command}\NormalTok{ traplab.py gui }\AttributeTok{{-}{-}page}\NormalTok{ mass\_landscape }\AttributeTok{{-}{-}standalone}
\end{Highlighting}
\end{Shaded}

\begin{Shaded}
\begin{Highlighting}[]
\ExtensionTok{./Run\_private.command}\NormalTok{ traplab.py gui }\AttributeTok{{-}{-}page}\NormalTok{ nuclear\_chart }\AttributeTok{{-}{-}standalone}
\end{Highlighting}
\end{Shaded}

\begin{Shaded}
\begin{Highlighting}[]
\ExtensionTok{./Run\_private.command}\NormalTok{ traplab.py gui }\AttributeTok{{-}{-}page}\NormalTok{ mass\_workbench }\AttributeTok{{-}{-}standalone}
\end{Highlighting}
\end{Shaded}

\begin{Shaded}
\begin{Highlighting}[]
\ExtensionTok{./Run\_private.command}\NormalTok{ apps/ame\_nubase.py }\AttributeTok{{-}{-}quantity}\NormalTok{ s2n }\AttributeTok{{-}{-}z}\NormalTok{ 50 }\AttributeTok{{-}{-}minimum}\NormalTok{ 60 }\AttributeTok{{-}{-}maximum}\NormalTok{ 74}
\end{Highlighting}
\end{Shaded}

\begin{Shaded}
\begin{Highlighting}[]
\ExtensionTok{./Run\_private.command}\NormalTok{ apps/mulberry\_mass.py }\AttributeTok{{-}{-}quantity}\NormalTok{ s2n }\AttributeTok{{-}{-}z}\NormalTok{ 50 }\AttributeTok{{-}{-}minimum}\NormalTok{ 60 }\AttributeTok{{-}{-}maximum}\NormalTok{ 74}
\end{Highlighting}
\end{Shaded}

\begin{Shaded}
\begin{Highlighting}[]
\ExtensionTok{./Run\_private.command}\NormalTok{ scripts/plot\_mass\_workbench.py }\AttributeTok{{-}{-}input}\NormalTok{ private\_examples/mass\_workbench\_v2/116Sn\_ME.csv }\AttributeTok{{-}{-}quantity}\NormalTok{ s2n }\AttributeTok{{-}{-}z}\NormalTok{ 50 }\AttributeTok{{-}{-}minimum}\NormalTok{ 60 }\AttributeTok{{-}{-}maximum}\NormalTok{ 74 }\AttributeTok{{-}{-}models}\NormalTok{ ws4rbf frdm2012 }\AttributeTok{{-}{-}output}\NormalTok{ outputs/mass\_workbench/FRESH\_RUN\_NAME }\AttributeTok{{-}{-}formats}\NormalTok{ png pdf svg}
\end{Highlighting}
\end{Shaded}

\begin{Shaded}
\begin{Highlighting}[]
\ExtensionTok{from}\NormalTok{ gsi\_mrtof.science import nuclear\_mass\_operator}\KeywordTok{;} \ExtensionTok{nuclear\_mass\_operator}\ErrorTok{(}\ExtensionTok{records,z,n,[}\ErrorTok{(}\ExtensionTok{dN,dZ,coefficient}\KeywordTok{)}\ExtensionTok{,...],offset\_keV=0,covariance=None}\KeywordTok{)}
\end{Highlighting}
\end{Shaded}

\begin{Shaded}
\begin{Highlighting}[]
\ExtensionTok{from}\NormalTok{ gsi\_mrtof.nuclear import apply\_energy\_operator}\KeywordTok{;} \VariableTok{result}\OperatorTok{=}\NormalTok{apply\_energy\_operator}\KeywordTok{(}\StringTok{\textquotesingle{}R1D1n1p\textquotesingle{}}\ExtensionTok{,table,be\_unit=}\StringTok{\textquotesingle{}keV\textquotesingle{}}\ExtensionTok{,covariance=None}\KeywordTok{)}
\end{Highlighting}
\end{Shaded}

\begin{Shaded}
\begin{Highlighting}[]
\ExtensionTok{./Run\_private.command} \AttributeTok{{-}m}\NormalTok{ gsi\_mrtof run ame\_analysis }\AttributeTok{{-}{-}inputs}\NormalTok{ @PATH\_TO\_INPUTS\_JSON}
\end{Highlighting}
\end{Shaded}

\begin{Shaded}
\begin{Highlighting}[]
\ExtensionTok{const}\NormalTok{ P=require}\ErrorTok{(}\StringTok{\textquotesingle{}./static/js/physics.js\textquotesingle{}}\KeywordTok{)}\ExtensionTok{,S=require}\ErrorTok{(}\StringTok{\textquotesingle{}./static/js/mass{-}surface.js\textquotesingle{}}\KeywordTok{);} \ExtensionTok{const}\NormalTok{ \{derive\}=S.create}\ErrorTok{(}\ExtensionTok{getter,}\DataTypeTok{\{MEn}\OperatorTok{,}\DataTypeTok{MEH}\OperatorTok{,}\DataTypeTok{MEa}\OperatorTok{,}\DataTypeTok{MEd\}}\ExtensionTok{,betaLookup,evaluatedLookup}\KeywordTok{);} \ExtensionTok{derive}\ErrorTok{(}\ExtensionTok{Z,N,}\StringTok{\textquotesingle{}ame\textquotesingle{}}\KeywordTok{)}\ExtensionTok{.s2n}\KeywordTok{;}
\end{Highlighting}
\end{Shaded}

\begin{itemize}
\tightlist
\item
  energies.py ·
  \texttt{8a8e6ad37aec99fabf5a5d87f5ec3ab70dc998589e990ba8ff7dc74191377677}
\item
  constants.py ·
  \texttt{76c6d2fe38d002795b2ff4997b1850fe125ad35bcb13aca8d6152c34194c37e3}
\item
  models.py ·
  \texttt{e6cc878e85d0bbdaf38fbe65da97776fe7b9dda802cb1c8b8ef0a33c28dc3c0f}
\item
  service.py ·
  \texttt{d1578556ae15e80f6ef3a325a1ddd0af52260a46c6644208c10bdfe87be8eb03}
\item
  decay.py ·
  \texttt{44337db5aa71ec5c0b3f11750673e57ef0ed1045dcfabc61b948926bd849e2a6}
\item
  \textbf{main}.py ·
  \texttt{4d66bd55cb44e79b39a2df3994907e268546e510583c40d6a1211ac191854c08}
\item
  nuclear\_reference.py ·
  \texttt{e92474c8cf387d607fcfb7404fbc3065bd2326e4dd05be23b095c9328bf2556d}
\item
  physics.py ·
  \texttt{ed1487953b7e3998bbb6d00a6952b8fe54d9bbfd4d2c1d81b3a8d97e26ee15cd}
\item
  mass\_workbench.py ·
  \texttt{384516621386296cb0739d796ea6d3a7ce32ec67479c4053ecc31ce893b8977c}
\item
  mass\_landscape.py ·
  \texttt{51b25b54eca939360911ae571572f99ce62b7a4f778229e364189542386a806a}
\item
  nuclear\_chart.py ·
  \texttt{b5247847f7836fe80e6292a2b4bc3fa7c4008ae9cd9d1dfda90fae38b9f772c4}
\item
  mass\_tools.py ·
  \texttt{9a87d2f3d46a816a32a90687b3cd83fc3af3ad71a0654b4dd0a7d3d1be62604e}
\item
  mass\_workbench\_ui.py ·
  \texttt{3f03ecf297f24db5fff2c99bc3d51ebd3f7029d6e8f5012686e31b4e3aa0b78f}
\item
  plot\_mass\_workbench.py ·
  \texttt{caacefbba4df1eda7887b75f33e2bd1bf491b58b5066cf2fcc1cbfaf6a42b6f9}
\item
  ame\_nubase.py ·
  \texttt{4fae5f046e983f4dfee8f765848127663811c70d0faef80c75d4a91df6dde1a4}
\item
  mulberry\_mass.py ·
  \texttt{b9cfa27bf5ba371d891deb24148b80653c7731b1e26f0ca8533166c52d41f157}
\item
  run\_individual.py ·
  \texttt{e1edc328dd460b0267498165827a3934690ad57ec60fea40c54d8e2e7846804a}
\item
  nuclear.py ·
  \texttt{721f27ce9a6a1d44af4d7e25bb7a285a10b73f03967ad04b8473dd4b07a5116b}
\item
  science.py ·
  \texttt{7988e5dbebdd7e5e9d0e33db2a9a138a1555d80d79b72f71d0e24164c75becb1}
\item
  \textbf{main}.py ·
  \texttt{31fc3c5fa727f82c7ed0c53b17feb4b28b06f75a0a3743f7c5332d9096757490}
\item
  nuclear\_operator\_metadata.csv ·
  \texttt{e7ac96558532287f83f84f145e2196f72129db10fe6574d29eec3abec52138af}
\item
  nuclear\_operator\_weights.csv ·
  \texttt{50691edddf6b7b1c8b9207ca281c719769c155468a1bb04b9b94e1b71dfc6865}
\item
  NUCLEAR\_REFERENCE\_EN.md ·
  \texttt{e8f684dc1273542210fbb9450ba1bcd0b24670788d2a375d9e3fc915471c9228}
\item
  NUCLEAR\_DATA\_EN.md ·
  \texttt{7c32914782e2b39bc0d008b865c95f67f6a6552ce4fab02cdf286adda03395a9}
\item
  MASS\_WORKBENCH\_EN.md ·
  \texttt{280072337fe3aed9eb9624e35cd2d76a068199c3a738144136b8fedb6c54772e}
\item
  AME\_MULBERRY\_FEATURE\_REVIEW\_EN.md ·
  \texttt{28fa3b49ccde73f90cc08299093a5d221d6db39962dd539c040120a13acd9e96}
\item
  mass-surface.js ·
  \texttt{0563dd3aa7ba2633d9867051dc4035db91195bd736106926ad3fc0963bdb8fcc}
\item
  physics.js ·
  \texttt{7bdf10025e7571b1dea722c2407305f2c2332a24cdcc184baaec317087c0746a}
\item
  ame-evaluated.js ·
  \texttt{e4fc2de8ef689dd32765e4c2ad2bfe86081c713179ce89e65c9e5f0deda70cfd}
\item
  nuclide-chart.js ·
  \texttt{7ea908222322320624d4e2c2b277d3cb23c5292fe6a2180afe03d0ac91870bd5}
\item
  mass-quantity-menu.js ·
  \texttt{68654dd057ac856e1caf35b8e35bb1be52be1eb49677c2885be22f12226f69ad}
\item
  measured-masses.js ·
  \texttt{730f1daed0677af8848b09c05d3bf709b42db0d2d3310e9a435c9c51409158c0}
\item
  ame2020\_reference.csv ·
  \texttt{f8cddd5544c1d7fa3f56213c58b3bde1d1e6cc7ec3c69e8b7933bfdb1612fca8}
\item
  ame2020\_reference\_metadata.csv ·
  \texttt{df9652aa9e6985686583eb3782bd568232356c8fbaace073778986d4cc5c448f}
\item
  index.json ·
  \texttt{71338c7c73837caa9b24092d43d93707ffa3b287302ae9d76162e5386c2ff834}
\item
  ame2020-reactions.json ·
  \texttt{902d19a60e7b42772044c06c2277615d9f7f9cbb90454baa85ebff2535b48ffc}
\item
  chart-typed-data.js ·
  \texttt{bb9aef72fbef7ef49817aee63e6b0882c869b63c4cc1f600195d85803e0208bc}
\item
  measured-masses-ui.js ·
  \texttt{e698b612fa6c5343d1a56b81d4e9f8c915f383022314868ccc1955cefd685046}
\item
  nuclide\_chart\_labels.yaml ·
  \texttt{f193f37649fb262df9d6ea92fd41b4a51eacbcef35ad01420e8e7e7b5aecfe28}
\item
  nuclide-chart.html ·
  \texttt{49904708401e5c3b12416314c5aee21bb7cb6943eb0c0637aa35b140911e3ba2}
\item
  chart\_typed\_data.test.cjs ·
  \texttt{e0553c10c99e76d72ee75fa8f65468d12dd9d6d872de902d7810156eb9988a3b}
\item
  measured\_mass\_ui.test.cjs ·
  \texttt{35f4330dffdee8dcbb863fee2b1916ddfc2bb34f7a80c402748c30e6dcd9ff5e}
\item
  chart\_controls.test.cjs ·
  \texttt{7ace99cded8be5a33c83f3c19692878dab005b271e22e819e8ed35635b9a4dc8}
\item
  test\_chart\_script\_dependencies.py ·
  \texttt{226d42e1744cc6587f3133e5781030c3dc2203b9b9c0a9db537e63475ef0820b}
\end{itemize}

\section{Modification reproduction and evidence
limits}\label{modification-reproduction-and-evidence-limits}

When equations, stencils, units or provenance change, update the
machine-readable inventory and run existing scientific regressions
before rebuilding the guide. Do not loosen independent reference
tolerances merely to pass a test. Results with different source or
uncertainty models may differ; retain the difference and its reason.
Module names identify implementation entry points; the companion
application is installed and run locally. Public-edition file hashes are
listed in MANIFEST.json.

This public edition provides portable Markdown, editable Word equations,
offline MathML HTML and standalone TeX source. Equations are editable
Word OMML, not images. The owner's explicit instruction prohibits
creating or saving screenshots of manual pages. OOXML structure, math
lint and optional vector-PDF text/page-bound checks do not establish
all-page visual layout acceptance in Word. Scientific-data verification,
browser interactions, downloads, deployment, Qt, Windows and hardware
acceptance are separate evidence levels.

\subsection{Reused bilingual
documentation}\label{reused-bilingual-documentation}

\begin{itemize}
\tightlist
\item
  AME audit
\item
  Mass workbench
\item
  Nuclear reference
\end{itemize}

\end{document}
